\documentclass[10pt]{amsart}
\usepackage{amsmath,amssymb,amsthm,mathtools}
\newtheorem{theorem}{Theorem}[section]
\newtheorem{lemma}[theorem]{Lemma}
\newtheorem{proposition}[theorem]{Proposition}
\newtheorem{corollary}[theorem]{Corollary}
\theoremstyle{definition}
\newtheorem{definition}[theorem]{Definition}
\newtheorem{problem}[theorem]{Problem}
\theoremstyle{remark}
\newtheorem{remark}[theorem]{Remark}
\newcommand{\NA}{\operatorname{NA}}
\newcommand{\supp}{\operatorname{supp}}
\newcommand{\sgn}{\operatorname{sign}}
\newcommand{\spn}{\operatorname{span}}
\newcommand{\dist}{\operatorname{dist}}
\newcommand{\R}{\mathbb R}
\newcommand{\N}{\mathbb N}
\newcommand{\Cset}{\mathcal C}
\newcommand{\Mset}{\mathcal M}
\newcommand{\Cert}{\operatorname{Cert}}
\newcommand{\Li}{\operatorname*{Li}}
\newcommand{\cl}{\operatorname{cl}}
\newcommand{\gap}{\operatorname{gap}}
\newcommand{\Def}{\operatorname{Def}}
\newcommand{\Rec}{\mathcal R}
\newcommand{\zh}{\hat z}
\newcommand{\ip}[2]{\langle #1,#2\rangle}
\newcommand{\lin}{\mathfrak l}
\newcommand{\Exc}{\operatorname{Exc}}
\title{Recovery of mates on Mart\'in norms\\and the resonance obstruction}
\author{}
\date{Research continuation, October 2026}
\begin{document}
\begin{abstract}
For Mart\'in's renormings $p$ of $c_0$ (canonical base, finitely or
infinitely many blocks) we study the density of norm-attaining operators
into $\ell_2^2$ through the mate fibres
$\Cset(f)=\{g:\ (f,g)\text{ is contractive}\}$.  Density is equivalent to the
statement that every pair $(f,\rho g)$ with $g\in\Cset(f)$ and $\rho<1$ is a
limit of norm-attaining pairs, and we reduce this to lower limits of the
fibres along norm-attaining approximants; finitely many blocks suffice.  We
introduce finite certificates with coordinatewise radii, prove that the
closed set they generate is transported along \emph{every} sequence
$f_n\to f$, and prove an averaging theorem over scales which covers weighted
directions, locally split mates and cross mates carried by off-peak
coordinates with non-summable gaps.  At tame supports we identify the sharp
second-order invariant, the mass-weighted coefficient $\Gamma_w$, and
realize it along canonical truncations by means of transfer peaks.  We then
construct an admissible operator $T$ and a non-attaining first row at which
all intrinsically recoverable mates lie on one line while $\Cset(f)$ spans an
infinite-dimensional space (the resonance obstruction); the canonical
truncations do not recover these mates.

This is a defect of the intrinsic mechanisms only.  Rebalancing through
transfer peaks at first order makes \emph{engineered} norm-attaining
approximants work without any steering: every mate with one-sided linear
decompositions of weighted coefficient at most one and non-negative
$d$-mismatch is recovered, for arbitrary contact sets and any number of
blocks, and the case of negative mismatch is reduced to a scrambling
condition.  By Fenchel duality we compute the exact one-sided second-order
coefficients at block-tame first rows, and deduce that every block-tame
first row is recoverable, whatever its contact set; in particular every mate
of the resonant example is recovered.  Finally we design an admissible
operator (private signatures, coarse-to-fine targets, a super-fast ladder
of weights) for which every first row with finite base support and
signature room is recoverable, and for which density is equivalent to a
single approximation statement, Lemma~Z.  The density problem itself
remains open; we formulate precisely what remains.
\end{abstract}
\maketitle
\tableofcontents

\section{Setting and standing facts}\label{sec:setting}

All spaces are real.  We write $s(t):=\sqrt{1+t^2}$ throughout.

\subsection{The base, the blocks and Mart\'in's lemmas}
We use the canonical base on $X=c_0$, as in the accompanying note
\cite{PrA}: $H$ is a Hilbert space, $U:H\to c_0$ is compact with dense
range, and
\[
 B_q=B_{c_0}+U(B_H),\qquad q^*(a)=\|a\|_1+\|U^*a\|_H\quad(a\in\ell_1=X^*).
\]
Since $U$ has dense range, $U^*$ is injective; it is compact, so
$\|U^*e_j^*\|_H\to0$.  The bidual ball is $B_{q^{**}}=B_{\ell_\infty}+U(B_H)$
(the sum of a weak$^*$-compact set and a norm compact set is weak$^*$
compact, contains $B_q$, and is contained in its weak$^*$ closure).  The
norm $q$ is a smoothing of the supremum norm in the sense of
Debs, Godefroy and Saint-Raymond \cite{DGS}.

Mart\'in's construction \cite{Martin} rests on two lemmas, which we recall
in the form stated there.

\begin{remark}[Mart\'in's Lemma A]\label{rem:lemmaA}
Let $\Phi$ be a sequence of positive numbers with $\|\Phi\|_1<1$ and let
$D_\Phi:\ell_2\to\ell_1$ be the diagonal operator with entries $\Phi(n)$.
The norm $|\cdot|_\Phi$ on $\ell_1$ with unit ball
$B_{\ell_1}+D_\Phi(B_{\ell_2})$ satisfies:
(a) $|f|_\Phi^*=\|f\|_\infty+\|D_\Phi f\|_2$ for $f\in\ell_\infty$;
(b) $|f|^*_\Phi\geq\|f\|_\infty\geq(1-\|\Phi\|_1)|f|^*_\Phi$;
(c) $|x|_\Phi\leq\|x\|_1\leq(1+\|\Phi\|_1)|x|_\Phi$;
(d) the dual norm is strictly convex, so $|\cdot|_\Phi$ is smooth;
(e) if $|f|^*_\Phi=1$, $\ip{f}{x}=|x|_\Phi$ and $|x(n)|>\Phi(n)|x|_\Phi$,
then $f(n)=\sgn(x(n))\|f\|_\infty$.
\end{remark}

\begin{remark}[Mart\'in's Lemma B and the choice of $T$]\label{rem:lemmaB}
Lemma B of \cite{Martin} reads: \emph{Let $Y$ be a separable operator
range.  Then there is a norm-one injective operator
$T:\ell_1(\N\times\N)\to Y$ such that for every $m$ the set
$\{T(e_{n,m})/\|T(e_{n,m})\|:n\in\N\}$ is dense in $S_Y$.}
Its proof takes a bijection $\sigma:\N\times\N\to\N$, a sequence
$(w'_n)$ in $S_Y$ such that every point of $S_Y$ is an accumulation point
of $(w'_n)$, and repeats the proof of \cite[Proposition~2.8]{KLMW20} with
the sequence $w_{\sigma(n,m)}=w'_n$.  Two features matter below.  First,
the sequence $(w'_n)$ is \emph{arbitrary} subject to the accumulation
property, so ``Mart\'in's space'' is a family of spaces indexed by the
admissible choices of $T$.  Second, every block $m$ follows the same
sequence $(w'_n)$, so that near duplicates $Te_{n,m}/\|Te_{n,m}\|\approx
Te_{n,m'}/\|Te_{n,m'}\|$ across blocks are built into the construction (up
to the perturbations produced in \cite{KLMW20}).  No other property of $T$
is used in \cite{Martin}, and no other property will be used here unless
it is stated explicitly.
\end{remark}

\begin{definition}[Admissible operators]\label{def:admissible}
An operator $T:\ell_1(\N\times\N)\to X^*=\ell_1$ is \emph{admissible} if
\begin{enumerate}
\item[(T-a)] $q^*(Te_{n,m})\le1$ for all $(n,m)$, with equality for some
$(n,m)$;
\item[(T-b)] $T$ is injective;
\item[(T-c)] $Y:=T(\ell_1(\N\times\N))$ satisfies $Y\cap c_{00}=\{0\}$;
\item[(T-d)] for every $m$, the set $\{u_{n,m}:n\in\N\}$, where
$u_{n,m}:=Te_{n,m}/q^*(Te_{n,m})$, is norm dense in $S_{q^*}$.
\end{enumerate}
\end{definition}

Since $\NA((X,q),\R)=c_{00}$ (Proposition~\ref{prop:smooth}(c) below),
(T-c) is the condition $Y\cap\NA(q)=\{0\}$ of \cite{Martin}.  As $S_{q^*}$
has no isolated points, (T-d) implies that every tail
$\{u_{n,m}:n\geq n_0\}$ is dense in $S_{q^*}$.

\subsection{The norm $p$}
Fix an admissible $T$.  For $m\in\N$ put
\begin{gather*}
 \Phi_m(k):=2^{-m-k}q^*(Te_{k,m})\le 2^{-m-k},\qquad
 \lambda_{k,m}:=m\Phi_m(k),\\
 (R_mx)(k):=\lambda_{k,m}u_{k,m}(x)\qquad(x\in X).
\end{gather*}
Let $D_m$ be the diagonal operator with entries $\Phi_m(k)$, let
$|\cdot|_m$ be the norm of Remark~\ref{rem:lemmaA} for $\Phi=\Phi_m$ on
$V_m=\ell_1$, and $N_m(w):=|w|_m^*=\|w\|_\infty+\|D_mw\|_2$.  For a nonempty
set $I\subseteq\N$ of \emph{blocks} put
\begin{gather*}
 V:=\Big(\bigoplus_{m\in I}V_m\Big)_{\ell_1},\qquad
 V^*=\Big(\bigoplus_{m\in I}(\ell_\infty,N_m)\Big)_{\ell_\infty},\\
 Lx:=(R_mx)_{m\in I},\qquad p(x):=q(x)+\|Lx\|_V .
\end{gather*}
For $I=\N$ this is Mart\'in's norm (with the canonical base); for
$I=\{1,\dots,N\}$ we obtain the finite-block norms $p_N$ of \cite{PrB}.
Since $|(R_mx)(k)|\le m2^{-m-k}q(x)$, $\|R_m\|\leq m2^{-m}$; each $R_m$ is a
norm limit of finite-rank operators, hence compact, and $L$ is compact.
For $w\in V^*$,
\begin{equation}\label{eq:Lstar}
 L^*w=\sum_mR_m^*w_m=T\Big(\sum_{m\in I}\sum_k m2^{-m-k}w_m(k)\,e_{k,m}\Big)\in Y,
\end{equation}
an absolutely convergent series; in particular $L^*$ is injective on
$V^*$.  For $\theta\in X^{**}=\ell_\infty$,
$(R_m^{**}\theta)(k)=\lambda_{k,m}\theta(u_{k,m})$, so
$R_m^{**}(X^{**})\subseteq\ell_1$.

\begin{remark}[Finite and infinite block sets]\label{rem:blocks}
By Lemma~\ref{lem:martintail} below (Mart\'in tails; cf.\
\cite[Remark~6.9]{PrB}), if $\NA((c_0,p_N),\mathcal F)$ is dense in
$\mathcal L((c_0,p_N),\mathcal F)$ for infinitely many $N$, then
$\NA((c_0,p),\mathcal F)$ is dense for Mart\'in's $p$ ($I=\N$).  Hence
finite block sets suffice for the density problem.  Sections~\ref{sec:setting}--\ref{sec:averaging} are
valid for every nonempty $I$ (certificates involve finitely many blocks),
with two exceptions.  Corollary~\ref{cor:decomposition} needs an extra
hypothesis when $I$ is infinite (Remark~\ref{rem:G2}).
Theorem~\ref{thm:locsplit} and Corollary~\ref{cor:splitcontractive} are
stated and proved only for finite $I$: their proof truncates every block,
and the admissible range of the block truncations shrinks with the block
index, so for infinitely many blocks it is not uniform.  By the first
part of this remark, finite $I$ is enough for the density problem, so
both results are covered by the reduction to $p_N$; we make no claim about
locally split mates for $I=\N$.  Sections~\ref{sec:second}--\ref{sec:designed} assume that $I$ is finite.
\end{remark}

\begin{lemma}[Mart\'in tails]\label{lem:martintail}
Fix an admissible $T$, let $p$ be the norm with $I=\N$ and $p_N$ the norm
with $I=\{1,\dots,N\}$, and let $\mathcal F$ be a Banach space.  If
$\NA((c_0,p_N),\mathcal F)$ is dense in $\mathcal L((c_0,p_N),\mathcal F)$
for infinitely many $N$, then $\NA((c_0,p),\mathcal F)$ is dense in
$\mathcal L((c_0,p),\mathcal F)$.
\end{lemma}

\begin{proof}
Put $s_N(x):=\sum_{m>N}|R_mx|_m$, so that $p=p_N+s_N$.  Since
$|R_mx|_m\le\|R_mx\|_1\le m2^{-m}q(x)$, we have $s_N\le\eta_Nq\le\eta_Np_N$
with $\eta_N:=\sum_{m>N}m2^{-m}=(N+2)2^{-N}$, hence $\|A\|_p\le\|A\|_{p_N}\le
(1+\eta_N)\|A\|_p$ for operators $A$.  Let $S\in\NA((c_0,p_N),\mathcal F)$
with $\|S\|_{p_N}=1=\|Sx_0\|$, $p_N(x_0)=1$.  Let $J$ be a norm-one
functional on $(\bigoplus_{m>N}V_m)_{\ell_1}$ norming $(R_mx_0)_{m>N}$
($J:=0$ if this vector is $0$), and $\varphi(x):=J((R_mx)_{m>N})$; then
$\varphi\in X^*$, $|\varphi|\le s_N$ and $\varphi(x_0)=s_N(x_0)$.  Put
$S'x:=Sx+\varphi(x)Sx_0$.  Then $\|S'x\|\le p_N(x)+s_N(x)=p(x)$ and
$\|S'x_0\|=1+s_N(x_0)=p(x_0)$, so $S'\in\NA((c_0,p),\mathcal F)$, and
$\|S'-S\|_p\le\sup_x|\varphi(x)|/p(x)\le\eta_N$.  Given $A\ne0$ and
$\varepsilon>0$, density for $p_N$ gives such an $S$ with
$\|S-A/\|A\|_{p_N}\|_{p_N}<\varepsilon$; then $\|A\|_{p_N}S'$ attains its
$p$-norm and $\|\,\|A\|_{p_N}S'-A\|_p\le(1+\eta_N)\|A\|_p(\eta_N+
\varepsilon)$.  Let $\varepsilon\to0$ and $N\to\infty$ along the given
$N$.
\end{proof}

\subsection{Elementary structure}

\begin{lemma}[Dual ball]\label{lem:dualball}
$B_{p^*}=B_{q^*}+L^*(B_{V^*})$, and for every $h\in X^*$
\[
 p^*(h)=\min\{\max(q^*(A),\|W\|_{V^*}):\ A+L^*W=h\},
\]
the minimum being attained.  Moreover $p^{**}(\theta)=q^{**}(\theta)+
\|L^{**}\theta\|_V$ for $\theta\in X^{**}$, and $L^{**}(X^{**})\subseteq V$.
\end{lemma}

\begin{proof}
If $q^*(A)\le1$ and $\|W\|\le1$ then $(A+L^*W)(x)\le q(x)+\|Lx\|=p(x)$.
The set $B_{q^*}+L^*(B_{V^*})$ is convex and weak$^*$ compact ($L^*$ is
weak$^*$-continuous), and its support function on $X$ is
$q(x)+\|Lx\|_V=p(x)$, which is also the support function of $B_{p^*}$; two
weak$^*$-closed convex sets with the same support function on $X$
coincide.  The formula for $p^*$ follows by homogeneity and compactness.
Taking suprema over $B_{p^*}$ gives the bidual formula; $L^{**}$ maps into
$V$ because $L$ is compact.
\end{proof}

\begin{lemma}[Block threshold; cf.\ {\cite[Lemma~6.5]{PrB}}]\label{lem:threshold}
Fix a block and drop the index $m$.  Let $\zeta\in\ell_1\setminus\{0\}$ and
$w=J(\zeta)$ its norming functional ($N(w)=1$, $w(\zeta)=|\zeta|$).  Put
$M=\|w\|_\infty$, $C=\|Dw\|_2$ and $P=\{k:|w(k)|=M\}$.  Then $M>0$, $C>0$,
$M+C=1$, and
\begin{gather*}
 \zeta/|\zeta|=\alpha+D^2w/C,\qquad \|\alpha\|_1=1,\\
 \supp\alpha\subseteq P,\qquad \sgn\alpha(k)=\sgn w(k)\quad(k\in\supp\alpha).
\end{gather*}
In particular $P\ne\emptyset$; off $P$ one has
$w(k)=C\zeta(k)/(\Phi(k)^2|\zeta|)$, and $|\zeta(k)|>|\zeta|\Phi(k)^2M/C$
implies $k\in P$ and $w(k)=M\sgn\zeta(k)$.
\end{lemma}

\begin{proof}
$w\neq0$ and $D$ is injective, so $M,C>0$.  The ball
$B_{\ell_1}+D(B_{\ell_2})$ is closed ($D$ is compact), so
$\zeta/|\zeta|=\alpha+D\beta$ with $\|\alpha\|_1,\|\beta\|_2\le1$.  Then
$1=\ip{w}{\alpha}+\ip{Dw}{\beta}\le M\|\alpha\|_1+C\|\beta\|_2\le M+C=1$,
so $\ip{w}{\alpha}=M=M\|\alpha\|_1$ and $\beta=Dw/C$; the first equality
forces $\|\alpha\|_1=1$ and $\alpha$ to live on $P$ with the signs of $w$.
The remaining claims follow by reading the identity coordinatewise.
\end{proof}

For a block vector $\zeta=R_m^{**}\theta$ the condition of the last
sentence reads $|\theta(u_{k,m})|>\vartheta_m\Phi_m(k)$ with the
\emph{threshold constant} $\vartheta_m:=|\zeta|_mM_m/(mC_m)$.  For $k\in P$
we define the \emph{margin}
$\mu_{k,m}:=|\theta(u_{k,m})|-\vartheta_m\Phi_m(k)\ge0$; reading
Lemma~\ref{lem:threshold} at $k\in P$ gives the identity
\begin{equation}\label{eq:margin}
 |\zeta|_m|\alpha_m(k)|=\lambda_{k,m}\mu_{k,m}\qquad(k\in P_m).
\end{equation}
Peaks with $\alpha_m(k)=0$ (equivalently $\mu_{k,m}=0$) are called
\emph{degenerate}.  We write $Q_m:=\N\setminus P_m$ (strict non-peaks) and
$\gap_m(k):=M_m-|w_m(k)|$.

\begin{lemma}[Rigidity and independence]\label{lem:rigidity}
\textup{(a)} If $b_++L^*W_+=b_-+L^*W_-$ with $b_+-b_-\in c_{00}$ and
$W_\pm\in V^*$, then $b_+=b_-$ and $W_+=W_-$.
\textup{(b)} Let $J\subseteq\N$ be cofinite, $I$ finite, and for each
$m\in I$ let $\Lambda_m\subseteq\N$ be finite and $P_m'\subseteq\N$ be
infinite, $s\in\{\pm1\}^{P'_m}$.  Put
$\pi_m:=\sum_{k\in P'_m}s_k\lambda_{k,m}u_{k,m}$.  Then the restrictions to
$c_{00}(J)$ of the functionals $u_{k,m}$ \textup{(}$k\in\Lambda_m$,
$m\in I$\textup{)} and $\pi_m$ \textup{(}$m\in I$\textup{)} are linearly
independent; consequently $\nu\mapsto(u_{k,m}(\nu),\pi_m(\nu))$ maps
$c_{00}(J)$ onto the corresponding finite-dimensional space.
\end{lemma}

\begin{proof}
(a) $b_+-b_-=L^*(W_--W_+)\in Y\cap c_{00}=\{0\}$, and $L^*$ is injective.
(b) A vanishing combination
$\sum c_{k,m}u_{k,m}+\sum_md_m\pi_m$ on $c_{00}(J)$ is an element of $Y$
supported in the finite set $\N\setminus J$, hence zero by (T-c).  By
(T-b) its coefficient at every $e_{k,m}$ vanishes: at the infinitely many
$k\in P'_m\setminus\Lambda_m$ this coefficient is $d_ms_km2^{-m-k}$, so
$d_m=0$, and then $c_{k,m}=0$.
\end{proof}

\begin{proposition}[Strict convexity and forced decomposition]\label{prop:forced}
\textup{(a)} Each $R_m^{**}$ is injective, and $p^{**}$ is strictly convex.

\textup{(b)} Every $f\in S_{p^*}$ has a unique normer $\xi\in S_{p^{**}}$.
Put $q_0:=q^{**}(\xi)$.  Then $0<q_0<1$, $\|L^{**}\xi\|_V=1-q_0$, every
$R_m^{**}\xi\ne0$, and there is exactly one decomposition $f=a+L^*w$ with
$q^*(a)\le1$ and $\|w\|_{V^*}\le1$; it satisfies
\[
 q^*(a)=1,\quad a(\xi)=q_0,\quad w_m=J_m(R_m^{**}\xi),\quad N_m(w_m)=1
 \quad(m\in I).
\]
\textup{(c)} Put $\nu:=\|U^*a\|_H>0$ and $e:=U^*a/\nu$.  Then
$\zh:=\xi/q_0=z+Ue$ with $z\in B_{\ell_\infty}$ and $z_j=\sgn a_j$ on
$F:=\supp a$; and $a(\zh)=1$, $b(\zh)=b(z)+\ip{U^*b}{e}$ for
$b\in\ell_1$.
\end{proposition}

\begin{proof}
(a) If $\theta(u_{k,m})=0$ for all $k$, then $\theta$ vanishes on a dense
subset of $S_{q^*}$, so $\theta=0$.  Let
$p^{**}(\xi)=p^{**}(\eta)=p^{**}((\xi+\eta)/2)=1$.  Since $q^{**}$ and each
$|R_m^{**}\cdot|_m$ are seminorms whose sum is $p^{**}$
(Lemma~\ref{lem:dualball}), we get
$|R_m^{**}(\xi+\eta)|_m=|R_m^{**}\xi|_m+|R_m^{**}\eta|_m$ for every $m$.
Fix $m$; $\zeta_1=R_m^{**}\xi$ and $\zeta_2=R_m^{**}\eta$ are nonzero, and a
norming functional $w$ of $\zeta_1+\zeta_2$ norms both.  If $\xi,\eta$ were
not positively proportional there would be $v\in S_{q^*}$ with
$\xi(v)>0>\eta(v)$ (use surjectivity of $v\mapsto(\xi(v),\eta(v))$ if they
are independent).  Choose $n_j\to\infty$ with $u_{n_j,m}\to v$.  Since
$\Phi_m(n_j)\to0$, Lemma~\ref{lem:threshold} applied to $\zeta_1$ and to
$\zeta_2$ gives $w(n_j)=\|w\|_\infty$ and $w(n_j)=-\|w\|_\infty$ for large
$j$, which is absurd.  Hence $\eta=\lambda\xi$ with $\lambda>0$, and
$\lambda=1$.

(b) The normers of $f$ form a convex subset of $S_{p^{**}}$, hence a point.
For any decomposition $f=a+L^*w$ with $q^*(a),\|w\|\le1$,
$1=f(\xi)=a(\xi)+w(L^{**}\xi)\le q^{**}(\xi)+\|L^{**}\xi\|=1$, so
$a(\xi)=q_0$ and $w_m(R_m^{**}\xi)=|R_m^{**}\xi|_m$ for each $m$.  As
$R_m^{**}\xi\ne0$ and $|\cdot|_m$ is smooth, $w_m=J_m(R_m^{**}\xi)$ is
determined, hence so is $a=f-L^*w$.  Finally $q_0>0$ because $q^{**}$ is a norm and
$\xi\ne0$, $q_0<1$ because $L^{**}\xi\ne0$, and
$q_0=a(\xi)\le q^*(a)q_0$ gives $q^*(a)=1$.

(c) $\xi/q_0\in B_{q^{**}}=B_{\ell_\infty}+U(B_H)$, so $\xi/q_0=z+Uh$ with
$\|z\|_\infty,\|h\|\le1$; then $1=a(z)+\ip{U^*a}{h}\le\|a\|_1+\|U^*a\|=1$
forces $a(z)=\|a\|_1$ and $h=e$.
\end{proof}

Throughout, the \emph{forced data} of $f\in S_{p^*}$ are
$(\xi,q_0,a,w,F,z,\zh,e,\nu)$ and, for each block,
$\zeta_m:=R_m^{**}\xi$, $\sigma_m:=|\zeta_m|_m$, $M_m$, $C_m$, $P_m$,
$\alpha_m$, $Q_m$, $\gap_m$, $\vartheta_m$, $\mu_{k,m}$.  Note
$q_0+\sum_m\sigma_m=1$.  The set $K:=\{j\notin F:|z_j|=1\}$ is the set of
\emph{contacts} (kinks) and $J:=\N\setminus(F\cup K)$ the set of
\emph{free coordinates}; the \emph{room} of $j\in J$ is $1-|z_j|>0$.

\subsection{Smoothness and norm-attaining functionals}

\begin{proposition}[Fr\'echet smoothness]\label{prop:smooth}
\textup{(a)} The norm $p$ is Fr\'echet differentiable at every $x\neq0$,
and $\nabla p(x)=\nabla q(x)+L^*J_V(Lx)$, where $J_V(Lx)=(J_m(R_mx))_m$.

\textup{(b)} $\NA((X,p),\R)\cap S_{p^*}=\nabla p(S_p)$, and the duality map
$\nabla p:S_p\to S_{p^*}$ is norm-to-norm continuous.

\textup{(c)} $\NA((X,q),\R)=c_{00}$.  More precisely, if $a'\in c_{00}$,
$q^*(a')=1$, and $z'\in c_0$ satisfies $\|z'\|_\infty\le1$ and
$z'_j=\sgn a'_j$ on $\supp a'$, then $y':=z'+U(U^*a'/\|U^*a'\|)$ satisfies
$q(y')=1$, $\nabla q(y')=a'$ and $\nabla p(y')=a'+L^*J_V(Ly')$.
Conversely every $f'\in\NA((X,p),\R)\cap S_{p^*}$ equals $\nabla p(y')$ for
such a pair $(a',z')$, with $a'=\nabla q(y')$.
\end{proposition}

\begin{proof}
\emph{Gateaux smoothness.}  The norm $q^*$ is strictly convex: if
$q^*(a)=q^*(b)=1=q^*(\frac{a+b}2)$, then $\|U^*a+U^*b\|=\|U^*a\|+\|U^*b\|$, so
$U^*b=\lambda U^*a$ with $\lambda\ge0$ (both are nonzero), hence $b=\lambda a$
by injectivity of $U^*$, and $\lambda=1$.  Thus $q$ is Gateaux smooth.  In
the same way $N_m$ is strictly convex ($D_m$ is injective), so $|\cdot|_m$
is smooth.  Let $x\neq0$ and let $f,f'\in S_{p^*}$ both norm $x$.  Write
$f=A+L^*W$ with $q^*(A),\|W\|\le1$ (Lemma~\ref{lem:dualball}).  Then
$p(x)=A(x)+W(Lx)\le q(x)+\|Lx\|=p(x)$ forces $A(x)=q(x)$ and
$W_m(R_mx)=|R_mx|_m$ for every $m$.  Since $R_m$ is injective (by the
argument of Proposition~\ref{prop:forced}(a)), $R_mx\ne0$, so $A=\nabla q(x)$
and $W_m=J_m(R_mx)$.  The same holds for $f'$, so $f=f'$.

\emph{Fr\'echet smoothness.}  By \v{S}mulyan's lemma it suffices to show:
if $x\in S_p$, $f_n\in B_{p^*}$ and $f_n(x)\to1$, then
$f_n\to\nabla p(x)$ in norm.  Write $f_n=a_n+L^*w_n$ with
$q^*(a_n),\|w_n\|\le1$.  As $a_n(x)\le q(x)$, $w_n(Lx)\le\|Lx\|$ and the
sum tends to $p(x)$, we get $a_n(x)\to q(x)$.  Write $x/q(x)=z+Uh$ with
$z\in B_{c_0}$, $\|h\|\le1$.  Then
\[
 \frac{a_n(x)}{q(x)}\le\sum_j|a_n(j)||z_j|+\|U^*a_n\|
 \le1-\sum_j|a_n(j)|(1-|z_j|),
\]
so $\sum_j|a_n(j)|(1-|z_j|)\to0$.  The set $G=\{j:|z_j|\ge\frac12\}$ is
finite, $\sum_{j\notin G}|a_n(j)|\to0$, and the coordinates in $G$ are
bounded; hence every subsequence of $(a_n)$ has a norm convergent
subsequence, whose limit $a'$ satisfies $q^*(a')\le1$, $a'(x)=q(x)$, so
$a'=\nabla q(x)$.  Every subsequence of $(w_n)$ has a weak$^*$ convergent
subsequence $w_n\to w$, and $L^*w_n\to L^*w$ in norm because $L$ is
compact.  Then $\nabla q(x)+L^*w\in B_{p^*}$ norms $x$, so it equals
$\nabla p(x)$.  This proves (a), and the continuity in (b).

(b) If $f\in S_{p^*}$ attains its norm at $x\in S_p$, then $f=\nabla p(x)$.

(c) If $a\in S_{q^*}$ attains its norm at $y\in S_q$, write $y=z+Uh$,
$z\in B_{c_0}$, $\|h\|\le1$; then $1=a(z)+\ip{U^*a}{h}$ forces
$|z_j|=1$ on $\supp a$, which is finite because $z\in c_0$.  For the
pair $(a',z')$: $y'\in B_q$ and
$a'(y')=\|a'\|_1+\|U^*a'\|=1=q^*(a')$, hence $q(y')=1$ and
$\nabla q(y')=a'$; then (a) gives $\nabla p(y')$.  Conversely, if
$f'=\nabla p(x')$ with $x'\in S_p$, then $a':=\nabla q(x')\in c_{00}$ and
$y':=x'/q(x')=z'+Uh'$ with $z'\in B_{c_0}$; the equality
$1=a'(z')+\ip{U^*a'}{h'}\le\|a'\|_1+\|U^*a'\|=1$ forces $z'=\sgn a'$ on
$\supp a'$ and $h'=U^*a'/\|U^*a'\|$; finally $\nabla p(y')=\nabla p(x')$.
\end{proof}

\begin{proposition}[Continuity of the forced data]\label{prop:continuity}
Let $f_n\to f$ in $S_{p^*}$ \textup{(}not necessarily norm attaining\textup{)}
with forced data indexed by $n$.  Then $\xi_n\to\xi$ weak$^*$;
$R_m^{**}\xi_n\to R_m^{**}\xi$ in $\ell_1$; $w_{n,m}\to w_m$ weak$^*$ and
$D_mw_{n,m}\to D_mw_m$ in $\ell_2$, so $M_{n,m}\to M_m$, $C_{n,m}\to C_m$,
$\sigma_{n,m}\to\sigma_m$ and $\vartheta_{n,m}\to\vartheta_m$;
$R_m^*w_{n,m}\to R_m^*w_m$ and $L^*w_n\to L^*w$ in norm; $a_n\to a$ in
$\ell_1$; $e_n\to e$ in $H$; $q_0^{(n)}\to q_0$; $\zh_n\to\zh$ and
$z_n\to z$ weak$^*$.
\end{proposition}

\begin{proof}
A weak$^*$ cluster point $\eta$ of $(\xi_n)$ lies in $B_{p^{**}}$ and
$f(\eta)=\lim f_n(\xi_n)=1$, so $\eta=\xi$; $B_{p^{**}}$ is weak$^*$
metrizable.  $R_m$ is compact, so $R_m^{**}$ maps bounded weak$^*$
convergent sequences to norm convergent ones.  A weak$^*$ cluster point
$w'$ of $(w_{n,m})_n$ has $N_m(w')\le1$ and
$w'(\zeta_m)=\lim w_{n,m}(\zeta_{n,m})=\sigma_m$, so $w'=w_m$.  Since
$\Phi_m\in\ell_2$, $D_m$ maps bounded weak$^*$ convergent sequences to
norm convergent ones; $M=1-C$.  $L^*$ is weak$^*$-to-norm continuous on
bounded sets, and $a_n=f_n-L^*w_n$.  Finally
$q_0^{(n)}=a_n(\xi_n)\to a(\xi)$, $e_n=U^*a_n/\|U^*a_n\|$, and
$z_n=\xi_n/q_0^{(n)}-Ue_n$.
\end{proof}

\begin{proposition}[Norm-attaining approximants]\label{prop:approximants}
Let $f\in S_{p^*}$ have forced data as above.  For a sequence
$f'_n\in\NA((X,p),\R)\cap S_{p^*}$ the following are equivalent:
\begin{enumerate}
\item[(i)] $f'_n\to f$ in norm;
\item[(ii)] $f'_n=\nabla p(y'_n)$ with $y'_n=z'_n+U(U^*a'_n/\|U^*a'_n\|)$,
where $a'_n\in c_{00}\cap S_{q^*}$, $a'_n\to a$ in $\ell_1$, $z'_n\in c_0$,
$\|z'_n\|_\infty\le1$, $z'_n=\sgn a'_n$ on $\supp a'_n$, and $z'_n\to z$
coordinatewise.
\end{enumerate}
In particular, for every $n$ the coordinates of $z'_n$ outside a finite
window $W_n\supseteq\supp a'_n$ may be chosen freely \textup{(}moduli
$\le1$, membership in $c_0$\textup{)} as long as the windows exhaust $\N$;
the window must contain $\supp a'_n$, where $z'_n=\sgn a'_n$ is forced.
\end{proposition}

\begin{proof}
(ii)$\Rightarrow$(i): $y'_n\to\zh$ weak$^*$ (bounded coordinatewise
convergence plus norm convergence of the $U$-part, as $U^*a'_n/\|U^*a'_n\|\to
e$).  By compactness of $L$, $p(y'_n)=1+\|Ly'_n\|\to1+\|L^{**}\zh\|=1/q_0$,
so $x_n:=y'_n/p(y'_n)\to\xi$ weak$^*$.  Then $\nabla q(x_n)=a'_n\to a$, and
$L^*J_V(Lx_n)\to L^*w$ in norm by the argument of
Proposition~\ref{prop:continuity} (a weak$^*$ cluster point of $J_V(Lx_n)$
norms $L^{**}\xi$, whose components are nonzero).
(i)$\Rightarrow$(ii): write $f'_n=\nabla p(x'_n)$ (Proposition
\ref{prop:smooth}); by Proposition~\ref{prop:continuity}, $x'_n\to\xi$
weak$^*$, $a'_n:=\nabla q(x'_n)\to a$, and Proposition~\ref{prop:smooth}(c)
gives the form of $y'_n=x'_n/q(x'_n)$; finally $z'_n=y'_n-U(\cdot)\to z$
weak$^*$.
\end{proof}

\subsection{Mates and the reduction}
For $f\in S_{p^*}$ we use, as in \cite{PrA},
\[
 \Cset(f):=\{g\in X^*:\ p^*(f+tg)\le s(t)\ \text{for all }t\in\R\}.
\]

\begin{proposition}[Reduction]\label{prop:reduction}
\textup{(a)} For $f\in S_{p^*}$ and $g\in X^*$ the operator
$(f,g):(X,p)\to\ell_2^2$ is contractive iff $g\in\Cset(f)$ iff
$f(y)^2+g(y)^2\le p(y)^2$ for all $y\in X$.
\textup{(b)} Every $S\in\mathcal L((X,p),\ell_2^2)$ of norm one is, after a
rotation of $\ell_2^2$, of the form $(f,g)$ with $f\in S_{p^*}$ and
$g\in\Cset(f)$.  Hence $\NA((X,p),\ell_2^2)$ is dense iff
$(f,\rho g)\in\cl\NA((X,p),\ell_2^2)$ for all $f\in S_{p^*}$,
$g\in\Cset(f)$ and $\rho\in(0,1)$.
\textup{(c)} If $g\in\Cset(f)$ then $p^*(g)\le1$, $g(\xi)=0$, and
$p^*(f+t\rho g)\le s(\rho t)<s(t)$ for $t\ne0$, $\rho\in(0,1)$.
\textup{(d)} If $f$ attains its norm at $x$ and $g\in\Cset(f)$, then
$(f,g)$ attains its norm at $x$.
\end{proposition}

\begin{proof}
(a) $\|(f,g)\|=\sup_\theta p^*(\cos\theta f+\sin\theta g)$ and
$p^*(\cos\theta f+\sin\theta g)=|\cos\theta|p^*(f+\tan\theta\,g)$ with
$|\cos\theta|=s(\tan\theta)^{-1}$; the case $\cos\theta=0$ is the limit
$t\to\infty$.  The primal form is
$\|(f,g)\|^2=\sup_{p(y)\le1}f(y)^2+g(y)^2$.
(b) $\theta\mapsto p^*(S^*(\cos\theta,\sin\theta))$ attains its maximum
$1$; rotate that direction to $e_1$.  Since $\NA$ is a cone and
$(f,g)=\lim_{\rho\to1}(f,\rho g)$, the reduction follows.
(c) Divide $p^*(f+tg)\le s(t)$ by $|t|\to\infty$; and
$1+tg(\xi)=(f+tg)(\xi)\le s(t)$ for all $t$ forces $g(\xi)=0$.
(d) $g(x)^2\le p(x)^2-f(x)^2=0$ (this observation is from
\cite[Section~13]{KLMW}).
\end{proof}

\begin{proposition}[Primal description of the fibre]\label{prop:primal}
For $f\in S_{p^*}$ put $r_f:=\sqrt{p^2-f^2}$ on $X$ and let $\tilde r_f$ be
its largest sublinear minorant,
$\tilde r_f(x)=\inf\{\sum_jr_f(y_j):\sum_jy_j=x\}$ \textup{(}finite
families\textup{)}.  Then $\tilde r_f$ is a seminorm,
$\Cset(f)=\{g\in X^*:g\le\tilde r_f\}$, $\Cset(f)$ is convex, symmetric
and weak$^*$ compact, and $\max_{g\in\Cset(f)}g(x)=\tilde r_f(x)$ for
$x\in X$.
\end{proposition}

\begin{proof}
The infimum $\varphi$ is subadditive, positively homogeneous, even and
$\le r_f$, and every sublinear minorant $s$ of $r_f$ satisfies
$s(x)\le\sum s(y_j)\le\sum r_f(y_j)$, so $\varphi=\tilde r_f$.  A linear $g$
satisfies $g\le r_f$ iff $g\le\tilde r_f$, and as $r_f$ is even this is
$|g|\le r_f$, i.e.\ $g\in\Cset(f)$ by Proposition~\ref{prop:reduction}(a).
Given $x$, the Hahn--Banach theorem gives a linear $g\le\tilde r_f$ with
$g(x)=\tilde r_f(x)$, continuous because $\tilde r_f\le p$.
\end{proof}

\begin{definition}[Recovery]\label{def:recovery}
For subsets $C_n\subseteq X^*$ let
$\Li_nC_n:=\{g:\dist(g,C_n)\to0\}$ (Kuratowski lower limit; it is closed,
and convex if the $C_n$ are convex).  A mate $g\in\Cset(f)$ is
\emph{recovered along} a sequence $f_n\to f$ in $S_{p^*}$ if
$g\in\Li_n\Cset(f_n)$.  We write
\[
 \Rec:=\{f\in S_{p^*}:\ (f,g)\in\cl\NA((X,p),\ell_2^2)\text{ for every }
 g\in\Cset(f)\}
\]
for the \emph{recoverable set}.  By Proposition~\ref{prop:reduction},
$\NA((X,p),\ell_2^2)$ is dense iff $\Rec=S_{p^*}$.
\end{definition}

\begin{lemma}[Recoverability of a single pair]\label{lem:pair}
Let $f\in S_{p^*}$, $g\in\Cset(f)$ and $\rho\in(0,1)$.  Then
$(f,\rho g)\in\cl\NA((X,p),\ell_2^2)$ iff there are
$f_n\in\NA((X,p),\R)\cap S_{p^*}$ with $f_n\to f$ and
$\rho g\in\Li_n\Cset(f_n)$.
\end{lemma}

\begin{proof}
Suppose $g_n\in\Cset(f_n)$ and $g_n\to\rho g$.  Then $(f_n,g_n)$ attains
its norm (Proposition~\ref{prop:reduction}(d)) and converges to
$(f,\rho g)$.
Conversely let $S_k\in\NA$ converge to $T:=(f,\rho g)$; $\|T\|=1$.  For a
unit vector $y=(\cos\theta,\sin\theta)$ with $\sin\theta\ne0$ we have
$p^*(T^*y)<1$: if $\cos\theta\ne0$,
\[
 p^*(T^*y)=|\cos\theta|\,p^*(f+\tan\theta\,\rho g)\le|\cos\theta|\,
 s(\rho\tan\theta)<1,
\]
and $p^*(\rho g)\le\rho<1$ otherwise.  Let $S_k$ attain its
norm at $x_k\in S_p$ and $y_k:=S_kx_k/\|S_kx_k\|$.  Then
$p^*(T^*y_k)\ge\|S_k\|-\|S_k-T\|\to1$, so every limit point of $(y_k)$ is
$\pm e_1$; replacing $x_k$ by $-x_k$ we may assume $y_k\to e_1$.  Let $Q_k$
be the rotation with $Q_ky_k=e_1$; then $Q_k\to I$, and $Q_kS_k=(\tilde
f_k,\tilde g_k)$ attains its norm at $x_k$ with $\tilde
f_k(x_k)=\|S_k\|$, $\tilde g_k(x_k)=0$.  The functionals
$f_k:=\tilde f_k/\|S_k\|$ and $g_k:=\tilde g_k/\|S_k\|$ satisfy
$\|(f_k,g_k)\|=1=f_k(x_k)$, so $f_k\in\NA\cap S_{p^*}$, $g_k\in\Cset(f_k)$,
and $(f_k,g_k)\to(f,\rho g)$.
\end{proof}

\section{The lower-limit theorem}\label{sec:lowerlimit}

\begin{theorem}[Global compactness; {\cite{PrA}}]\label{thm:compact}
If $f_n\to f$ in $S_{p^*}$ and $g_n\in\Cset(f_n)$, then $(g_n)$ has a norm
convergent subsequence and every limit of such a subsequence lies in
$\Cset(f)$.  Consequently every $\Cset(f)$ is norm compact, $f\mapsto\Cset(f)$
is upper semicontinuous, and $\Cset(f)\cup\bigcup_n\Cset(f_n)$ is norm
compact.
\end{theorem}

\begin{proof}
$p^*(g_n)\le1$; pass to a weak$^*$ convergent subsequence $g_n\to g$.
Weak$^*$ lower semicontinuity of $p^*$ gives $g\in\Cset(f)$.  Let
$\Delta:=\limsup_n\|g_n-g\|_1$, realized along a further subsequence.  Fix
$t>0$ and write $f_n+tg_n=A_n+L^*W_n$ with $q^*(A_n),\|W_n\|\le s(t)$
(Lemma~\ref{lem:dualball}).  Passing to a subsequence, $L^*W_n\to k$ in
norm, $W_n\to W$ weak$^*$, and $A_n\to A:=f+tg-k$ weak$^*$.  Since
$k(\xi)=W(L^{**}\xi)\le s(t)(1-q_0)$ we get $A(\xi)\ge1-s(t)(1-q_0)$ and
$q^*(A)\ge A(\xi)/q_0$.  Moreover $U^*A_n\to U^*A$ in norm, and for
bounded coordinatewise convergent sequences in $\ell_1$ one has
$\|A_n\|_1-\|A_n-A\|_1\to\|A\|_1$.  Since
$t(g_n-g)=(A_n-A)+(L^*W_n-k)-(f_n-f)$,
\begin{align*}
 t\Delta&=\limsup_n\|A_n-A\|_1=\limsup_n(q^*(A_n)-q^*(A))\\
 &\le s(t)-\frac{1-s(t)(1-q_0)}{q_0}=\frac{s(t)-1}{q_0}.
\end{align*}
Letting $t\to0$ gives $\Delta=0$.  The remaining assertions follow
(for the union, a sequence either has infinitely many terms in one compact
fibre or a subsequence to which the first part applies).
\end{proof}

\begin{theorem}[Lower-limit theorem]\label{thm:lowerlimit}
Let $f_n\to f$ in $S_{p^*}$ \textup{(}not necessarily norm
attaining\textup{)} and $\rho\in[0,1]$.  The following are equivalent:
\begin{enumerate}
\item[(i)] $\rho\,\Cset(f)\subseteq\Li_n\Cset(f_n)$;
\item[(ii)] $\liminf_n\tilde r_{f_n}(x)\ge\rho\,\tilde r_f(x)$ for every
$x\in X$;
\item[(iii)] \textup{(ii)} holds for all $x$ in a dense subset of $X$.
\end{enumerate}
\end{theorem}

\begin{proof}
(i)$\Rightarrow$(ii): choose $g\in\Cset(f)$ with $g(x)=\tilde r_f(x)$
(Proposition~\ref{prop:primal}) and $g_n\in\Cset(f_n)$ with $g_n\to\rho
g$; then $\tilde r_{f_n}(x)\ge g_n(x)\to\rho\tilde r_f(x)$.
(ii)$\Leftrightarrow$(iii): each $\tilde r_h$ is a seminorm bounded by $p$,
so $|\tilde r_h(x)-\tilde r_h(y)|\le p(x-y)$ uniformly in $h$.
(ii)$\Rightarrow$(i): suppose $g\in\rho\Cset(f)$ and
$\dist(g,\Cset(f_{n_j}))\ge\delta>0$.  All $\Cset(f_{n_j})$ lie in the
compact set of Theorem~\ref{thm:compact}; by Blaschke's selection theorem
we may assume $\Cset(f_{n_j})\to A$ in the Hausdorff metric.  Then $A$ is
compact and convex and $\dist(g,A)\ge\delta$.  For $x\in X$ the support
functions satisfy $|h_C(x)-h_{C'}(x)|\le d_H(C,C')\|x\|_\infty$, so
$h_A(x)=\lim_j\tilde r_{f_{n_j}}(x)\ge\rho\tilde r_f(x)\ge g(x)$.  As $A$
is weak$^*$ compact and convex, $g\notin A$ would give $x\in c_0$ with
$g(x)>h_A(x)$.  Contradiction.
\end{proof}

\begin{corollary}[Finite, diagonal and mate versions]\label{cor:finite}
For $f\in S_{p^*}$ the following are equivalent:
\begin{enumerate}
\item[(a)] there are $f_n\in\NA\cap S_{p^*}$ with $f_n\to f$ and
$\Cset(f)\subseteq\Li_n\Cset(f_n)$;
\item[(b)] for every $\varepsilon>0$, $\rho<1$ and $x_1,\dots,x_k\in X$ there
is $f'\in\NA\cap S_{p^*}$ with $p^*(f'-f)<\varepsilon$ and
$\tilde r_{f'}(x_i)\ge\rho\tilde r_f(x_i)-\varepsilon$ for $i\le k$;
\item[(c)] for every $\varepsilon>0$, $\rho<1$ and $g_1,\dots,g_k\in\Cset(f)$
there is $f'\in\NA\cap S_{p^*}$ with $p^*(f'-f)<\varepsilon$ and
$\dist(\rho g_i,\Cset(f'))<\varepsilon$ for $i\le k$;
\item[(d)] as \textup{(c)}, uniformly:
$\sup_{g\in\Cset(f)}\dist(\rho g,\Cset(f'))<\varepsilon$.
\end{enumerate}
If they hold, then $f\in\Rec$.
\end{corollary}

\begin{proof}
(a)$\Rightarrow$(d): $\dist(g,\Cset(f_n))\to0$ for each $g$, uniformly on
the compact set $\Cset(f)$ (finite nets and $1$-Lipschitz dependence on
$g$); and $\dist(\rho g,\Cset(f_n))\le\rho\dist(g,\Cset(f_n))$ because
$\Cset(f_n)$ is convex and contains $0$.  (d)$\Rightarrow$(c) is trivial.
(c)$\Rightarrow$(b): pick $g_i\in\Cset(f)$ with $g_i(x_i)=\tilde r_f(x_i)$
and apply (c) with $\varepsilon/\max(1,\max_i\|x_i\|)$.
(b)$\Rightarrow$(a): let $(x_i)$ be dense in $X$; for each $n$ apply (b)
with $\varepsilon=1/n$, $\rho=1-1/n$ and $x_1,\dots,x_n$ to obtain $f_n$.
Then $\liminf_n\tilde r_{f_n}(x_i)\ge\tilde r_f(x_i)$ for every $i$, and
Theorem~\ref{thm:lowerlimit} with $\rho=1$ gives (a).  The
last claim is Lemma~\ref{lem:pair}.
\end{proof}

\begin{theorem}[Transitivity]\label{thm:transitivity}
$\NA\cap S_{p^*}\subseteq\Rec$.  Let $f\in S_{p^*}$ and suppose that for
every $\rho<1$ there is a sequence $f^\rho_n\in\Rec$ with $f_n^\rho\to f$
and $\liminf_n\tilde r_{f_n^\rho}(x)\ge\rho\,\tilde r_f(x)$ on a dense set
of $x\in X$.  Then $f\in\Rec$.
\end{theorem}

\begin{proof}
The first claim is Proposition~\ref{prop:reduction}(d).  By
Theorem~\ref{thm:lowerlimit}, $\rho\Cset(f)\subseteq\Li_n\Cset(f^\rho_n)$.
Given $g\in\Cset(f)$ choose $g_n\in\Cset(f_n^\rho)$ with $g_n\to\rho g$;
each $(f_n^\rho,g_n)$ lies in $\cl\NA$, hence so does $(f,\rho g)$; let
$\rho\to1$.
\end{proof}

\begin{remark}[Several rows and the residual set]\label{rem:joint}
For $d\ge1$ let
$\Cset_d(f):=\{(g_1,\dots,g_d):f(y)^2+\sum_ig_i(y)^2\le p(y)^2\ (y\in X)\}$;
equivalently $p^*(f+\sum_it_ig_i)\le s(|t|)$ for $t\in\R^d$.
Theorem~\ref{thm:compact}, Theorem~\ref{thm:lowerlimit},
Corollary~\ref{cor:finite} and Theorem~\ref{thm:transitivity} hold verbatim
for $\Cset_d$ (support functions on $X^d$, separation by $(c_0)^d$), with
$\ell_2^{d+1}$ in place of $\ell_2^2$.  By the residual uniform recovery
theorem of \cite{PrA} there
is a dense $G_\delta$ set $\mathcal O\subseteq S_{p^*}$ at which every $\Cset_d$
is Hausdorff continuous; in particular $\mathcal O\subseteq\Rec$.  This does not
decide density: Hausdorff continuity may fail on a meagre set of first
rows, and the remaining sections study what can be recovered there.
\end{remark}

\section{Certificates and the transport theorem}\label{sec:certificates}

Throughout this section $f\in S_{p^*}$ has forced data as in
Section~\ref{sec:setting}; block indices are dropped when a single block is
involved.  For $b\in\ell_1$ we write $\|b/a\|_\infty:=\sup_{j\in F}|b_j|/|a_j|$
(when $\supp b\subseteq F$), and $P^\perp$ denotes the orthogonal projection
of $H$ onto $e^\perp$ (throughout the paper; it is not used for other
projections).

\subsection{Definitions and exact expansions}

\begin{definition}[Finite certificates]\label{def:certificate}
A \emph{finite certificate} at $f$ is a pair $c=(b,\omega)$,
$\omega=(\omega_m)_{m\in I}$, such that
\begin{enumerate}
\item[(C1)] $b\in\ell_1$, $\supp b\subseteq F$, $\|b/a\|_\infty<\infty$ and
$b(\zh)=0$;
\item[(C2)] each $\omega_m$ is finitely supported in $Q_m$, and
$\omega_m=0$ for all but finitely many $m$ (the \emph{active} blocks).
\end{enumerate}
Put $d_m:=\ip{D_mw_m}{D_m\omega_m}/C_m$ and define the \emph{direction}
$g_c:=b+\sum_mR_m^*(\omega_m-d_mw_m)$, the coefficients
\begin{gather*}
 h(b):=\frac{\|P^\perp U^*b\|^2}{\nu},\qquad
 H_m(\omega_m):=\frac{\|D_m\omega_m\|_2^2-d_m^2}{C_m},\\
 H(c):=\max\big(h(b),\max_mH_m(\omega_m)\big),
\end{gather*}
the constant $\kappa(c):=\max\{2\beta/\nu,\ \max_m2|d_m|M_m/C_m\}$ with
$\beta:=\|U^*b\|$, and the \emph{coordinatewise radius}
\[
 r(c):=\min\Big\{\frac1{\|b/a\|_\infty},\ \frac\nu{2\beta},\
 \min_m\min\Big(\min_{k\in\supp\omega_m}\frac{\gap_m(k)}{2|\omega_m(k)|},\
 \frac1{2|d_m|},\ \frac{C_m}{2|d_m|M_m}\Big)\Big\}
\]
(with $1/0=\infty$).  Finally
\[
 \Cert(f):=\{g_c:\ c\text{ a finite certificate at }f,\ H(c)\le1,\ g_c\in\Cset(f)\}.
\]
\end{definition}

Let $P_m^\perp$ denote the orthogonal projection of $\ell_2$ onto
$(D_mw_m)^\perp$.  Then
\[
 \|P_m^\perp D_m\omega_m\|^2=\|D_m\omega_m\|^2-d_m^2,
\]
so $H_m\ge0$.  The radius is the
\emph{coordinatewise} one: each off-peak coordinate is compared with its own
gap.  (A cruder radius $\min_k\gap(k)/(2\max_k|\omega(k)|)$ would not suffice
for Theorem~\ref{thm:weighted} below.)

\begin{lemma}[Algebra]\label{lem:algebra}
Finite certificates form a vector space and $c\mapsto g_c$ is linear; $H$ is
convex and $H(\lambda c)=\lambda^2H(c)$; $r(\lambda c)=r(c)/|\lambda|$,
$\kappa(\lambda c)=|\lambda|\kappa(c)$, and $\kappa(c)r(c)\le1$.  Every
direction satisfies $g_c(\xi)=0$, and $\Cert(f)$ is convex and symmetric.
\end{lemma}

\begin{proof}
$d_m$ is linear in $\omega_m$, and $H$ is a maximum of positive
semidefinite quadratic forms.  The inequality $\kappa r\le1$ follows from
the terms $\nu/(2\beta)$ and $C_m/(2|d_m|M_m)$ of the radius.  We have
$b(\xi)=q_0b(\zh)=0$ and, by Lemma~\ref{lem:threshold},
\begin{align*}
 \frac{\ip{\omega_m-d_mw_m}{\zeta_m}}{\sigma_m}
 &=\ip{\omega_m}{\alpha_m}+\frac{\ip{D_m\omega_m}{D_mw_m}}{C_m}
 -d_m\Big(\ip{w_m}{\alpha_m}+\frac{\|D_mw_m\|^2}{C_m}\Big)\\
 &=0+d_m-d_m(M_m+C_m)=0,
\end{align*}
because $\alpha_m$ lives on $P_m$ and $\omega_m$ vanishes there.  Convexity
of $\Cert(f)$ follows from linearity, convexity of $H$ and of $\Cset(f)$.
\end{proof}

\begin{lemma}[Exact base identity]\label{lem:base}
Put $\Psi(h):=\|e+h\|-1-\ip{e}{h}$ for $h\in H$.  For $B\in\ell_1$ and
$t\in\R$,
\begin{equation}\label{eq:baseidentity}
 q^*(a+tB)=1+tB(\zh)+\mathrm{Fl}_B(t)+\sum_{j\notin F}\big(|tB_j|-z_jtB_j\big)
 +\nu\Psi(tU^*B/\nu),
\end{equation}
where $\mathrm{Fl}_B(t):=\sum_{j\in F}2(-\sgn(a_j)tB_j-|a_j|)_+$; all terms
after $tB(\zh)$ are nonnegative.  Moreover
$\Psi(h)=\|h_\perp\|^2/(\|e+h\|+1+\ip eh)$ with $h_\perp=P^\perp h$, and
for $\|h\|\le\frac12$:
$\|h_\perp\|^2/(2(1+\|h\|))\le\Psi(h)\le\|h_\perp\|^2/(2(1-\|h\|))$.
Consequently:
\begin{enumerate}
\item[(a)] if $B(\zh)=0$ and $|t|\,\|U^*B\|\le\nu/2$, then
$q^*(a+tB)-1\ge\frac{t^2}{2}h(B)/(1+|t|\|U^*B\|/\nu)$;
\item[(b)] if $b$ satisfies \textup{(C1)} and
$|t|\le\min(1/\|b/a\|_\infty,\nu/(2\beta))$, then
$q^*(a+tb)\le1+\frac{t^2}2h(b)(1+2|t|\beta/\nu)$.
\end{enumerate}
\end{lemma}

\begin{proof}
For real $x\ne0$ and $y$: $|x+y|=|x|+\sgn(x)y+2(-\sgn(x)y-|x|)_+$.  Summing
over $j\in F$ (where $z_j=\sgn a_j$) gives
$\sum_{j\in F}|a_j+tB_j|=\|a\|_1+\sum_{j\in F}z_jtB_j+\mathrm{Fl}_B(t)$,
and $\sum_{j\notin F}|tB_j|=\sum_{j\notin F}z_jtB_j+\sum_{j\notin
F}(|tB_j|-z_jtB_j)$.  Next $\|U^*a+tU^*B\|=\nu\|e+h\|$ with $h=tU^*B/\nu$
and $\nu\ip eh=t\ip{U^*B}{e}$.  Adding and using $\|a\|_1+\nu=1$ and
$B(z)+\ip{U^*B}{e}=B(\zh)$ gives \eqref{eq:baseidentity}.  Nonnegativity of
$\Psi$: $\|e+h\|\ge\ip{e}{e+h}$.  The formula for $\Psi$ follows from
$\|e+h\|^2-(1+\ip eh)^2=\|h\|^2-\ip eh^2$; for $\|h\|\le\frac12$ the
denominator lies in $[2(1-\|h\|),2(1+\|h\|)]$.  With $h=tU^*B/\nu$ one has
$\nu\|h_\perp\|^2=t^2h(B)$; (a) follows by dropping the nonnegative terms,
and (b) because no flip occurs for $|t|\le1/\|b/a\|_\infty$, the sum over
$j\notin F$ vanishes, and $1/(1-x)\le1+2x$ for $x\le\frac12$.
\end{proof}

\begin{lemma}[Block expansion]\label{lem:block}
Fix a block.  Let $\omega:\N\to\R$ vanish on $P$ with $D\omega\in\ell_2$,
$d:=\ip{Dw}{D\omega}/C$, $h_\perp:=D\omega-(d/C)Dw$ \textup{(}so
$\|h_\perp\|^2=\|D\omega\|^2-d^2$\textup{)}, $W(\sigma):=(1-d\sigma)w+
\sigma\omega$ and $Y:=C+d\sigma M$.  Then:
\begin{enumerate}
\item[(a)] $\|DW(\sigma)\|_2=\sqrt{Y^2+\sigma^2\|h_\perp\|^2}$;
\item[(b)] if $1-d\sigma\ge0$ and $Y>0$, then
$N(W(\sigma))\ge1+\sigma^2\|h_\perp\|^2/(\sqrt{Y^2+\sigma^2\|h_\perp\|^2}+Y)$;
\item[(c)] if $|d\sigma|\le\frac12$ and $|\sigma\omega(k)|\le\gap(k)/2$ for
all $k$ with $\omega(k)\ne0$, then $\|W(\sigma)\|_\infty$ equals
$(1-d\sigma)M$; if moreover $Y>0$, equality holds in \textup{(b)};
\item[(d)] if $\omega\in c_{00}$ and $|\sigma|\le\min\big(\min_{k}
\gap(k)/(2|\omega(k)|),1/(2|d|),C/(2|d|M)\big)$, then
$N(W(\sigma))\le1+\frac{\sigma^2}2H(\omega)(1+2|d\sigma|M/C)$.
\end{enumerate}
\end{lemma}

\begin{proof}
(a) Using $1/C-1=M/C$, $(1-d\sigma)Dw+\sigma D\omega=(1+d\sigma M/C)Dw
+\sigma h_\perp$, and $h_\perp\perp Dw$.  (b) At a peak $k$,
$|W(\sigma)(k)|=(1-d\sigma)M$, and $(1-d\sigma)M+Y=M+C=1$.  (c) Off
$\supp\omega$, $|W(k)|=(1-d\sigma)|w(k)|$; on $\supp\omega$,
$|W(k)|\le(1-d\sigma)(M-\gap(k))+\gap(k)/2\le(1-d\sigma)M$.  (d) Here $Y\ge
C/2$, the bracket in (c) is at most $\sigma^2\|h_\perp\|^2/(2Y)$, and
$C/Y\le1/(1-|d\sigma|M/C)\le1+2|d\sigma|M/C$.
\end{proof}

\begin{proposition}[Certificate expansion]\label{prop:expansion}
For a finite certificate $c$ at $f$ and $|\sigma|\le r(c)$,
\[
 p^*(f+\sigma g_c)\le1+\frac{\sigma^2}2H(c)\big(1+\kappa(c)|\sigma|\big).
\]
\end{proposition}

\begin{proof}
$f+\sigma g_c=(a+\sigma b)+\sum_mR_m^*W_m(\sigma)$ with
$W_m(\sigma)=(1-d_m\sigma)w_m+\sigma\omega_m$ ($=w_m$ on inactive blocks).
Apply Lemma~\ref{lem:dualball}, Lemma~\ref{lem:base}(b) and
Lemma~\ref{lem:block}(d).
\end{proof}

\begin{lemma}[Slack and validity window]\label{lem:slack}
\textup{(a)} For $\rho\in(0,1)$ and real $t$,
$s(t)-s(\rho t)\ge(1-\rho^2)\min(t^2,|t|)/3$.
\textup{(b)} If $H(c)\le1-\delta$ with $0<\delta\le1$ and
$|\sigma|\le r_*(c,\delta):=\min(r(c),\delta/(2\kappa(c)+2),\sqrt\delta)$,
then $p^*(f+\sigma g_c)\le s(\sigma)$.
\end{lemma}

\begin{proof}
(a) $s(t)-s(\rho t)=(1-\rho^2)t^2/(s(t)+s(\rho t))\ge(1-\rho^2)t^2/(2s(t))$
and $2s(t)\le3\max(1,|t|)$.  (b) $(1-\delta)(1+\kappa|\sigma|)\le
(1-\delta)(1+\delta/2)\le1-\delta/2$, and $1+\frac{\sigma^2}2(1-\frac\delta
2)\le1+\frac{\sigma^2}2-\frac{\sigma^4}8\le s(\sigma)$ since
$\sigma^2\le\delta$ and $\sqrt{1+x}\ge1+\frac x2-\frac{x^2}8$ for $x\ge0$.
\end{proof}

\begin{proposition}[Scale criterion]\label{prop:scale}
Let $f,f'\in S_{p^*}$, $g\in\Cset(f)$, $\rho\in(0,1)$, and let $c'$ be a
finite certificate at $f'$ with $H(c')\le1-\delta$, $0<\delta\le1$.  Put
$r_*:=r_*(c',\delta)$, $\varepsilon:=p^*(f'-f)$ and
$\eta:=p^*(g_{c'}-\rho g)$.  If $\varepsilon\le(1-\rho^2)r_*^2/6$ and
$\eta\le(1-\rho^2)r_*/6$, then $g_{c'}\in\Cert(f')$ and
$\dist_{p^*}(\rho g,\Cset(f'))\le\eta$.
\end{proposition}

\begin{proof}
For $|t|\le r_*$ use Lemma~\ref{lem:slack}(b) at $f'$.  For $|t|\ge r_*$,
$p^*(f'+tg_{c'})\le p^*(f+t\rho g)+\varepsilon+|t|\eta\le
s(\rho t)+\varepsilon+|t|\eta$; for $r_*\le|t|\le1$ both $\varepsilon$
and $|t|\eta$ are at most $(1-\rho^2)t^2/6$, and for $|t|\ge1$ their sum
is at most $(1-\rho^2)|t|/3$; conclude with Lemma~\ref{lem:slack}(a).
\end{proof}

Thus the slack alone certifies the scales $|t|\gtrsim\sqrt{\varepsilon/(1-
\rho^2)}$, and the smaller scales must be certified by the local structure
of $f'$; this is the multi-scale nature of the problem.

\subsection{Transport along every sequence}

\begin{theorem}[Transport]\label{thm:transport}
Let $f_n\to f$ in $S_{p^*}$ be an arbitrary sequence \textup{(}norm
attaining or not\textup{)} and let $c=(b,\omega)$ be a finite certificate
at $f$.  Put
\begin{gather*}
 G_n:=\{j\in F:\ \sgn a_n(j)=\sgn a_j,\ |a_n(j)|\ge|a_j|/2\},\\
 b'_n:=b\,1_{G_n},\qquad \tau_n:=b'_n(\zh_n),\qquad b_n:=b'_n-\tau_na_n,
\end{gather*}
and $c_n:=(b_n,\omega)$ \textup{(}the same block coefficients; $d_{n,m}$ is
computed at $f_n$\textup{)}.  Then for $n$ large $c_n$ is a finite
certificate at $f_n$, $\|g_{c_n}-g_c\|\to0$, $H(c_n)\to H(c)$,
$\kappa(c_n)\to\kappa(c)$ and $\liminf_nr(c_n)\ge r(c)/2$.  If moreover
$g_c\in\Cset(f)$ and $H(c)\le1$, then for every $\rho\in(0,1)$,
$\rho g_{c_n}\in\Cert(f_n)$ for all large $n$.
\end{theorem}

\begin{proof}
\emph{Base.}  By Proposition~\ref{prop:continuity}, $a_n\to a$ in $\ell_1$,
$\nu_n\to\nu$, $e_n\to e$ and $z_n\to z$ weak$^*$; also $z_n=\sgn a_n$ on
$\supp a_n$.  Every $j\in F$ eventually lies in $G_n$, so $b'_n\to b$ in
$\ell_1$ by dominated convergence.  On $G_n$,
$z_n(j)=\sgn a_n(j)=\sgn a_j=z_j$, hence
$\tau_n=\sum_{j\in G_n}b_jz_j+\ip{U^*b'_n}{e_n}\to b(z)+\ip{U^*b}{e}=
b(\zh)=0$, and $b_n\to b$.  Further $\supp b_n\subseteq\supp a_n$,
$\|b_n/a_n\|_\infty\le2\|b/a\|_\infty+|\tau_n|$, and
$b_n(\zh_n)=\tau_n-\tau_na_n(\zh_n)=0$.  Thus (C1) holds at $f_n$,
$h_n(b_n)\to h(b)$ and $\|U^*b_n\|\to\beta$.

\emph{Blocks.}  For the finitely many active $m$ and $k\in\supp\omega_m$,
$w_{n,m}(k)\to w_m(k)$ and $M_{n,m}\to M_m$, so
$M_{n,m}-|w_{n,m}(k)|\to\gap_m(k)>0$: eventually $\supp\omega_m\subseteq
Q_{n,m}$ and the gaps converge.  Also $d_{n,m}\to d_m$, $C_{n,m}\to C_m$,
$D_mw_{n,m}\to D_mw_m$, so $H(c_n)\to H(c)$, $\kappa(c_n)\to\kappa(c)$ and
$\liminf r(c_n)\ge r(c)/2$ (the base ratio is at most halved).

\emph{Directions.}  $g_{c_n}-g_c=(b_n-b)+\sum_m(d_mR_m^*w_m-d_{n,m}R_m^*w_{n,m})
\to0$, since $R_m^*w_{n,m}\to R_m^*w_m$ in norm.

\emph{Mates.}  Let $\delta:=(1-\rho^2)/2$.  Then $H(\rho c_n)=\rho^2H(c_n)
\to\rho^2H(c)\le1-2\delta$, so $H(\rho c_n)\le1-\delta$ for large $n$, and
$r_*(\rho c_n,\delta)$ is bounded below by some $r_0>0$.  Apply
Proposition~\ref{prop:scale} with $f'=f_n$, $c'=\rho c_n$ and $g=g_c$:
$p^*(f_n-f)\to0$ and $p^*(g_{\rho c_n}-\rho g_c)=\rho p^*(g_{c_n}-g_c)\to0$.
\end{proof}

\begin{corollary}[Lower semicontinuity of the certified fibre]\label{cor:lsc}
For every sequence $f_n\to f$ in $S_{p^*}$,
\[
 \cl\Cert(f)\subseteq\Li_n\cl\Cert(f_n)\subseteq\Li_n\Cset(f_n).
\]
Consequently: \textup{(a)} every $g\in\cl\Cert(f)$ is recovered along every
norm-attaining sequence $f_n\to f$, finitely many such mates simultaneously,
and $(f,g)$ lies in $\cl\NA((X,p),\ell_2^2)$; \textup{(b)} if
$\Cset(f)=\cl\Cert(f)$, then $f\in\Rec$ and $\Cset$ is Hausdorff continuous
at $f$; \textup{(c)} if for every $f$ the set
$\Cset(f)\setminus\cl\Cert(f)$ is contained in $\Li_n\Cset(f_n)$ for \emph{some} norm-attaining sequence
$f_n\to f$, then $\NA((X,p),\ell_2^2)$ is dense.
\end{corollary}

\begin{proof}
By Theorem~\ref{thm:transport}, $\rho\Cert(f)\subseteq\Li_n\Cert(f_n)$ for
each $\rho<1$; lower limits are closed.  (a) follows from
Lemma~\ref{lem:pair} and Corollary~\ref{cor:finite} (applied to finitely
many mates).  (b) Upper semicontinuity (Theorem~\ref{thm:compact}) and
$\Cset(f)\subseteq\Li\Cset(f_n)$ give Kuratowski, hence Hausdorff,
convergence inside a compact set.  (c) Then $\Cset(f)\subseteq
\Li\Cset(f_n)$ for that sequence (the lower limit is closed and convex);
apply Corollary~\ref{cor:finite}.
\end{proof}

\begin{corollary}[Base and block directions]\label{cor:check}
\textup{(i)} If $\omega$ is finitely supported strictly below the peak of
$w_m$, $d=\ip{D_mw_m}{D_m\omega}/C_m$, $g=R_m^*(\omega-dw_m)\in\Cset(f)$
and $H_m(\omega)\le1$, then $g\in\Cert(f)$; in particular $(f,g)$ is
recovered along every sequence.  This is the finite local-certificate
theorem quoted from the unavailable note \cite{Check} in \cite{PrA}.
\textup{(ii)} If $b$ satisfies \textup{(C1)}, then $b/\Lambda\in\Cert(f)$
for every
\[
 \Lambda\ge\max\big(1,\sqrt{2h(b)},p^*(b)r_0/(s(r_0)-1)\big),
 \qquad r_0:=\min\big(r(c),1/(4\kappa(c)+4),1/\sqrt2\big),
\]
where $c=(b,0)$.
\end{corollary}

\begin{proof}
(i) is the definition of $\Cert(f)$.  (ii) $H(c/\Lambda)=h(b)/\Lambda^2\le
\frac12$, $\kappa(c/\Lambda)\le\kappa(c)$ and $r(c/\Lambda)\ge r(c)$, so
$r_*(c/\Lambda,\frac12)\ge r_0$ and Lemma~\ref{lem:slack}(b) gives
$p^*(f+\sigma b/\Lambda)\le s(\sigma)$ for $|\sigma|\le r_0$.  For
$|\sigma|\ge r_0$, $p^*(f+\sigma b/\Lambda)\le1+|\sigma|p^*(b)/\Lambda\le
s(\sigma)$, since $(s(\sigma)-1)/|\sigma|$ is increasing in $|\sigma|$.
\end{proof}

\subsection{Shifted certificates}
Since $f=a+L^*w$, one may move a second-order amount of mass between the
base and the blocks: $f=\lambda(a+L^*w)-(\lambda-1)f$ with
$\lambda=1+\theta\tau^2$.

\begin{definition}[Shifted certificates]\label{def:shifted}
A \emph{shifted certificate} is $c=(b,\omega,\theta)$, where $(b,\omega)$ is
a finite certificate and $\theta=(\theta_m)$ is a finitely supported real
family.  Put $v_\theta:=\sum_m\theta_mR_m^*w_m\in\ell_1$,
$\kappa_q(v):=\sum_{j\notin F}(|v_j|+z_jv_j)\in[0,2\|v\|_1]$ and
\[
 H^{\rm sh}(c):=\max\Big(h(b)-\frac2{q_0}\sum_m\theta_m\sigma_m
 +2\kappa_q(v_\theta),\ \max_m\big(H_m(\omega_m)+2\theta_m\big)\Big).
\]
The direction is $g_c:=g_{(b,\omega)}$, and
$\Cert^{\rm sh}(f):=\{g_c\in\Cset(f):\ H^{\rm sh}(c)\le1\}$.
\end{definition}

Note $v_\theta(\zh)=\sum_m\theta_m\ip{w_m}{\zeta_m}/q_0=\sum_m\theta_m
\sigma_m/q_0$.  The associated decomposition is
\begin{equation}\label{eq:shiftdec}
 \begin{gathered}
 A(\tau):=a+\tau b-\tau^2v_\theta,\qquad W_m(\tau):=(1+\theta_m\tau^2)w_m+
 \tau(\omega_m-d_mw_m),\\
 A(\tau)+L^*W(\tau)=f+\tau g_c.
 \end{gathered}
\end{equation}
For $\rho>0$ put $\rho.c:=(\rho b,\rho\omega,\rho^2\theta)$; then
$g_{\rho.c}=\rho g_c$ and $H^{\rm sh}(\rho.c)=\rho^2H^{\rm sh}(c)$.
$H^{\rm sh}$ is jointly convex, so $\Cert^{\rm sh}(f)$ is convex,
symmetric and contains $\Cert(f)$ ($\theta=0$).

\begin{proposition}[Shifted expansion]\label{prop:shifted}
For a shifted certificate $c$ at $f$ there is $\varepsilon_c(\tau)\to0$ with
$p^*(f+\tau g_c)\le1+\frac{\tau^2}2(H^{\rm sh}(c)+\varepsilon_c(\tau))$.
\end{proposition}

\begin{proof}
\emph{Blocks.} $W_m(\tau)=(1+\theta_m\tau^2)[w_m+\tau'(\omega_m-d_mw_m)]$
with $\tau'=\tau/(1+\theta_m\tau^2)$, so Lemma~\ref{lem:block}(d) gives
$N_m(W_m(\tau))\le1+\tau^2(\theta_m+H_m/2)+O(|\tau|^3)$.
\emph{Base.} Apply \eqref{eq:baseidentity} with $t=\tau$,
$B=b-\tau v_\theta$: $\tau B(\zh)=-\tau^2v_\theta(\zh)$; off $F$,
$|\tau B_j|-z_j\tau B_j=\tau^2(|v_j|+z_jv_j)$; for $|\tau|\le
1/(2\|b/a\|_\infty)$ each flip term is at most
$2(\tau^2|v_j|-|a_j|/2)_+$, so
$\mathrm{Fl}\le2\tau^2\sum_{j\in F,\,|v_j|>|a_j|/(2\tau^2)}|v_j|=o(\tau^2)$;
and $\nu\Psi(\tau U^*B/\nu)=\frac{\tau^2}2(h(b)+O(\tau))$.  Hence
$q^*(A(\tau))\le1+\frac{\tau^2}2\big(h(b)-2v_\theta(\zh)+2\kappa_q(v_\theta)
\big)+o(\tau^2)$.  Conclude with Lemma~\ref{lem:dualball}.
\end{proof}

\begin{theorem}[Shifted transport]\label{thm:shifted}
Let $f_n\to f$ in $S_{p^*}$ be arbitrary, $c=(b,\omega,\theta)$ a shifted
certificate at $f$, and $c_n:=(b_n,\omega,\theta)$ with $b_n$ as in
Theorem~\ref{thm:transport}.  For every $\eta>0$ there are $\tau_\eta>0$ and
$n_\eta$ with $p^*(f_n+\tau g_{c_n})\le1+\frac{\tau^2}2(H^{\rm sh}(c)+\eta)$
for $n\ge n_\eta$, $|\tau|\le\tau_\eta$; moreover $\limsup_nH^{\rm
sh}_n(c_n)\le H^{\rm sh}(c)$ \textup{(}computed at $f_n$\textup{)}.
Consequently, if $g_c\in\Cset(f)$ and $H^{\rm sh}(c)\le1$, then
$\rho g_{c_n}\in\Cert^{\rm sh}(f_n)$ for every $\rho<1$ and $n$ large, and
\[
 \cl\Cert^{\rm sh}(f)\subseteq\Li_n\cl\Cert^{\rm sh}(f_n)\subseteq
 \Li_n\Cset(f_n)\quad\text{for every sequence }f_n\to f.
\]
\end{theorem}

\begin{proof}
The blocks are treated as in Proposition~\ref{prop:shifted}, with the
uniform bounds of Theorem~\ref{thm:transport}.  For the base let
$v_n:=\sum_m\theta_mR_m^*w_{n,m}\to v_\theta$ in $\ell_1$ and apply
\eqref{eq:baseidentity} at $f_n$ with $B=b_n-\tau v_n$.  For
$|\tau|\le1/(2\|b_n/a_n\|_\infty)$ and $|\tau\tau_n|\le1$ one checks
coordinatewise: for $j\in G_n$ the flip term is at most
$2\tau^2|v_{n,j}|1[|v_{n,j}|>|a_j|/(4\tau^2)]$; for $j\in\supp a_n\setminus
G_n$ one has $b_n(j)=-\tau_na_n(j)$ and the flip term is at most
$\tau^2(|v_{n,j}|+z_n(j)v_{n,j})$; off $\supp a_n$ the kink term equals
$\tau^2(|v_{n,j}|+z_n(j)v_{n,j})$.  Hence the flip and kink terms are at
most $\tau^2$ times
\[
 2\sum_{j\in F}|v_{n,j}|1\big[|v_{n,j}|>\tfrac{|a_j|}{4\tau^2}\big]
 +\sum_{j\notin F}(|v_{n,j}|+z_n(j)v_{n,j})+2\sum_{j\in F\setminus
 G_n}|v_{n,j}|.
\]
As $n\to\infty$ and $\tau\to0$ the first and third sums tend to $0$
(compare $v_n$ with $v_\theta$ in $\ell_1$, use dominated convergence and
that each $j\in F$ eventually lies in $G_n$), and the second tends to
$\kappa_q(v_\theta)$ ($z_n\to z$ coordinatewise, $|z_n|\le1$).  Also
$v_n(\zh_n)\to v_\theta(\zh)$, $\sigma_{n,m}\to\sigma_m$,
$q_0^{(n)}\to q_0$, and the Hilbert term is
$\frac{\tau^2}2(h_n(b_n)+O(\tau))$ with $h_n(b_n)\to h(b)$.  This gives the
uniform estimate, and the same computation without $\tau$ gives the
$\limsup$ of $H^{\rm sh}_n(c_n)$.  For the mates, fix $\rho<1$ and
$\eta\le(1-\rho)/\rho$; for $|\tau|\le t_1:=\min(\tau_\eta/\rho,
2\sqrt{1-\rho})$ we get $p^*(f_n+\tau\rho g_{c_n})\le1+\rho\tau^2/2\le
s(\tau)$, and for $|\tau|\ge t_1$ we use the slack
(Lemma~\ref{lem:slack}(a)) with $p^*(f_n-f)\to0$ and
$p^*(g_{c_n}-g_c)\to0$.  Finally $H^{\rm sh}_n(\rho.c_n)=\rho^2H^{\rm
sh}_n(c_n)\le1$ for large $n$.
\end{proof}

\begin{remark}\label{rem:shiftopen}
Shifts relax the condition $H\le1$ to $H^{\rm sh}\le1$.  Whether
$\cl\Cert^{\rm sh}(f)$ is strictly larger than $\cl\Cert(f)$ for some $f$ in
Mart\'in's space is \emph{open}: no mate with $H>1\ge H^{\rm sh}$ outside
$\cl\Cert(f)$ has been exhibited (finite-dimensional models suggest that
shifts help, but they cannot distinguish $\Cert(f)$ from its closure).
Section~\ref{sec:second} identifies the sharp second-order invariant at
tame supports, $\Gamma_w$, and one checks that $\min_\theta H^{\rm
sh}(b,\omega,\theta)\ge\Gamma_w$ (Remark~\ref{rem:gammaw}).
\end{remark}

\subsection{Several rows}

\begin{theorem}[Joint fibres]\label{thm:joint}
Let $c_1,\dots,c_d$ be finite certificates at $f$ with
$\mathbf g:=(g_{c_1},\dots,g_{c_d})\in\Cset_d(f)$ and
$\sup_{\theta\in S^{d-1}}H(\sum_i\theta_ic_i)\le1$.  Then for every sequence
$f_n\to f$ in $S_{p^*}$ and every $\rho<1$ the transported tuple
$\rho\mathbf g_n:=(\rho g_{c_{1,n}},\dots,\rho g_{c_{d,n}})$ lies in $\Cset_d(f_n)$
for $n$ large, and $\mathbf g_n\to\mathbf g$.  In particular, if the $f_n$ attain their
norms, then $(f,\rho\mathbf g)$ is a limit of norm-attaining operators into
$\ell_2^{d+1}$.
\end{theorem}

\begin{proof}
The construction of Theorem~\ref{thm:transport} is linear in $c$ (the index sets
$G_n$ of Theorem~\ref{thm:transport} do not depend on $b$), so $c(\theta):=\sum_i\theta_ic_i$ is
transported to $c_n(\theta)=\sum_i\theta_ic_{i,n}$.  On $S^{d-1}$ the
coefficients $H(c_n(\theta))$ converge to $H(c(\theta))$ uniformly (they
are maxima of quadratic forms in $\theta$ with convergent coefficients),
$\kappa(c_n(\theta))$ is bounded, and $r(c_n(\theta))\ge r_0>0$ uniformly
for $n$ large, because only the finitely many coordinates of
$\bigcup_i\supp\omega_i$ and the bounded quantities
$\|b_{i,n}/a_n\|_\infty$, $\|U^*b_{i,n}\|$, $|d_{i,n,m}|$ enter.  Now run
the proof of Proposition~\ref{prop:scale} along each ray $t=|t|\theta$:
for $|t|\le r_0$ use Lemma~\ref{lem:slack}(b) for $\rho c_n(\theta)$; for
$|t|\ge r_0$ use $p^*(f+\rho\sum_it_ig_{c_i})\le s(\rho|t|)$ and the slack
in $|t|$.  If $f_n$ attains its norm at $x_n$, then each row of $\rho\mathbf g_n$
vanishes at $x_n$ and $(f_n,\rho\mathbf g_n)$ attains its norm there.
\end{proof}

\section{Averaging over scales}\label{sec:averaging}

\subsection{The averaging theorem}

\begin{theorem}[Weighted averaging]\label{thm:averaging}
Let $g\in\Cset(f)$ and $\rho\in(0,1)$.  Let $c_1,\dots,c_n$ be finite
certificates at $f$ and $\mu_1,\dots,\mu_n\ge0$ weights with
$\sum_i\mu_i=1$; put $r_i:=r(c_i)$ and $\epsilon_i:=p^*(g-g_{c_i})$.  Let
$s_0\in(0,\min(1,\sqrt{1-\rho^2})]$ and $\eta_0,\kappa_0\ge0$ satisfy
$\rho^2(1+\eta_0)(1+\kappa_0)\le(1+\rho^2)/2$, and assume
\begin{enumerate}
\item[(i)] $H(c_i)\le1+\eta_0$ and $\kappa(c_i)\min(r_i,s_0)\le\kappa_0$
for all $i$;
\item[(ii)] $\sum_{i:\,r_i<\rho s}\mu_i\epsilon_i\le(1-\rho^2)s/(12\rho)$ for every
$s\in(0,s_0]$;
\item[(iii)] $\sum_i\mu_i\epsilon_i\le(1-\rho^2)s_0/(3\rho)$.
\end{enumerate}
Then $c:=\sum_i\mu_ic_i$ satisfies $\rho g_c\in\Cert(f)$, and
$p^*(\rho g_c-\rho g)\le\rho\sum_i\mu_i\epsilon_i$.  Consequently, if for every
$\rho<1$ and $\varepsilon>0$ such families exist with
$\sum_i\mu_i\epsilon_i\le\varepsilon$, then $g\in\cl\Cert(f)$.
\end{theorem}

\begin{proof}
By linearity $g_c=\sum_i\mu_ig_{c_i}$, and by convexity of $p^*$ we
have $p^*(f+s\rho g_c)\le\sum_i\mu_ip^*(f+s\rho g_{c_i})$ for real $s$.  We
use two bounds for each $i$.  (A) If $\rho|s|\le r_i$ and $|s|\le s_0$,
Proposition~\ref{prop:expansion} and (i) give
$p^*(f+s\rho g_{c_i})\le1+\frac{\rho^2s^2}2(1+\eta_0)(1+\kappa_0)\le
1+\frac{s^2}4(1+\rho^2)$, since $\kappa(c_i)\rho|s|\le\kappa(c_i)
\min(r_i,s_0)$.  (B) Always, since $g\in\Cset(f)$,
\[
 p^*(f+s\rho g_{c_i})\le p^*(f+s\rho g)+\rho|s|\epsilon_i\le s(\rho s)
 +\rho|s|\epsilon_i.
\]

\emph{Case $|s|\le s_0$.}  Use (A) for $r_i\ge\rho|s|$ and (B) otherwise.
By (ii),
\[
 p^*(f+s\rho g_c)\le\max\Big(1+\frac{s^2}4(1+\rho^2),\,s(\rho s)\Big)
 +\frac{(1-\rho^2)s^2}{12}.
\]
Now $1+\frac{s^2}4(1+\rho^2)+\frac{(1-\rho^2)s^2}{12}=1+\frac{s^2}{6}
(2+\rho^2)=1+\frac{s^2}2-\frac{(1-\rho^2)s^2}6\le1+\frac{s^2}2-\frac{s^4}8
\le s(s)$ because $s^2\le1-\rho^2$; and $s(\rho s)+(1-\rho^2)s^2/12\le s(s)$
by Lemma~\ref{lem:slack}(a).

\emph{Case $|s|\ge s_0$.}  Use (B) for every $i$ and (iii):
$p^*(f+s\rho g_c)\le s(\rho s)+(1-\rho^2)s_0|s|/3\le s(s)$, because
$\min(s^2,|s|)\ge s_0|s|$.

Thus $\rho g_c\in\Cset(f)$.  By convexity of $H$,
$H(\rho c)\le\rho^2\max_iH(c_i)\le\rho^2(1+\eta_0)\le1$, so $\rho
g_c\in\Cert(f)$.  The distance estimate is the triangle inequality.
\end{proof}

The mechanism is that at each scale $s$ only the certificates that are too
fine for $s$ contribute an error, of order $s$ each, and these errors are
summable against the weights.

\begin{corollary}[Ladder form]\label{cor:ladder}
Let $g\in\Cset(f)$.  Suppose there are $c_0>0$, $\kappa_1,K<\infty$, $t_*>0$
and a function $\eta(t)\to0$ $(t\to0^+)$ such that for every $t\in(0,t_*]$
there is a finite certificate $c_t$ at $f$ with
\[
 H(c_t)\le1+\eta(t),\quad r(c_t)\ge c_0t,\quad \kappa(c_t)\le\kappa_1,\quad
 p^*(g-g_{c_t})\le Kt.
\]
Then $g\in\cl\Cert(f)$; hence $g$ is recovered along every sequence
$f_n\to f$ \textup{(}Corollary~\ref{cor:lsc}\textup{)}.  It suffices that
the hypotheses hold for the scales $t=t_*2^{1-i}$, $i\ge1$.
\end{corollary}

\begin{proof}
Fix $\rho<1$; choose $\eta_0,\kappa_0>0$ with
$\rho^2(1+\eta_0)(1+\kappa_0)\le(1+\rho^2)/2$, then
$s_0\in(0,\min(1,\sqrt{1-\rho^2})]$ with $\kappa_1s_0\le\kappa_0$, an integer
$n\ge24\rho^2K/(c_0(1-\rho^2))$, and $t_1\in(0,t_*]$ with $\eta(t)\le\eta_0$
for $t\le t_1$ and $2\rho Kt_1/n\le(1-\rho^2)s_0/3$; in the case of the last
sentence of the statement we take $t_1$ of the form $t_*2^{1-i_0}$, so that
all the scales used below are among the given ones.  Put $t_i:=t_12^{1-i}$,
$c_i:=c_{t_i}$ and $\mu_i:=1/n$ ($i\le n$).  Condition (i) holds.  For (ii),
$r_i<\rho s$ implies $c_0t_i<\rho s$, and the $t_i$ with this property sum
to less than $2\rho s/c_0$, so the left side of (ii) is at most
$(K/n)(2\rho s/c_0)\le(1-\rho^2)s/(12\rho)$.  For (iii),
$\sum_i\mu_iKt_i\le2Kt_1/n$.  Theorem~\ref{thm:averaging} gives $\rho
g_{c}\in\Cert(f)$ within $2\rho Kt_1/n$ of $\rho g$; let $t_1\to0$ and then
$\rho\to1$.
\end{proof}

\begin{corollary}[Decomposition form]\label{cor:decomposition}
Let $g\in\Cset(f)$.  Suppose that for every small $t>0$ there are a finite
certificate $c_t=(b_t,\omega_t)$ with $r(c_t)\ge c_0t$ and
$\kappa(c_t)\le\kappa_1$, and $e_t:=g-g_{c_t}$ with $p^*(e_t)\le Kt$, such
that for $\tau=t$ or for $\tau=-t$ the decomposition
\[
 f+\tau g=(a+\tau b_t+\tau e_t)+\sum_mR_m^*\big((1-d_{t,m}\tau)w_m+\tau
 \omega_{t,m}\big)
\]
has all component norms at most $s(\tau)+o(\tau^2)$.  Assume moreover that
$t^2/C_m\to0$ uniformly over the active blocks of $c_t$ \textup{(}this is
automatic if $I$ is finite\textup{)}.  Then $H(c_t)\le1+o(1)$, and
Corollary~\ref{cor:ladder} gives $g\in\cl\Cert(f)$.
\end{corollary}

\begin{proof}
\emph{Blocks.}  Since $\kappa(c_t)\le\kappa_1$, $|d_{t,m}|\le
\kappa_1C_m/(2M_m)$, so $Y=C_m+d_{t,m}\tau M_m=C_m(1+O(t))$ and
$1-d_{t,m}\tau>0$ for small $t$.  Put $X:=|\tau|\|h_\perp\|$.
Lemma~\ref{lem:block}(b) gives $\sqrt{Y^2+X^2}-Y\le\delta:=s(\tau)-1+o(\tau^2)$,
hence $X^2\le\delta(2Y+\delta)$ and
\[
 H_m(\omega_{t,m})=\frac{\|h_\perp\|^2}{C_m}\le\frac\delta{\tau^2}
 \Big(\frac{2Y}{C_m}+\frac\delta{C_m}\Big)=1+o(1)+O(\tau^2/C_m).
\]
\emph{Base.}  $B:=b_t+e_t$ satisfies $B(\zh)=0$, because $g(\xi)=0$ and
$g_{c_t}(\xi)=0$ (Lemma~\ref{lem:algebra}).  Also $\|U^*b_t\|\le
\kappa_1\nu/2$ and $\|e_t\|_1\le q^*(e_t)\le(1+\|L\|)p^*(e_t)\le(1+\|L\|)Kt$.
Lemma~\ref{lem:base}(a) gives $h(B)\le1+o(1)$, and since $\sqrt h$ is a
seminorm and $h(e_t)\le\|U\|^2\|e_t\|_1^2/\nu=O(t^2)$, also $h(b_t)\le
1+o(1)$.
\end{proof}

\begin{remark}\label{rem:G2}
For $I=\N$ the extra hypothesis in Corollary~\ref{cor:decomposition} cannot
simply be dropped: $C_m\le\|\Phi_m\|_2\le2^{-m}$, and a block with
$C_m\ll\tau^2$ can remain admissible at scale $\tau$ with
$\|h_\perp\|\approx\tau/2+C_m/\tau$, i.e.\ with $H_m\approx\tau^2/(4C_m)\gg1$.
Corollary~\ref{cor:ladder} itself has no such restriction.
\end{remark}

\subsection{Weighted directions}
For a block $m$, a function $\omega:\N\to\R$ vanishing on $P_m$ and
$\sigma>0$, put
\begin{gather*}
 E_m(\sigma;\omega):=\{k\in Q_m:\sigma|\omega(k)|>\gap_m(k)/2\},\\
 T_m(\sigma;\omega):=\sum_{k\in E_m(\sigma;\omega)}\lambda_{k,m}|\omega(k)|.
\end{gather*}

\begin{lemma}[Box tails]\label{lem:boxtails}
$T_m(\sigma;\omega)=o(\sigma)$ as $\sigma\to0$ in each of the following cases:
\textup{(a)} $\sum_k\lambda_{k,m}\omega(k)^2/\gap_m(k)<\infty$;
\textup{(b)} $\omega$ is bounded and $|\omega(k)|\le K\sqrt{\gap_m(k)}$
whenever $\gap_m(k)\le\gamma_0$, for some $K,\gamma_0>0$;
\textup{(c)} $\gap_m\ge M_m/2$ on $\supp\omega$ and
$\sum_k\lambda_{k,m}\omega(k)^2<\infty$ \textup{(}the weighted directions
of the weighted-directions theorem of \cite{PrA}\textup{)}.
\end{lemma}

\begin{proof}
(a) On $E_m(\sigma)$, $\lambda_k|\omega(k)|<2\sigma\lambda_k\omega(k)^2/
\gap(k)$, and the sets $E_m(\sigma)$ decrease to $\emptyset$ as
$\sigma\downarrow0$; use dominated convergence.  (b) If
$\sigma\|\omega\|_\infty<\gamma_0/2$, then $k\in E_m(\sigma)$ forces
$\gap(k)\le\gamma_0$, hence $\gap(k)<2\sigma K\sqrt{\gap(k)}$, i.e.\
$\gap(k)<4K^2\sigma^2$ and $|\omega(k)|<2K^2\sigma$; so $T_m(\sigma)\le
2K^2\sigma\sum_{\gap(k)<4K^2\sigma^2}\lambda_k=o(\sigma)$.  (c) is a case
of (a).
\end{proof}

\begin{theorem}[Weighted directions]\label{thm:weighted}
Let $b\in\ell_1$ with $\supp b\subseteq F$, $b(\zh)=0$ and
$\sum_{j\in F,\,|b_j|>|a_j|/(2t)}|b_j|=O(t)$ \textup{(}e.g.\
$\sum_{j\in F}b_j^2/|a_j|<\infty$\textup{)}.  For finitely many blocks $m$
let $\omega_m:\N\to\R$ vanish on $P_m$, with
$\sum_k\lambda_{k,m}|\omega_m(k)|<\infty$ and $T_m(\sigma;\omega_m)=O(\sigma)$.
Put $d_m:=\ip{D_mw_m}{D_m\omega_m}/C_m$,
$g:=b+\sum_mR_m^*(\omega_m-d_mw_m)$ and $H:=\max(h(b),\max_mH_m(\omega_m))$.
If $g\in\Cset(f)$ and $H\le1$, then $g\in\cl\Cert(f)$; in particular $g$
is recovered along every sequence $f_n\to f$.
\end{theorem}

\begin{proof}
Since $\Phi_m\le\lambda_{\cdot,m}$, $D_m\omega_m\in\ell_1\subseteq\ell_2$.
For small $t>0$ choose $J_t\to\infty$ so large that the tails of $b$ and of
$\lambda_{\cdot,m}\omega_m$ beyond $J_t$ have $\ell_1$-norm at most $t$, and
put
\begin{gather*}
 G_t:=F\cap[1,J_t]\cap\{|b_j|\le|a_j|/(2t)\},\qquad
 b_t:=b1_{G_t}-\varkappa_ta,\qquad \varkappa_t:=(b1_{G_t})(\zh),\\
 \omega_{m,t}:=\omega_m1_{\{k\le J_t\}\setminus E_m(t;\omega_m)}.
\end{gather*}
Then $c_t:=(b_t,(\omega_{m,t}))$ is a finite certificate.  Since
$|\varkappa_t|=|(b1_{G_t^c})(\zh)|\le q^*(b1_{F\setminus
G_t})\le(1+\|U\|)\|b1_{F\setminus G_t}\|_1=O(t)$, we get
$\|b_t/a\|_\infty\le1/(2t)+O(t)$, so the base part of $r(c_t)$ is at least
$t$ for small $t$; the block coordinates kept satisfy
$t|\omega_{m,t}(k)|\le\gap_m(k)/2$, and the remaining terms of the radius
are bounded below (the $d_{m,t}$ and $\|U^*b_t\|$ are bounded).  Hence
$r(c_t)\ge t$ and $\kappa(c_t)$ is bounded.  The remainder is
\[
 g-g_{c_t}=(b-b_t)+\sum_mR_m^*\big((\omega_m-\omega_{m,t})-(d_m-d_{m,t})w_m
 \big),
\]
with $\|b-b_t\|_1\le\|b1_{F\setminus G_t}\|_1+|\varkappa_t|=O(t)$,
$\sum_k\lambda_{k,m}|\omega_m-\omega_{m,t}|(k)\le T_m(t;\omega_m)+t=O(t)$ and
$|d_m-d_{m,t}|\le\|D_m(\omega_m-\omega_{m,t})\|_2=O(t)$.  So
$p^*(g-g_{c_t})\le Kt$.  Finally $b_t\to b$ in $\ell_1$ and
$D_m\omega_{m,t}\to D_m\omega_m$ in $\ell_2$, so $H(c_t)\to H\le1$.
Corollary~\ref{cor:ladder} applies.
\end{proof}

By Lemma~\ref{lem:boxtails}(c), Theorem~\ref{thm:weighted} contains
the weighted-directions theorem of \cite{PrA} (and removes its dependence on \cite{Check});
near-peak supports are allowed as long as the box tails are $O(\sigma)$.
In particular the non-attaining examples of \cite[Section~3]{PrA}, which
admit no bounded first-order lift, are recovered along every sequence.

\subsection{Locally split mates}

\begin{definition}\label{def:locsplit}
A mate $g\in\Cset(f)$ is \emph{locally split} if there are $\tau>0$,
$b\in X^*$ and $v\in V^*$ with $g=b+L^*v$ and, for $|t|\le\tau$,
\[
 q^*(a+tb)\le s(t)\quad\text{and}\quad\|w+tv\|_{V^*}\le s(t).
\]
\end{definition}

Equivalently, the linear path $f+tg=(a+tb)+L^*(w+tv)$ is admissible at
level $s(t)$ near $t=0$.  Every finite certificate direction $g_c$ with
$H(c)<1$ is locally split: by Lemma~\ref{lem:base}(b),
Lemma~\ref{lem:block}(d) and the argument of Lemma~\ref{lem:slack}(b),
each component is at most $s(t)$ for $|t|\le r_*(c,1-H(c))$.  (For
$H(c)=1$ this can fail: the expansion only gives
$1+\frac{t^2}2(1+\kappa(c)|t|)$, which exceeds $s(t)$ when $\kappa(c)>0$.)
The converse fails, since $b$ and $v$ may have infinite supports.

\begin{lemma}[Base truncation with exact first-order correction]\label{lem:basetrunc}
Let $b\in\ell_1$ satisfy $q^*(a+tb)\le s(t)$ for $|t|\le\tau$.  Then
$\supp b\subseteq F$ and $b(\zh)=0$.  Put
$\tau_a:=\frac14\min\big(\tau,1,\nu/(2\|U^*b\|+1)\big)$.  For every
$\delta>0$ there is $N$ such that $b'':=P_Nb+\beta_Na$, with $P_N$ the
coordinate projection onto $[1,N]$ and $\beta_N:=((I-P_N)b)(\zh)$,
satisfies $b''(\zh)=0$, $\|b''-b\|_1<\delta$, and
$q^*(a+tb'')\le s(t)+\delta t^2$ for $|t|\le\tau_a$.
\end{lemma}

\begin{proof}
\emph{First order.}  The function $\varphi(t)=q^*(a+tb)$ is convex with
$\varphi(0)=1$ and $\varphi(t)\le1+t^2/2$ on $[-\tau,\tau]$, so both one-sided
derivatives at $0$ vanish.  The right derivative is
$\sum_{j\in F}\sgn(a_j)b_j+\sum_{j\notin F}|b_j|+\ip{U^*b}{e}$ and the left
derivative is $\sum_{j\in F}\sgn(a_j)b_j-\sum_{j\notin
F}|b_j|+\ip{U^*b}{e}$; hence $\sum_{j\notin F}|b_j|=0$ and $b(\zh)=0$.

\emph{Truncation.}  Let $r:=(I-P_N)b$ and $\varepsilon_N:=\|U^*r\|$; note
that $\beta_N=r(\zh)=\sum_{j>N}z_jr_j+\ip{U^*r}{e}$ tends to $0$.  For $|s|\le2\tau_a$ and
$j>N$ in $F$, $|a_j|\le|a_j+sr_j|-sz_jr_j$; hence
$\|a+sP_Nb\|_1\le\|a+sb\|_1-s\sum_{j>N}z_jr_j$.  For the Hilbert part we use
the elementary inequality
\begin{equation}\label{eq:hilbertineq}
 \|y-x\|\le\|y\|-\ip{y}{x}/\|y\|+\|x\|^2/(2\|y\|)\qquad(y\ne0),
\end{equation}
which follows from
$(\ip{y}{x}/\|y\|^2-\|x\|^2/(2\|y\|^2))^2\ge0$.  With
$y_s:=U^*(a+sb)$ (so $\|y_s\|\ge\frac34\nu$ and
$\|y_s/\|y_s\|-e\|\le2|s|\|U^*b\|/\nu$) and $x=sU^*r$,
$\|U^*(a+sP_Nb)\|\le\|y_s\|-s\ip{e}{U^*r}+C_2s^2\varepsilon_N$, with $C_2$
depending only on $\nu$, $\|U^*b\|$.  Adding,
\[
 q^*(a+sP_Nb)\le q^*(a+sb)-s\beta_N+C_2s^2\varepsilon_N.
\]
Now $a+tb''=(1+t\beta_N)(a+sP_Nb)$ with $s=t/(1+t\beta_N)$; for $N$ large
($|\beta_N|\le1$) and $|t|\le\tau_a$ we have $|s|\le2\tau_a$, and
\[
 q^*(a+tb'')\le(1+t\beta_N)s(s)-t\beta_N+C_3t^2\varepsilon_N
 =1+F_t(\beta_N)+C_3t^2\varepsilon_N,
\]
where $F_t(\beta):=\sqrt{(1+t\beta)^2+t^2}-(1+t\beta)$.  Since $F_t(0)=
s(t)-1$ and $|\partial_\beta F_t|=|t|\,|(1+t\beta)/\sqrt{(1+t\beta)^2+t^2}-1|
\le2|t|^3$ for $|t\beta|\le\frac12$, we get $q^*(a+tb'')\le s(t)+
2|\beta_N||t|^3+C_3t^2\varepsilon_N\le s(t)+\delta t^2$ for $N$ large.
Finally $\|b''-b\|_1\le\|r\|_1+|\beta_N|\|a\|_1\to0$ and
$b''(\zh)=b(\zh)-r(\zh)+\beta_N=0$.
\end{proof}

\begin{lemma}[Block truncation]\label{lem:blocktrunc}
Fix a block \textup{(}drop $m$\textup{)} and let $v\in\ell_\infty$ satisfy
$N(w+tv)\le s(t)$ for $|t|\le\tau$.  Put $c:=\ip{Dw}{Dv}/C$, $d:=c/M$ and
$\omega:=v+dw$.  Then $\omega=0$ on $P$, $|\sgn(w(k))v(k)+c|\le
\sqrt{2\gap(k)}$ whenever $w(k)\neq0$ and $\gap(k)\le\tau^2/2$, and
$d=\ip{Dw}{D\omega}/C$.  Moreover there is $\tau_3>0$ such that for every
$\delta>0$ there is a finitely supported $\omega''$, vanishing on $P$, with
$\min_{\supp\omega''}\gap>0$ and $\ip{Dw}{D\omega''}=\ip{Dw}{D\omega}$, such
that $v'':=\omega''-dw$ satisfies $\|R^*(v''-v)\|_1<\delta$ and
$N(w+tv'')\le N(w+tv)+\delta t^2$ for $|t|\le\tau_3$.
\end{lemma}

\begin{proof}
\emph{Structure.}  $\chi(t):=\|D(w+tv)\|$ is convex and differentiable at
$0$ with $\chi'(0)=c$, so $\chi(t)\ge C+tc$, and $\|w+tv\|_\infty\le s(t)-
\chi(t)\le M-tc+t^2/2$ for $|t|\le\tau$.  For $k$ with $w(k)\neq0$ and
$\varsigma_k:=\sgn w(k)$ this gives
$|w(k)|+t\varsigma_kv(k)\le M-tc+t^2/2$, i.e.\
$t(\varsigma_kv(k)+c)\le\gap(k)+t^2/2$ for $|t|\le\tau$.  For a peak
($\gap=0$) both signs of $t$ give $\varsigma_kv(k)=-c$; for $\gap(k)\le
\tau^2/2$ take $t=\pm\sqrt{2\gap(k)}$.  Then $\omega(k)=\varsigma_k(\varsigma_k
v(k)+c)=0$ on $P$, and off $P$ with $w(k)\ne0$,
$\omega(k)=\varsigma_k\big((\varsigma_kv(k)+c)-c\gap(k)/M\big)$, which tends to
$0$ along every sequence with $\gap(k)\to0$; $\omega$ is bounded.
Finally $\ip{Dw}{D\omega}=cC+dC^2=cC(M+C)/M=dC$.

\emph{Truncation.}  If $w(k)=0$ for every $k\in Q$, then
$\ip{Dw}{D\omega}=0$, $d=c=0$, and we put $k_0:=$ none, $\varkappa:=0$.
Otherwise fix $k_0\in Q$ with $w(k_0)\neq0$, $\gamma_0:=\gap(k_0)$.  For
$\gamma\in(0,\gamma_0)$ and $K\ge k_0$ let $S:=\{k\le K:\gap(k)\ge\gamma\}$
(finite, contains $k_0$) and
\[
 \omega'':=\omega1_S+\varkappa e_{k_0},\qquad
 \varkappa:=\ip{Dw}{D(\omega-\omega1_S)}/(\Phi(k_0)^2w(k_0)),
\]
so that $\ip{Dw}{D\omega''}=\ip{Dw}{D\omega}$.  As $\gamma\to0$ and
$K\to\infty$, $\omega-\omega1_S\to0$ pointwise and boundedly ($\omega=0$ on
$P$), hence $\Delta:=D(\omega''-\omega)\to0$ in $\ell_2$, $\varkappa\to0$ and
$\|R^*(\omega''-\omega)\|_1\le\sum_k\lambda_k|\omega''(k)-\omega(k)|\to0$.

\emph{Sup part.}  Let $W:=w+tv=(1-dt)w+t\omega$ and
$W'':=w+tv''=(1-dt)w+t\omega''$.  They agree on $S\setminus\{k_0\}$.  For
$k\notin S\cup\{k_0\}$, $|W''(k)|=|1-dt||w(k)|\le|1-dt|M=|W(k_*)|$ for any
peak $k_*$ ($\omega(k_*)=0$), so $|W''(k)|\le\|W\|_\infty$.  At $k_0$,
$|W''(k_0)|\le|1-dt|(M-\gamma_0)+|t|(\|\omega\|_\infty+1)\le|1-dt|M$ for
$|t|\le t_0:=\gamma_0/(4(|d|M+\|\omega\|_\infty+2))$ (once $|\varkappa|\le1$).
Hence $\|W''\|_\infty\le\|W\|_\infty$ for $|t|\le t_0$.

\emph{Hilbert part.}  Put $A_t:=DW$.  Since $\ip{Dw}{\Delta}=0$,
$\ip{A_t}{\Delta}=t\ip{D\omega}{\Delta}$ and
\begin{gather*}
 \|A_t+t\Delta\|^2=\|A_t\|^2+2t^2\ip{D\omega}{\Delta}+t^2\|\Delta\|^2,
 \quad\text{so}\\
 \|A_t+t\Delta\|\le\|A_t\|+\frac{2t^2(\|D\omega\|\|\Delta\|+\|\Delta\|^2)}{C}
\end{gather*}
for $|t|\le t_1:=C/(2\|D(\omega-dw)\|+2)$ (then $\|A_t\|\ge C/2$).  Adding,
$N(W'')\le N(W)+o(1)t^2$ on $|t|\le\tau_3:=\min(\tau,t_0,t_1)$, where
$o(1)\to0$ as $\gamma\to0$, $K\to\infty$.
\end{proof}

\begin{theorem}[Locally split mates]\label{thm:locsplit}
Let the block set $I$ be finite.  Every locally split mate
$g\in\Cset(f)$ lies in $\cl\Cert(f)$; in particular $(f,g)\in
\cl\NA((c_0,p),\ell_2^2)$ and $g$ is recovered along every sequence
$f_n\to f$ in $S_{p^*}$.
\end{theorem}

\begin{proof}
Let $g=b+L^*v$ be as in Definition~\ref{def:locsplit}.  Fix $\rho\in(0,1)$.
Let $\tau_a$ be the range of Lemma~\ref{lem:basetrunc} for $b$ and
$\tau_{3,m}$ the range of Lemma~\ref{lem:blocktrunc} for $v_m$, and put
$\tau_5:=\min(\tau_a,\min_{m\in I}\tau_{3,m})>0$; since $I$ is finite and
these ranges do not depend on the truncation accuracy, $\tau_5$ does not
depend on $\delta$ below.  Fix $\delta>0$ and apply
Lemma~\ref{lem:basetrunc} to $b$ and Lemma~\ref{lem:blocktrunc} to each
$v_m$, $m\in I$, with accuracy $\delta$: we obtain a finite certificate
$c''=(b'',(\omega''_m)_{m\in I})$ \textup{(}its $d$-coefficients are the
$d_m$ of Lemma~\ref{lem:blocktrunc}, because
$\ip{D_mw_m}{D_m\omega''_m}=\ip{D_mw_m}{D_m\omega_m}$ with
$\omega_m:=v_m+d_mw_m$ as in that lemma\textup{)} with
\[
 p^*(g_{c''}-g)\le q^*(g_{c''}-g)\le(1+\|U\|)\|g_{c''}-g\|_1
 \le\delta':=(1+\|U\|)(|I|+1)\delta,
\]
and, for $|t|\le\tau_5$,
\[
 q^*(a+tb'')\le s(t)+\delta t^2,\qquad
 N_m(w_m+t(\omega''_m-d_mw_m))\le s(t)+\delta t^2 .
\]
\emph{Coefficient.}  Lemma~\ref{lem:base}(a) gives
$\frac{t^2}2h(b'')/(1+|t|\|U^*b''\|/\nu)\le s(t)-1+\delta t^2$; dividing
by $t^2$ and letting $t\to0$, $h(b'')\le1+2\delta$.  Similarly
Lemma~\ref{lem:block}(b) gives $H_m(\omega''_m)\le1+2\delta$.  So
$H(c'')\le1+2\delta$.

\emph{Mate.}  For $|t|\le\tau_5/\rho$ we have
\[
 p^*(f+t\rho g_{c''})\le s(\rho t)+\delta\rho^2t^2\le s(t)
 \quad\text{as soon as}\quad
 \delta\rho^2\max(1,\tau_5/\rho)\le(1-\rho^2)/3
\]
(Lemma~\ref{lem:slack}(a)).  For $|t|\ge\tau_5/\rho$ we have
\[
 p^*(f+t\rho g_{c''})\le s(\rho t)+\rho|t|\delta'\le s(t)
 \quad\text{as soon as}\quad
 \rho\delta'\le(1-\rho^2)\min(1,\tau_5/\rho)/3.
\]
Since $\tau_5$ does not depend on $\delta$ and $\delta'\to0$ with $\delta$,
both conditions hold for $\delta$ small; then $\rho g_{c''}\in\Cset(f)$ and $H(\rho c'')\le
\rho^2(1+2\delta)\le1$, i.e.\ $\rho g_{c''}\in\Cert(f)$, at distance at most
$\rho\delta'$ from $\rho g$.  Let $\delta\to0$ and $\rho\to1$.  The last
assertions follow from Corollary~\ref{cor:lsc}.
\end{proof}

\begin{corollary}[Split-contractive operators]\label{cor:splitcontractive}
Let the block set $I$ be finite.  Let $S\in\mathcal L((c_0,q),\ell_2^2)$
and $\Phi\in\mathcal L(V,\ell_2^2)$ with $\|S\|_q\le1$, $\|\Phi\|\le1$, and
suppose $A_0:=S+\Phi\circ L$ has $\|A_0\|_p=1$.  Then
$A_0\in\cl\NA((c_0,p),\ell_2^2)$.
\end{corollary}

\begin{proof}
After a rotation $A_0=(f,g)$ with $p^*(f)=1$, $f=s_1+L^*\varphi_1$,
$g=s_2+L^*\varphi_2$, where $(s_1,s_2)$ is $q$-contractive and
$(\varphi_1,\varphi_2)$ is $V$-contractive.  Then $q^*(s_1),\|\varphi_1\|\le
1$, so by Proposition~\ref{prop:forced}, $s_1=a$ and $\varphi_1=w$; the
contractivity of the components gives $q^*(a+ts_2)\le s(t)$ and
$\|w+t\varphi_2\|\le s(t)$ for all $t$.  So $g$ is locally split, and
Theorem~\ref{thm:locsplit} applies.
\end{proof}

Since the norm-attaining operators of the split form $S+\Phi\circ L$ are
not dense \cite[Section~3]{PrA}, the corollary shows that split operators
are recovered by non-split norm-attaining ones.

\subsection{Carriers with non-summable gaps}

\begin{theorem}[Carrier mates]\label{thm:carriers}
Let $g\in\Cset(f)$, let $c^0=(b^0,\omega^0)$ be a finite certificate, fix a
block $m$ and $c\in\R\setminus\{0\}$, and let $k_1,k_2,\dots$ be distinct
strict non-peaks of $w_m$ outside $\supp\omega^0_m$, with
$\lambda_i:=\lambda_{k_i,m}$ and $\gamma_i:=\gap_m(k_i)$, such that
\begin{enumerate}
\item[(a)] $\gamma_{i+1}\lambda_{i+1}\le\gamma_i\lambda_i/2$ for all $i$;
\item[(b)] $p^*(g-g_{c^0}-c\,u_{k_i,m})\le K\lambda_i$ for all $i$;
\item[(c)] $H_*:=\max\big(h(b^0),\ \max_{m'\ne m}H_{m'}(\omega^0_{m'}),\
H_m(\omega^0_m)+c^2/(m^2C_m)\big)\le1$;
\item[(d)] $\sum_i\gamma_i=\infty$.
\end{enumerate}
Then $g\in\cl\Cert(f)$.  In particular this holds when $\inf_i\gamma_i>0$,
$\|u_{k_i,m}-v\|\le K'\lambda_i$ for some $v$ with $g=g_{c^0}+cv$, and
\textup{(c)} holds \textup{(}after passing to a subsequence for which
\textup{(a)} holds\textup{)}: cross mates with rates through carriers whose
gaps are bounded below are recovered along every sequence, provided
\textup{(c)} holds.
\end{theorem}

\begin{proof}
Let $c_i:=c^0+(0,(c/\lambda_i)e_{k_i}\text{ in block }m)$.  Its
$d$-coefficient in block $m$ is $d_{i}=d^0_m+c\Phi_m(k_i)w_m(k_i)/(mC_m)
\to d^0_m$, and $\|D_m\omega_{i}\|^2=\|D_m\omega^0_m\|^2+c^2/m^2$ (because
$\Phi_m(k_i)/\lambda_i=1/m$), so $\limsup_iH(c_i)\le H_*\le1$.  The
coordinatewise radius is $r_i=\min(r_*,\gamma_i\lambda_i/(2|c|))$ with
$r_*>0$ bounded below for large $i$, $\kappa(c_i)$ is bounded, and
\[
 \epsilon_i:=p^*(g-g_{c_i})\le p^*(g-g_{c^0}-cu_{k_i,m})+|d_i-d^0_m|p^*(R_m^*w_m)
 \le K'\lambda_i .
\]
Fix $\rho<1$, choose $\eta_0,\kappa_0,s_0$ as in Theorem~\ref{thm:averaging}
with moreover $s_0\sup_i\kappa(c_i)\le\kappa_0$.  For a stretch
$i_0\le i\le i_1$ put $\Gamma:=\sum_{i_0}^{i_1}\gamma_i$ and
$\mu_i:=\gamma_i/\Gamma$.  If $s\le r_*/\rho$, then by (a)
\[
 \sum_{r_i<\rho s}\mu_i\epsilon_i\le\frac{K'}{\Gamma}\sum_{\gamma_i\lambda_i<2|c|\rho
 s}\gamma_i\lambda_i\le\frac{4K'|c|\rho s}{\Gamma}\le\frac{(1-\rho^2)s}
 {12\rho}
\]
once $\Gamma\ge48K'|c|\rho^2/(1-\rho^2)$, which is possible with $i_0$
arbitrarily large by (d).  If $r_*/\rho<s\le s_0$, then
$\sum_i\mu_i\epsilon_i\le K'\sup_{i\ge i_0}\lambda_i\le(1-\rho^2)r_*/(12\rho^2)$
for $i_0$ large (note $\lambda_i\to0$), which gives (ii) as well; (iii) holds similarly, and (i) holds for
$i_0$ large because $\limsup H(c_i)\le1$ and we may take $\eta_0>0$.
Theorem~\ref{thm:averaging} gives $\rho g_c\in\Cert(f)$ within
$\rho K'\sup_{i\ge i_0}\lambda_i$ of $\rho g$.  For the last assertion,
$\epsilon_i\le|c|K'\lambda_i+O(\lambda_i)$, and a subsequence with
$\lambda_{i+1}\le\lambda_i\inf\gamma/2$ satisfies (a).
\end{proof}

\begin{remark}\label{rem:summablegaps}
If $\sum_i\gamma_i<\infty$ the argument breaks down: condition (ii) at
$s\approx\gamma_i\lambda_i$ forces $\mu_i\lesssim\gamma_i$.  Whether
switching mates carried by such ``super-near-threshold'' coordinates, or by
weak peaks (small margins), lie in $\cl\Cert(f)$ is open; their limit
directions lie in the closed span of the non-peak vectors, so the linear
obstruction of Section~\ref{sec:resonance} does not apply to them.
\end{remark}

\section{Second-order rebalancing and tame supports}\label{sec:second}

In this section the block set $I$ is finite.  Let $f\in S_{p^*}$ have a
normer $\eta\in S_{p^{**}}$; either $\eta=\xi$ is the bidual normer of an
arbitrary $f$, or $\eta=x'\in S_p$ when $f=f'$ attains its norm.  We use the
forced data of Section~\ref{sec:setting} at $\eta$, writing
$c_\eta:=q^{**}(\eta)$ (so $c_\eta=q_0$ when $\eta=\xi$),
$\zh:=\eta/c_\eta=z+Ue$, $\zeta_m:=R_m^{**}\eta$ and
$\sigma_m:=|\zeta_m|_m$, so that $c_\eta+\sum_m\sigma_m=1$.

\subsection{Slices, excess and tame points}

\begin{lemma}[Bidual slice]\label{lem:slice}
Let $g\in\Cset(f)$.  Then $g(\eta)=0$, and for every $\chi\in X^{**}$ with
$\chi(f)=0$,
\[
 \chi(g)^2\le\varphi(\chi)\big(2+\varphi(\chi)\big),\qquad
 \varphi(\chi):=p^{**}(\eta+\chi)-1 .
\]
In particular, if $\varphi(s\chi)=o(s^2)$ as $s\to0^+$, then $\chi(g)=0$.
\end{lemma}

\begin{proof}
For $\theta\in X^{**}$ and $t\in\R$,
$\theta(f)+t\theta(g)=\theta(f+tg)\le p^{**}(\theta)s(t)$; the supremum of the
left side divided by $s(t)$ over $t$ is $\sqrt{\theta(f)^2+\theta(g)^2}$, so
$\theta(f)^2+\theta(g)^2\le p^{**}(\theta)^2$.  With $\theta=\eta$ we get
$g(\eta)=0$, and with $\theta=\eta+\chi$ we get
$1+\chi(g)^2\le(1+\varphi(\chi))^2$.  For the last claim apply this to
$s\chi$ and divide by $s^2$.
\end{proof}

\begin{lemma}[Excess splitting]\label{lem:excess}
For $\chi\in X^{**}$ put
$\Delta(\eta+\chi):=p^{**}(\eta+\chi)-f(\eta+\chi)\ge0$.  Then
\begin{gather*}
 \Delta(\eta+\chi)=E_q(\chi)+\sum_mE_m(R_m^{**}\chi),\\
 E_q(\chi):=q^{**}(\eta+\chi)-c_\eta-a(\chi),\qquad
 E_m(\delta):=|\zeta_m+\delta|_m-\sigma_m-w_m(\delta),
\end{gather*}
and all these terms are nonnegative.  If $f(\chi)=0$, then
$\Delta(\eta+\chi)=\varphi(\chi)$.
\end{lemma}

\begin{proof}
$p^{**}(\eta)=1=f(\eta)=c_\eta+\sum_m\sigma_m$, $f=a+\sum_mR_m^*w_m$ and
$p^{**}=q^{**}+\sum_m|R_m^{**}\cdot|_m$; nonnegativity holds because $a$ and
$w_m$ have norm one.
\end{proof}

\begin{lemma}[Flatness on free coordinates]\label{lem:flat}
If $\chi\in\ell_\infty$ is supported in $J$ and $|\chi_j|\le
c_\eta(1-|z_j|)$ for all $j$, then $E_q(\chi)=0$ and $a(\chi)=0$.
\end{lemma}

\begin{proof}
$\eta+\chi=c_\eta((z+\chi/c_\eta)+Ue)$ with $\|z+\chi/c_\eta\|_\infty\le1$,
so $q^{**}(\eta+\chi)\le c_\eta$; and $q^{**}(\eta+\chi)\ge a(\eta+\chi)=
c_\eta$ because $J\cap F=\emptyset$.
\end{proof}

\begin{lemma}[Overshoot bound]\label{lem:overshoot}
Fix a block and drop $m$.  Let $\delta=\mu\zeta+\delta'$ with $\mu>-1$,
$\delta'\in\ell_1$ supported in $P$ and $\sum_{k\in P}s_k\delta'_k=0$, where
$s_k:=\sgn w(k)$.  Then
\[
 E(\delta)\le2\sum_{k\in P}\big(-s_k\delta'_k-(1+\mu)|\zeta||\alpha(k)|
 \big)_+\le2\sum_{k\in P}\big(|\delta'_k|-(1+\mu)\lambda_k\mu_k\big)_+ .
\]
\end{lemma}

\begin{proof}
By Lemma~\ref{lem:threshold}, $\zeta+\delta=[(1+\mu)|\zeta|\alpha+\delta']
+(1+\mu)|\zeta|D(Dw/C)$.  Using $|x+y|=|x|+sy+2(-sy-|x|)_+$ for $s\in\{\pm1\}$
with $sx=|x|$, the $\ell_1$-norm of the bracket is $(1+\mu)|\zeta|+
\sum_Ps_k\delta'_k+2\,\mathrm{ov}=(1+\mu)|\zeta|+2\,\mathrm{ov}$, where
$\mathrm{ov}$ is the sum in the statement; the second term lies in
$(1+\mu)|\zeta|D(B_{\ell_2})$.  Hence $|\zeta+\delta|\le(1+\mu)|\zeta|
+2\,\mathrm{ov}$, while $w(\delta)=\mu|\zeta|+M\sum_Ps_k\delta'_k=\mu|\zeta|$.
The second inequality is \eqref{eq:margin}.
\end{proof}

Note that $|(R_m\chi)(k)|=\lambda_{k,m}|u_{k,m}(\chi)|\le\lambda_{k,m}q(\chi)$
for $\chi\in X$.

\begin{definition}[Tame points, certificate span]\label{def:tame}
Let $P^\circ_m:=\{k\in P_m:\mu_{k,m}>0\}$ and $\mathrm{Dg}_m:=P_m\setminus P_m^\circ$
(degenerate peaks), $\bar Q_m:=Q_m\cup \mathrm{Dg}_m$.  The normer $\eta$ satisfies
\emph{margin sparsity} \textup{(MS)} if
$\sum_{k\in P_m^\circ,\ \mu_{k,m}<s}\Phi_m(k)=o(s)$ as $s\to0^+$ for every
$m$.  It is \emph{C-tame} if $a\in c_{00}$, $K$ is finite, every
$\bar Q_m$ is finite, and \textup{(MS)} holds.  (At $\eta=x'\in c_0$ the
first two conditions are automatic.)  The \emph{certificate span} at $f$ is
\begin{gather*}
 S(f):=\Big\{b+\sum_mR_m^*(\omega_m-d_m(\omega_m)w_m):\ b\in\ell_1(F),\
 b(\eta)=0,\ \omega_m\in c_{00}(Q_m)\Big\},\\
 d_m(\omega):=\frac{\ip{D_mw_m}{D_m\omega}}{C_m},
\end{gather*}
and its elements are called \emph{balanced finite certificates} (the base
part need not satisfy $\|b/a\|_\infty<\infty$ when $F$ is infinite).
Equivalently
$S(f)=(\ell_1(F)\cap\eta^\perp)+\spn\{y_{k,m}:k\in Q_m,m\in I\}$ with
$y_{k,m}:=\lambda_{k,m}u_{k,m}-(\Phi_m(k)^2w_m(k)/C_m)R_m^*w_m$.  For
$g\in S(f)$ with data $(b,\omega)$ put $H_b:=h(b)$, $H_m:=H_m(\omega_m)$ and
\[
 \Gamma_w(g):=c_\eta H_b+\sum_m\sigma_mH_m,\qquad
 \Gamma_{\max}(g):=\max(H_b,\max_mH_m).
\]
\end{definition}

When $F$ and all $Q_m$ are finite the data $(b,\omega)$ of $g\in S(f)$ are
unique (Lemma~\ref{lem:rigidity}), so $\Gamma_w$ is well defined; in general
we regard $\Gamma_w$ as a function of the data.  Since $c_\eta+\sum_m
\sigma_m=1$, $\Gamma_w\le\Gamma_{\max}$.

\begin{theorem}[Fibres at C-tame points]\label{thm:tamefibre}
If the normer $\eta$ of $f$ is C-tame, then $\Cset(f)\subseteq S(f)$; that
is, every mate is a balanced finite certificate
$g=b+\sum_mR_m^*(\omega_m-d_mw_m)$ with $\supp b\subseteq F$, $b(\eta)=0$
and $\omega_m$ supported in the finite set $Q_m$.  In particular
$\Cset(f)$ is finite dimensional.
\end{theorem}

\begin{proof}
$J$ is cofinite.  By Lemma~\ref{lem:rigidity}(b) with $P'_m:=P^\circ_m$
(cofinite, hence infinite) and $\pi_m:=\sum_{k\in P_m}s_k\lambda_{k,m}
u_{k,m}$ (so that $\pi_m(\chi)=\sum_{P_m}s_k(R_m\chi)_k$; the degenerate
peaks are among the finitely many $u_{k,m}$, $k\in\bar Q_m$), the
restrictions to $c_{00}(J)$ of the finite family
\[
 \Psi:=\{u_{k,m}:k\in\bar Q_m,m\in I\}\cup\{\pi_m:m\in I\}
\]
are linearly independent.  Note $R_m^*w_m=M_m\pi_m+\sum_{k\in
Q_m}w_m(k)\lambda_{k,m}u_{k,m}$.  We use repeatedly: if $\chi\in X$ and
$R_m\chi=\mu\zeta_m+\delta'$ where $\delta'$ vanishes on $\bar Q_m$ and
has zero peak sum, then by Lemma~\ref{lem:overshoot} and \textup{(MS)},
$E_m(t\mu\zeta_m+t\delta')\le2m\sum_{k\in P^\circ_m}\Phi_m(k)\big(|t|B-
\tfrac12\mu_{k,m}\big)_+=o(t^2)$ for $|t\mu|\le\frac12$, where $B$ bounds
$|\delta'_k|/\lambda_{k,m}$.

\emph{Step 1: $\Cset(f)\subseteq W:=\spn\{e_j^*:j\in F\cup K\}+\spn\Psi$.}
Let $\nu\in c_{00}(J)$ with $\psi(\nu)=0$ for all $\psi\in\Psi$.  Then
$a(\nu)=0$, $R_m\nu$ vanishes on $\bar Q_m$, and its peak sum is
$\pi_m(\nu)=0$; also $f(\nu)=a(\nu)+\sum_mw_m(R_m\nu)=0$.  For small $|t|$,
$E_q(t\nu)=0$ (Lemma~\ref{lem:flat}: $\nu$ is finitely supported in $J$) and
$E_m(tR_m\nu)=o(t^2)$, so $\varphi(t\nu)=o(t^2)$ and $g(\nu)=0$ for every
$g\in\Cset(f)$ (Lemma~\ref{lem:slice}).  By linear algebra, the restriction
of $g$ to $c_{00}(J)$ is a combination of the restrictions of $\Psi$; hence
$g$ minus that combination is supported in $F\cup K$.

\emph{Step 2: unique form.}  Since $\pi_m$ and the $u_{k,m}$ are images
under $R_m^*$ of $s1_{P_m}$ and $\lambda_{k,m}^{-1}e_k$, every $g\in W$ can
be written $g=b+\sum_mR_m^*(\omega_m+c_mw_m)$ with $\supp b\subseteq F\cup
K$ and $\omega_m\in c_{00}(\bar Q_m)$; the representation is unique by
Lemma~\ref{lem:rigidity}(a) and because $w_m\ne0$ on the infinite set
$P_m^\circ$.

\emph{Step 3: no contact components.}  Let $j\in K$.  Choose
$\nu_J\in c_{00}(J)$ with $\psi(\nu_J)=z_j\psi(e_j)$ for $\psi\in\Psi$ and
put $\nu:=-z_je_j+\nu_J$, so $\psi(\nu)=0$ for $\psi\in\Psi$.  For small
$s>0$, $\eta+s\nu=c_\eta\big((z-sz_je_j/c_\eta+s\nu_J/c_\eta)+Ue\big)$ has
cube part of norm $\le1$ (the contact moves inward), so
$q^{**}(\eta+s\nu)\le c_\eta=a(\eta+s\nu)$ and $E_q(s\nu)=0$.  As in
Step~1, $E_m(sR_m\nu)=o(s^2)$ and $f(\nu)=0$.  Hence $g(\nu)=0$.  But
$g(\nu)=-z_jb_j$, because $b(\nu_J)=0$ and
$(\omega_m+c_mw_m)(R_m\nu)=0$.  So $b_j=0$.

\emph{Step 4: no degenerate components.}  Let $k_0\in \mathrm{Dg}_{m_0}$ and
$s_0:=\sgn w_{m_0}(k_0)$.  Choose $\nu\in c_{00}(J)$ with
$u_{k_0,m_0}(\nu)=s_0$ and $\psi(\nu)=0$ for all other $\psi\in\Psi$.  Then
$R_{m_0}\nu$ vanishes on $\bar Q_{m_0}\setminus\{k_0\}$, equals
$s_0\lambda_{k_0,m_0}$ at $k_0$, and has zero peak sum.  For $s>0$ the
first bound of Lemma~\ref{lem:overshoot} contributes nothing at $k_0$
(there $-s_0\cdot s\,s_0\lambda_{k_0,m_0}<0$) and nothing at the other
degenerate peaks; the non-degenerate peaks give $o(s^2)$ by (MS).  Other
blocks and the base are as in Step~1.  So $\varphi(s\nu)=o(s^2)$ for
$s>0$ and $g(\nu)=0$; but
$g(\nu)=\omega_{m_0}(k_0)\lambda_{k_0,m_0}s_0$.  So $\omega_{m_0}(k_0)=0$.

\emph{Step 5: balance.}  Fix $m_0$.  Choose $\nu_J\in c_{00}(J)$ with
$\psi(\nu_J)=-\psi(\eta)/\sigma_{m_0}$ for $\psi\in\{u_{k,m_0}:k\in\bar
Q_{m_0}\}\cup\{\pi_{m_0}\}$ and $\psi(\nu_J)=0$ for the other elements of
$\Psi$, and put $\nu:=\eta+\nu_J\in X^{**}$.  Then
$f(\nu)=1+\sum_mw_m(R_m\nu_J)=1-w_{m_0}(\zeta_{m_0})/\sigma_{m_0}=0$.  The
base: $\eta+t\nu=(1+t)(\eta+t\nu_J/(1+t))$, so $E_q(t\nu)=0$ for small
$|t|$ by Lemma~\ref{lem:flat} and homogeneity.  Blocks $m\ne m_0$:
$R_m^{**}(t\nu)=t\zeta_m+tR_m\nu_J$ with $R_m\nu_J$ vanishing on $\bar Q_m$
and of zero peak sum.  Block $m_0$: $R^{**}_{m_0}(t\nu)=t(1-1/\sigma_{m_0})
\zeta_{m_0}+t\delta''$ with $\delta'':=R_{m_0}\nu_J+\zeta_{m_0}/\sigma_{m_0}$
vanishing on $\bar Q_{m_0}$ and of zero peak sum.  Hence $\varphi(t\nu)=
o(t^2)$, and $g(\nu)=0$ by Lemma~\ref{lem:slice}.  Writing
$v_m:=\omega_m+c_mw_m$ we compute $g(\nu)=g(\eta)+b(\nu_J)+\sum_m
v_m(R_m\nu_J)=-v_{m_0}(\zeta_{m_0})/\sigma_{m_0}$, since $v_m(\delta)=0$
whenever $\delta$ vanishes on $\bar Q_m$ and has zero peak sum.  Thus
$v_m(\zeta_m)=0$ for all $m$.  As $\omega_m$ lives on $Q_m$, where
$\zeta_m(k)=\sigma_m\Phi_m(k)^2w_m(k)/C_m$, we have
$\omega_m(\zeta_m)=\sigma_md_m(\omega_m)$, and $v_m(\zeta_m)=\sigma_m(d_m+
c_m)$; hence $c_m=-d_m$.  Finally $b(\eta)=g(\eta)-\sum_mv_m(\zeta_m)=0$.
\end{proof}

\begin{remark}\label{rem:tameNA}
At a norm-attaining $f'$ with C-tame normer $x'$ (i.e.\ finitely many
strict non-peaks and degenerate peaks, and (MS)), Theorem~\ref{thm:tamefibre}
says $\Cset(f')\subseteq S(f')$.  This is a \emph{linear obstruction at the
approximants}: if $f_n\to f$ are norm attaining with C-tame normers, then
$\Li_n\Cset(f_n)\subseteq\Li_nS(f_n)$.  It will be used in
Section~\ref{sec:resonance}.
\end{remark}

\subsection{Transfer peaks and the mass-weighted coefficient}
In this subsection $\eta=\xi$.  The \emph{canonical truncations} of $f$
are defined when $a\in c_{00}$: let $z^N_j:=z_j$ for $j\in F\cup[1,N]$ and
$z^N_j:=0$ otherwise, $\tilde x_N:=z^N+Ue\in c_0$, $x'_N:=\tilde
x_N/p(\tilde x_N)$ and $f'_N:=\nabla p(x'_N)$.  By
Proposition~\ref{prop:smooth}(c), $q(\tilde x_N)=1$, $\nabla q(\tilde
x_N)=a$ and $f'_N=a+\sum_mR_m^*w'_{m,N}$ with $w'_{m,N}=J_m(R_mx'_N)$; by
Proposition~\ref{prop:approximants}, $f'_N\to f$.  We write
$c_N:=q(x'_N)\to q_0$, and $M'_{m,N},C'_{m,N},\sigma'_{m,N},\dots$ for the
data at $x'_N$; they converge to those at $\xi$
(Proposition~\ref{prop:continuity}).

\begin{theorem}[Transfer peaks]\label{thm:transfer}
Let $f\in S_{p^*}$ with $a\in c_{00}$, and let
$g=b+\sum_mR_m^*(\omega_m-d_mw_m)\in S(f)$ be a balanced finite certificate
such that $(f,g)$ is contractive and $\Gamma_w(g)\le1$.  Then for every
$\rho\in(0,1)$ there are $g'_N\in X^*$ with $g'_N\to g$ and
$\rho g'_N\in\Cset(f'_N)$ for $N$ large.  In particular
$(f,\rho g)\in\cl\NA((c_0,p),\ell_2^2)$, and
all such mates are recovered simultaneously along the canonical truncations,
which do not depend on $g$.
\end{theorem}

\begin{proof}
Fix $\rho\in(0,1)$ and $\varepsilon_0\in(0,\frac12]$ with $\rho^2(\Gamma_w(g)+4\varepsilon_0)\le\rho$, and
put $\ell_*:=\Gamma_w(g)/2+\varepsilon_0$.  Let $\gamma_0:=\frac12\min_m\min_{k\in\supp
\omega_m}\gap_m(k)>0$.

\emph{Step 1: the transported natural certificate.}  Put
$b'_N:=b-(b(x'_N)/c_N)a$, $d'_{m,N}:=\ip{D_mw'_{m,N}}{D_m\omega_m}/C'_{m,N}$,
$v'_{m,N}:=\omega_m-d'_{m,N}w'_{m,N}$ and
$g'_N:=b'_N+\sum_mR_m^*v'_{m,N}$.  Since $b(x'_N)\to b(\xi)=0$,
$d'_{m,N}\to d_m$ and $R_m^*w'_{m,N}\to R_m^*w_m$, we have $g'_N\to g$.  For
$N$ large every $k\in\supp\omega_m$ has $M'_{m,N}-|w'_{m,N}(k)|\ge\gamma_0$.
\emph{Base:} $b'_N$ is supported in the finite set $F$, where
$z^N=\sgn a$.  The computation in the proof of Lemma~\ref{lem:base}, with
$\zh$ replaced by $\tilde x_N=z^N+Ue$, gives
$q^*(a+tb'_N)=1+t\,b'_N(\tilde x_N)+\nu\Psi(tU^*b'_N/\nu)$ as long as no
coordinate of $a+tb'_N$ changes sign on $F$; and
$b'_N(\tilde x_N)=b'_N(x'_N)/c_N=0$.  Since $P^\perp U^*b'_N=P^\perp U^*b$
and $b'_N\to b$, the bounds for $\Psi$ give
$q^*(a+tb'_N)\le1+\frac{t^2}2H_b\theta(t)$, where here and below
$\theta(t)\to1$ as $t\to0$, uniformly in large $N$.
\emph{Blocks:} Lemma~\ref{lem:block}(c),(d) at $x'_N$ give
$\|w'_{m,N}+tv'_{m,N}\|_\infty\le(1-td'_{m,N})M'_{m,N}$ and
$\|D_m(w'_{m,N}+tv'_{m,N})\|\le C'_{m,N}+td'_{m,N}M'_{m,N}+\frac{t^2}2
H'_{m,N}\theta(t)$, with $H'_{m,N}:=(\|D_m\omega_m\|^2-d'^2_{m,N})/
C'_{m,N}\to H_m$.

\emph{Step 2: transfer peaks at $\xi$.}  Fix $m$ and put
$\varsigma_m:=\sgn(H_m/2-\ell_*)\in\{-1,0,1\}$; if $\varsigma_m=0$ set
$\tau_m:=0$ and skip this step.  Choose $\gamma_m\in(\max_{k\in\supp\omega_m}
|w_m(k)|,M_m)$ with $|w_m(k)|\ne\gamma_m$ for all $k$, let
$\mathfrak L_m:=\{k\notin\supp\omega_m:|w_m(k)|\ge\gamma_m\}\supseteq P_m$,
$y^0_m:=\sgn(w_m)1_{\mathfrak L_m}$, $\hat\pi_m:=R_m^*y^0_m$ and
$\tau_{*,m}:=(\hat\pi_m(\xi)/q_0)a-\hat\pi_m$, so $\tau_{*,m}(\xi)=0$.  For
$\delta_0\in(0,1]$ let $\mathcal T_m:=\varsigma_m\tau_{*,m}+\delta_0(q^*(\tau_{*,m})+1)
a$ and $\ell_m:=q^*(\mathcal T_m)$; then $\mathcal T_m(\xi)=\delta_0(q^*(\tau_{*,m})+1)q_0>0$.
For $\delta_1\in(0,\mathcal T_m(\xi)/(2\ell_mq_0))$ choose, by density of the tails
of $(u_{k,m})_k$, an index $k_*=k_{*,m}\notin\supp\omega_m$ with
$q^*(u_{k_*,m}-\mathcal T_m/\ell_m)\le\delta_1$, so large that $\ell_m\Phi_m(k_*)
M_m/(mC_m)\le\frac12$ and $\vartheta_m\Phi_m(k_*)<\mathcal T_m(\xi)/(2\ell_m)$.  Then
$u_{k_*,m}(\xi)\ge \mathcal T_m(\xi)/\ell_m-\delta_1q_0\ge \mathcal T_m(\xi)/(2\ell_m)>
\vartheta_m\Phi_m(k_*)$, so $k_*$ is a non-degenerate peak with
$w_m(k_*)=M_m$ and margin $\mu_*>0$.  Put
\begin{gather*}
 y_m:=y^0_m+\varsigma_m\frac{\ell_m}{\lambda_{k_*,m}}e_{k_*},\qquad
 R_m^*y_m=\hat\pi_m+\varsigma_m\ell_mu_{k_*,m},\\
 e_{\infty,m}:=R_m^*y_m-\frac{(R_m^*y_m)(\xi)}{q_0}a .
\end{gather*}
Since $R_m^*y_m=(\text{multiple of }a)+\varsigma_m(\ell_mu_{k_*,m}-\mathcal T_m)$, we get
$e_{\infty,m}=\varsigma_m\big[(\ell_mu_{k_*,m}-\mathcal T_m)-\frac{(\ell_mu_{k_*,m}-
\mathcal T_m)(\xi)}{q_0}a\big]$ and $q^*(e_{\infty,m})\le2\ell_m\delta_1$.  Also
$\ell_m\mu_*\le\ell_mu_{k_*,m}(\xi)\le \mathcal T_m(\xi)+\ell_m\delta_1q_0$.  Hence the
\emph{inefficiency} $\epsilon_{1,m}:=\ell_m\mu_*/q_0+q^*(e_{\infty,m})\le
\delta_0(q^*(\tau_{*,m})+1)+3\ell_m\delta_1$ can be made as small as we
wish by choosing first $\delta_0$ and then $\delta_1$ (note $\ell_m\le
(1+\delta_0)(q^*(\tau_{*,m})+1)$).  Put
$e_{y,m}:=1+\ip{D_my_m}{D_mw_m}/C_m$.  As $y^0_mw_m\ge0$ and
$\ell_m\Phi_m(k_*)M_m/(mC_m)\le\frac12$, we have $e_{y,m}\ge\frac12$.
Finally let $\tau_m:=(H_m/2-\ell_*)/e_{y,m}$, whose sign is $\varsigma_m$, and
note $|\tau_m|\le H_m+2$ independently of $\delta_0,\delta_1,k_*$.  We
choose $\delta_0,\delta_1$ so that $\sum_m(H_m+2)\epsilon_{1,m}\le\varepsilon_0/2$.

\emph{Step 3: the limiting levels.}  Put
$\Lambda_b:=H_b/2+\sum_m\big(\tau_m\sigma_me_{y,m}/q_0+|\tau_m|
\epsilon_{1,m}\big)$ and $\Lambda_m:=H_m/2-\tau_me_{y,m}=\ell_*$.  Using
$\sum_m\sigma_m=1-q_0$,
\[
 q_0\Lambda_b=\frac{\Gamma_w}2-\ell_*(1-q_0)+q_0\sum_m|\tau_m|\epsilon_{1,m},
 \quad\text{so}\quad
 \Lambda_b=\ell_*-\frac{\varepsilon_0}{q_0}+\sum_m|\tau_m|\epsilon_{1,m}\le \ell_* .
\]

\emph{Step 4: transport of the transfer data.}  For $N$ large let
$\mathfrak L_{m,N}:=\{k\notin\supp\omega_m:|w'_{m,N}(k)|\ge\gamma_m\}$ and
$y_{m,N}:=\sgn(w'_{m,N})1_{\mathfrak L_{m,N}}+\varsigma_m(\ell_m/\lambda_{k_*,m})
e_{k_*}$.  For large $N$, $k_*$ is a peak of $w'_{m,N}$ of sign $+1$ with
margin $\mu'_{*,N}\to\mu_*$, and the peak set of $w'_{m,N}$ is contained
in $\mathfrak L_{m,N}$.  Since $w'_{m,N}\to w_m$ coordinatewise and
$|w_m(k)|\ne\gamma_m$, $y_{m,N}\to y_m$ coordinatewise and boundedly; hence
$R_m^*y_{m,N}\to R_m^*y_m$ in $\ell_1$ and $D_my_{m,N}\to D_my_m$ in
$\ell_2$.  Put $\beta_{m,N}:=((R_m^*y_{m,N})(x'_N)/c_N)a$,
$e_{m,N}:=R_m^*y_{m,N}-\beta_{m,N}\to e_{\infty,m}$, and
$e_{y,m,N}:=1+\ip{D_my_{m,N}}{D_mw'_{m,N}}/C'_{m,N}\to e_{y,m}$.  Writing
$R_mx'_N=\sigma'_{m,N}(\alpha'+D_m^2w'_{m,N}/C'_{m,N})$
(Lemma~\ref{lem:threshold}) and using \eqref{eq:margin} at $k_*$, one obtains
the bookkeeping identity
\begin{equation}\label{eq:bookkeeping}
 y_{m,N}(R_mx'_N)=\sigma'_{m,N}e_{y,m,N}+\varsigma_m\ell_m\mu'_{*,N}.
\end{equation}
Indeed $\ip{\sgn(w')1_{\mathfrak L_{m,N}}}{\alpha'}=1$, and
\[
 \frac{\ell_m}{\lambda_{k_*}}(R_mx'_N)(k_*)=\ell_mu_{k_*}(x'_N)=
 \ell_m\mu'_{*,N}+\frac{\ell_m}{\lambda_{k_*}}\,\frac{\sigma'\Phi(k_*)^2M'}{C'},
\]
the last term being the contribution of $e_{k_*}$ to
$\sigma'\ip{D y_{m,N}}{Dw'}/C'$.

\emph{Step 5: the decomposition.}  For $t\in\R$ put
$\varepsilon_m:=\tau_m\rho^2t^2$ and
\[
 A(t):=a+t\rho b'_N+\sum_m\varepsilon_mR_m^*y_{m,N},\qquad
 W_m(t):=w'_{m,N}+t\rho v'_{m,N}-\varepsilon_my_{m,N},
\]
so that $A(t)+\sum_mR_m^*W_m(t)=f'_N+t\rho g'_N$.
\emph{Base.}  By \eqref{eq:bookkeeping},
$A(t)=(1+E)a+t\rho b'_N+\sum_m\varepsilon_me_{m,N}$ with
$E:=\rho^2t^2\sum_m\big(\tau_m\sigma'_{m,N}e_{y,m,N}+|\tau_m|\ell_m
\mu'_{*,N}\big)/c_N$.  By homogeneity and Step~1,
\[
 q^*(A(t))\le1+E+\frac{\rho^2t^2}{2(1+E)}H_b\theta(t)
 +\rho^2t^2\sum_m|\tau_m|q^*(e_{m,N})
 =1+\rho^2t^2\big(\Lambda_b+o(1)\big),
\]
where $o(1)\to0$ as $t\to0$ and $N\to\infty$.
\emph{Blocks, sup norm.}  For $|t|$ small and $N$ large,
$\|W_m(t)\|_\infty\le(1-t\rho d'_{m,N})M'_{m,N}-\varepsilon_m$: on
$\mathfrak L_{m,N}\setminus\{k_*\}$ the coordinates are
$(1-t\rho d')w'(k)-\varepsilon_m\sgn w'(k)$, whose moduli are at most this
bound; at $k_*$ the value is $(1-t\rho d')M'-\varepsilon_m(1+\varsigma_m
\ell_m/\lambda_{k_*})$, which lies between $\pm((1-t\rho
d')M'-\varepsilon_m)$ once $|\varepsilon_m|\ell_m/\lambda_{k_*}\le M'$ (here
$\varsigma_m\varepsilon_m\ge0$ is used); on $\supp\omega_m$ the moduli are at
most $(1-t\rho d')(M'-\gamma_0)+\rho|t|\|\omega_m\|_\infty$; elsewhere at
most $(1-t\rho d')\gamma_m$; both are below the bound for small $|t|$ and
large $N$ because $M'_{m,N}\to M_m>\gamma_m$.
\emph{Blocks, Hilbert part.}  With $h_0:=D_m(w'+t\rho v')$, inequality
\eqref{eq:hilbertineq} gives
$\|h_0-\varepsilon_mD_my_{m,N}\|\le\|h_0\|-\varepsilon_m\ip{h_0/\|h_0\|}{D_m
y_{m,N}}+O(t^4)$, and $\ip{h_0/\|h_0\|}{D_my_{m,N}}=
\ip{D_mw'}{D_my_{m,N}}/C'+O(|t|)$ uniformly in $N$.  Together with Step~1,
\[
 N_m(W_m(t))\le1+\rho^2t^2\Big(\frac{H'_{m,N}\theta(t)}2-\tau_me_{y,m,N}
 \Big)+O(|t|^3)=1+\rho^2t^2(\ell_*+o(1)).
\]
By Lemma~\ref{lem:dualball} and Step~3, there are $t_4>0$ and $N_4$ with
$p^*(f'_N+t\rho g'_N)\le1+\rho^2t^2(\ell_*+\varepsilon_0)\le1+\frac{t^2}2\rho^2
(\Gamma_w+4\varepsilon_0)\le1+\rho t^2/2\le s(t)$ for $|t|\le t_5:=\min(t_4,
2\sqrt{1-\rho},1)$ and $N\ge N_4$.

\emph{Step 6: large $t$.}  Since $(f,g)$ is contractive,
$p^*(f'_N+t\rho g'_N)\le s(\rho t)+\varepsilon_N+\rho|t|\varepsilon'_N$ with
$\varepsilon_N:=p^*(f'_N-f)\to0$ and $\varepsilon'_N:=p^*(g'_N-g)\to0$; by
Lemma~\ref{lem:slack}(a) this is at most $s(t)$ for $|t|\ge t_5$ and $N$
large.  Hence $\rho g'_N\in\Cset(f'_N)$ for $N$ large, and $g'_N\to g$.
\end{proof}

\begin{remark}\label{rem:transfer}
(a) If $\Gamma_{\max}(g)\le1$, no transfer peaks are needed ($\tau_m=0$);
this is the natural-certificate case, and together with
Corollary~\ref{cor:check} it re-proves the local-certificate theorem of
\cite{Check}.  (b) The transfer peak is the genuinely infinite-dimensional
ingredient: lowering a block's sup norm requires lowering every near-peak
coordinate, and the image $\hat\pi_m$ of such a uniform lowering is a fixed
vector of $Y$ which the base cannot absorb; density of the tails of
$(u_{k,m})_k$ provides a robust deep peak whose functional converts it into
a multiple of $a$ with arbitrarily small inefficiency.  In finite-dimensional
models of the construction this mechanism is absent, and numerical
computations in such models (numerical evidence only, not a theorem) give
second-order coefficients strictly larger than $\Gamma_w$.
(c) For every balanced finite certificate $g$ at $f$ with $a\in c_{00}$,
$\limsup_{t\to0}2(p^*(f+tg)-1)/t^2\le\Gamma_w(g)$; this is proved in
Proposition~\ref{prop:onesidedupper}, by running the rebalancing at $\xi$
instead of $x'_N$, and Theorem~\ref{thm:onesided} identifies the exact
one-sided coefficients at \textup{(BT)} points.
\end{remark}

\begin{remark}[$\Gamma_w$ and shifts]\label{rem:gammaw}
For every shifted certificate $c=(b,\omega,\theta)$ at $f$ one has
$H^{\rm sh}(c)\ge\Gamma_w(g_c)$.  Indeed, if $\Lambda:=H^{\rm sh}(c)$, then
$2\theta_m\le\Lambda-H_m$ for every $m$, and
$\Lambda\ge h(b)-2\sum_m\theta_m\sigma_m/q_0\ge h(b)-\sum_m(\Lambda-H_m)
\sigma_m/q_0$; multiplying by $q_0$ and using $q_0+\sum_m\sigma_m=1$ gives
$\Lambda\ge\Gamma_w$.  Thus, when $a\in c_{00}$, Theorem~\ref{thm:transfer}
covers every $g_c\in\Cset(f)$ with $c$ a shifted certificate and
$H^{\rm sh}(c)\le1$ (the direction of a finite certificate always lies in
$S(f)$).
\end{remark}

\subsection{The sharp second-order coefficient at tame supports}

\begin{lemma}[Exact local block formula]\label{lem:localformula}
Fix a block and drop $m$.  Recall that $\sum_k\Phi(k)^2<1$.  Let $y\in\ell_1$, let
$P\ne\emptyset$ be a set of indices with complement $Q$, and
$s\in\{\pm1\}^P$.  Put $S:=\sum_{k\in P}s_ky_k$,
$Y:=(\sum_{k\in Q}y_k^2/\Phi(k)^2)^{1/2}$ and $F_2:=\sum_{k\in P}\Phi(k)^2\in
(0,1)$, and assume $0<S$, $Y<\infty$ and $Y\le S$.  Define
\[
 \Lambda(S,Y):=\frac{S-\sqrt{F_2}\sqrt{S^2-(1-F_2)Y^2}}{1-F_2},\qquad
 \lambda:=\Lambda(S,Y),\qquad \rho_0:=(S/\lambda-1)/F_2 .
\]
Then $\lambda>0$ and
\begin{equation}\label{eq:overshootformula}
 |y|\le\lambda+2\sum_{k\in P}\big(\lambda\rho_0\Phi(k)^2-s_ky_k\big)_+ .
\end{equation}
Moreover, if $\zeta\ne0$ with norming functional $w$, and $(P,s)$ are its
peak set and signs, then with $c_Q:=\|D(w1_Q)\|_2^2$:
$S_P(\zeta)=|\zeta|(1+F_2M/C)$, $Y_Q(\zeta)=|\zeta|\sqrt{c_Q}/C<S_P(\zeta)$,
$\Lambda(S_P(\zeta),Y_Q(\zeta))=|\zeta|$, $\lambda\rho_0=|\zeta|M/C$, and at
this point
\[
 \partial_S\Lambda=M,\qquad \partial_Y\Lambda=CY/|\zeta|,\qquad
 \partial_S^2\Lambda=\frac{\sqrt{F_2}\,Y^2}{(S^2-(1-F_2)Y^2)^{3/2}}
 =\frac{C\,c_Q}{|\zeta|F_2}.
\]
\end{lemma}

\begin{proof}
$\Delta:=S^2-(1-F_2)Y^2\ge F_2S^2>0$; $\lambda$ is the smaller root of
$(1-F_2)l^2-2Sl+(S^2+F_2Y^2)=0$, i.e.\ of $(S-l)^2+F_2Y^2=F_2l^2$, and
$\lambda>0$ because $\sqrt{F_2\Delta}<S$.  The root equation reads
$\rho_0^2F_2+Y^2/\lambda^2=1$.  Define $h\in\ell_2$ by
$h_k:=\rho_0\Phi(k)s_k$ ($k\in P$), $h_k:=y_k/(\lambda\Phi(k))$ ($k\in Q$);
then $\|h\|_2=1$, and $\alpha:=y/\lambda-Dh$ vanishes on $Q$, with
$\|\alpha\|_1=\sum_Ps_k\alpha_k+2\sum_P(-s_k\alpha_k)_+=S/\lambda-\rho_0F_2
+\frac2\lambda\sum_P(\lambda\rho_0\Phi(k)^2-s_ky_k)_+$, and
$S/\lambda-\rho_0F_2=1$.  Hence $y=\lambda\alpha+\lambda Dh\in\lambda
\|\alpha\|_1(B_{\ell_1}+D(B_{\ell_2}))$, which is \eqref{eq:overshootformula}.

At $\zeta$, Lemma~\ref{lem:threshold} gives $s_k\zeta_k=|\zeta|(|\alpha_k|+
\Phi(k)^2M/C)$ on $P$ and $\zeta_k=|\zeta|\Phi(k)^2w_k/C$ on $Q$, whence
the formulas for $S$ and $Y$; $Y<S$ because $\sqrt{c_Q}\le C$ and $F_2M>0$.
Using $C^2=M^2F_2+c_Q$ one checks that $l=|\zeta|$ satisfies
$(S-l)^2+F_2Y^2=F_2l^2$; since $|\zeta|(1-F_2)<S$, it is the smaller root.
Then $\rho_0=M/C$, and from the root equation
$\sqrt{F_2\Delta}=S-(1-F_2)\lambda=\lambda F_2(1+\rho_0)$, i.e.\
$\sqrt\Delta=|\zeta|\sqrt{F_2}/C$.  Differentiating,
$\partial_S\Lambda=(1-\sqrt{F_2}S/\sqrt\Delta)/(1-F_2)=(1-C-F_2M)/(1-F_2)=M$,
$\partial_Y\Lambda=\sqrt{F_2}Y/\sqrt\Delta=CY/|\zeta|$, and
$\partial^2_S\Lambda=\sqrt{F_2}(S^2-\Delta)/((1-F_2)\Delta^{3/2})=
\sqrt{F_2}Y^2/\Delta^{3/2}=Cc_Q/(|\zeta|F_2)$.
\end{proof}

\begin{proposition}[Lower bound]\label{prop:lowerbound}
Let $I$ be finite and let $f\in S_{p^*}$ satisfy: $a\in c_{00}$, $K$ is
finite, every $Q_m$ is finite, there are no degenerate peaks, and
\textup{(MS)} holds at $\xi$.  Then for every balanced finite certificate
$g\in S(f)$,
\[
 \liminf_{t\to0}\frac{p^*(f+tg)-1}{t^2}\ge\frac{\Gamma_w(g)}2 .
\]
In particular $\Gamma_w(g)\le1$ for every $g\in\Cset(f)\cap S(f)$.
\end{proposition}

\begin{proof}
Write $g=b+\sum_mR_m^*v_m$, $v_m=\omega_m-d_mw_m$.  For $\chi,\chi_2\in X$
put $\eta_t:=\xi+t\chi+t^2\chi_2$.  Since $g(\xi)=0$ and
$p^{**}(\eta_t)=f(\eta_t)+\Delta(\eta_t)$ with $f(\eta_t)=1+O(t)$,
\begin{equation}\label{eq:testratio}
 p^*(f+tg)\ge\frac{(f+tg)(\eta_t)}{p^{**}(\eta_t)}
 =1+t^2g(\chi)-\Delta(\eta_t)+O(t^3)
 \quad\text{whenever }\Delta(\eta_t)=O(t^2).
\end{equation}
By Lemma~\ref{lem:excess},
$\Delta(\eta_t)=E_q(t\chi+t^2\chi_2)+\sum_mE_m(t\delta_{1,m}+t^2\delta_{2,m})$
with $\delta_{i,m}:=R_m\chi_i$.

\emph{Base test.}  Let $E_F:=\spn\{U^*e_j^*:j\in F\}\ni e$ and fix
$k_\perp\in E_F$ with $\ip{k_\perp}e=0$.  Put
$\kappa_2:=\nu\|k_\perp\|^2/(2q_0)$, $\kappa_2':=-\kappa_2\|a\|_1/\nu$, and
let $k_{2\perp}\in E_F$ solve $\ip{k_{2\perp}}{U^*e_j^*}=-\kappa_2z_j-
\kappa'_2(Ue)_j$ for $j\in F$ (the Gram matrix of the independent vectors
$U^*e_j^*$ is invertible); then $\ip{k_{2\perp}}{e}=(-\kappa_2\|a\|_1-
\kappa'_2\nu)/\nu=0$.  Put $k_2:=\kappa'_2e+k_{2\perp}$, so that
$(Uk_2)_j=-\kappa_2z_j$ on $F$.  Let $\chi:=Uk_\perp+\chi_c$ with
$\chi_c\in c_{00}(J)$ to be chosen, and $\chi_2:=(Uk_2)1_{\N\setminus F}$.
With $\mathbf h:=q_0e+tk_\perp+t^2k_2$ and $Z:=\eta_t-U\mathbf h=q_0z+t
\chi_c-t^2(Uk_2)1_F$ we have $|Z_j|=q_0+t^2\kappa_2$ on $F$ and $|Z_j|\le
q_0$ off $F$ for small $t$ ($\chi_c$ moves finitely many free
coordinates), while
$\|\mathbf h\|=q_0+t^2(\kappa'_2+\|k_\perp\|^2/(2q_0))+O(t^3)=q_0+t^2\kappa_2
(1-\|a\|_1)/\nu+O(t^3)=q_0+t^2\kappa_2+O(t^3)$, using $\|a\|_1+\nu=1$.
Since $B_{q^{**}}=B_{\ell_\infty}+U(B_H)$, $q^{**}(\eta_t)\le q_0+t^2\kappa_2
+O(t^3)$; and $a(\chi)=\ip{U^*a}{k_\perp}=0=a(\chi_2)$.  Hence
$E_q(t\chi+t^2\chi_2)\le t^2\nu\|k_\perp\|^2/(2q_0)+O(t^3)$, and
$b(\chi)=\ip{U^*b}{k_\perp}$.

\emph{Block test.}  Fix $m$ and drop it.  Apply \eqref{eq:overshootformula}
to $y:=\zeta+t\delta_1+t^2\delta_2$ with the peak set and signs of $\zeta$.
As $Q$ is finite, $G(y):=\Lambda(S_P(y),Y_Q(y))$ is a smooth function of the
finitely many numbers $(S_P(y),(y_k)_{k\in Q})$ near those of $\zeta$
($\Lambda$ depends smoothly on $(S,Y^2)$).  By Lemma~\ref{lem:localformula}
its differential at $\zeta$ is $\delta\mapsto MS_P(\delta)+\sum_{k\in
Q}w_k\delta_k=w(\delta)$, so
$G(y)=|\zeta|+w(t\delta_1+t^2\delta_2)+\frac{t^2}2\mathcal H(\delta_1)+
O(t^3)$, where $\mathcal H$ is the Hessian of $G$ at $\zeta$.  The overshoot:
$\lambda(y)\rho_0(y)=|\zeta|M/C+O(t)$ and, by \eqref{eq:margin},
$s_ky_k=\lambda_k\mu_k+|\zeta|\Phi(k)^2M/C+ts_k\delta_{1,k}+t^2s_k
\delta_{2,k}$ on $P$, with $|\delta_{i,k}|\le\lambda_kq(\chi_i)$; so each
overshoot term is at most $\lambda_k(|t|B-\mu_k)_+$ for a constant $B$,
and their sum is $o(t^2)$ by (MS), all peaks being non-degenerate.  Hence
$E_m(t\delta_1+t^2\delta_2)\le\frac{t^2}2\mathcal H_m(\delta_{1,m})+o(t^2)$.

\emph{Decoupling.}  $\mathcal H_m(\delta)$ and $v_m(\delta)$ depend only on
$\sigma_1:=S_{P_m}(\delta)$ and $\delta|_{Q_m}$.  Since $F\cup K$ is finite,
Lemma~\ref{lem:rigidity}(b) (with $P'_m=P_m$, which is cofinite) allows us to
choose $\chi_c\in c_{00}(J)$ so that, for the fixed $k_\perp$, the finitely
many numbers $u_{k,m}(\chi)$ ($k\in Q_m$) and $\pi_m(\chi)=S_{P_m}(R_m\chi)$
take arbitrary prescribed values.  By \eqref{eq:testratio},
\begin{multline*}
 \liminf_{t\to0}\frac{p^*(f+tg)-1}{t^2}\\
 \ge\frac12\Big[\sup_{k_\perp}
 \Big(2\ip{U^*b}{k_\perp}-\frac{\nu}{q_0}\|k_\perp\|^2\Big)+\sum_m
 \sup_{\delta}\big(2v_m(\delta)-\mathcal H_m(\delta)\big)\Big].
\end{multline*}
\emph{Base supremum.}  $U^*b\in E_F$; the maximizer is
$k_\perp=(q_0/\nu)P^\perp U^*b\in E_F\cap e^\perp$, with value
$(q_0/\nu)\|P^\perp U^*b\|^2=q_0H_b$.

\emph{Block suprema.}  Drop $m$.  Suppose first $c_Q>0$, so $Y>0$.  Write
$\zeta':=D^{-1}\delta_Q$, $\hat h:=D^{-1}\zeta_Q/Y=D w_Q/\sqrt{c_Q}$,
$\sigma_2:=\ip{\hat h}{\zeta'}$, $\zeta'_\perp:=\zeta'-\sigma_2\hat h$ and
$r:=Y\sigma_1-S\sigma_2$.  Since $\Lambda$ is $1$-homogeneous, its Hessian
is $(\partial_S^2\Lambda/Y^2)(Y,-S)^{\otimes2}$, and
$Y_Q(\zeta+\delta)=Y+\sigma_2+\|\zeta'_\perp\|^2/(2Y)+O(\|\delta\|^3)$;
with Lemma~\ref{lem:localformula},
$\mathcal H(\delta)=(C/|\zeta|)\|\zeta'_\perp\|^2+(\partial_S^2\Lambda/Y^2)
r^2$.  Next, with $\tilde\omega:=(\omega-dw)1_Q$,
$v(\delta)=\ip{D\tilde\omega}{\zeta'_\perp}+\big(\ip{D\tilde\omega}{\hat h}
\sigma_2-dM\sigma_1\big)$.  The linear form in brackets vanishes at
$(\sigma_1,\sigma_2)=(S,Y)$, since there it equals $v(\zeta)=\omega(\zeta)
-d|\zeta|=0$ ($\omega(\zeta)=|\zeta|d$ because $\zeta=|\zeta|D^2w/C$ on $Q$);
hence it equals $-(dM/Y)r$.  As $\zeta'_\perp\in\hat h^\perp$ and $r$ are
independent variables,
\begin{align*}
 \sup_\delta(2v-\mathcal H)&=\frac{|\zeta|}C\|P_{\hat h^\perp}D\tilde\omega
 \|^2+\frac{d^2M^2}{\partial_S^2\Lambda}
 =\frac{|\zeta|}C\Big(\|D\omega\|^2-\frac{d^2C^2}{c_Q}\Big)+
 \frac{d^2M^2|\zeta|F_2}{Cc_Q}\\
 &=\frac{|\zeta|}C(\|D\omega\|^2-d^2),
\end{align*}
where we used $\ip{D\omega}{Dw}=dC$, $\|P_{\hat h^\perp}D\tilde\omega\|^2=
\|D\omega\|^2-d^2C^2/c_Q$ (a direct computation) and $C^2=c_Q+M^2F_2$.  This
is $\sigma_mH_m$.  If $c_Q=0$, then $w_Q=0$, $d=0$, $\Lambda(S,\cdot)$
depends on $Y^2$ with $\partial\Lambda/\partial(Y^2)=C/(2|\zeta|)$ at
$\zeta$, $\mathcal H(\delta)=(C/|\zeta|)\|D^{-1}\delta_Q\|^2$,
$v(\delta)=\ip{D\omega}{D^{-1}\delta_Q}$, and the supremum is
$(|\zeta|/C)\|D\omega\|^2=\sigma_mH_m$ directly.

Summing, $\liminf_{t\to0}(p^*(f+tg)-1)/t^2\ge\frac12(q_0H_b+\sum_m\sigma_m
H_m)=\Gamma_w(g)/2$.  If $g\in\Cset(f)$, then $(p^*(f+tg)-1)/t^2\le
(s(t)-1)/t^2\to\frac12$.
\end{proof}

\begin{theorem}[Tame supports are recoverable]\label{thm:tame}
Let $I$ be finite and let $f\in S_{p^*}$ have a C-tame normer without
degenerate peaks, i.e.\ $a\in c_{00}$, $K$ finite, every $Q_m$ finite,
$\mathrm{Dg}_m=\emptyset$ and \textup{(MS)}.  Then
$\Cset(f)\subseteq\Li_N\Cset(f'_N)$ for the canonical truncations
$f'_N$ of $f$.  In particular $f\in\Rec$.
\end{theorem}

\begin{proof}
By Theorem~\ref{thm:tamefibre}, every $g\in\Cset(f)$ is a balanced finite
certificate; by Proposition~\ref{prop:lowerbound}, $\Gamma_w(g)\le1$; by
Theorem~\ref{thm:transfer}, $\rho g\in\Li_N\Cset(f'_N)$ for every $\rho<1$,
and the lower limit is closed.  Corollary~\ref{cor:finite} gives $f\in\Rec$.
\end{proof}

\begin{remark}[On the hypotheses]\label{rem:kinks}
(a) The finiteness of $K$ enters Theorem~\ref{thm:tame} twice: in Step~1 of
Theorem~\ref{thm:tamefibre} and in the decoupling step of
Proposition~\ref{prop:lowerbound}, both of which need the set $J$ of free
coordinates to be cofinite so that Lemma~\ref{lem:rigidity}(b) applies.
Lemma B does not exclude elements of $\spn\{u_{k,m},\pi_m\}\subseteq Y$
supported in $F\cup K$ when $K$ is infinite, and then the lower bound may
fail.  In a finite-dimensional analogue in which every off-support
coordinate is a contact, an exact numerical computation (in that finite
model only; this is not a theorem about Mart\'in's space) gives a
second-order coefficient of a pure block certificate of at most $0.296$ on
both sides, while $\Gamma_w=0.485$; the mechanism is that re-splitting
through the contacts is free at first order for one sign of $t$.  Whether a finite certificate with $\Gamma_2\le1<\Gamma_w$ exists at a
support with infinitely many contacts in Mart\'in's space is \emph{open}.
Theorem~\ref{thm:transfer} would not recover it, but
Corollary~\ref{cor:BTrecovered} does whenever $f$ satisfies \textup{(BT)}:
by Theorem~\ref{thm:onesided} the exact one-sided coefficients are the
minima $\gamma^\pm$ of $\Gamma_w$ over side-admissible decompositions, which
re-split through contacts with the cheap sign, and the proof of the lower
bound there needs no decoupling.  (b) Earlier formulations
assumed $K=\emptyset$; Step~3 of Theorem~\ref{thm:tamefibre} removes this.
(c) The existence of tame non-attaining supports for a \emph{given}
admissible $T$ is not known in general.  In the open part of the bidual
face where $K=\emptyset$, where the residuality argument of
\cite[Section~3]{PrA} for half-peaks applies, they form a meagre set.
For the operator $T$ of Section~\ref{sec:resonance} the canonical
truncations of the special first row are C-tame norm-attaining points
(Proposition~\ref{prop:nonrecovery}).
\end{remark}

\section{The resonance obstruction}\label{sec:resonance}

Throughout this section $I$ is finite.

\subsection{The linear obstruction}
For $f\in S_{p^*}$ let $\mathfrak B(f):=\{g\in\Cset(f)\cap S(f):\Gamma_w(g)
\le1\}$ if $a\in c_{00}$ (then the data of $g\in S(f)$ are unique by
Lemma~\ref{lem:rigidity}(a)), and $\mathfrak B(f):=\emptyset$ otherwise.
The \emph{intrinsic class} $\mathcal K(f)$ is the closed convex hull of
$\Cert^{\rm sh}(f)\cup\mathfrak B(f)$, and the \emph{defect} is
$\Def(f):=\Cset(f)\setminus\mathcal K(f)$.

\begin{proposition}[Intrinsic classes]\label{prop:intrinsic}
\textup{(a)} $\mathcal K(f)$ contains every mate shown to be recoverable in
Sections~\ref{sec:certificates}--\ref{sec:second}: $\cl\Cert(f)$ (hence the
mates of Corollaries~\ref{cor:check} and~\ref{cor:ladder} and
Theorems~\ref{thm:weighted}, \ref{thm:locsplit}, \ref{thm:carriers}),
$\cl\Cert^{\rm sh}(f)$ and $\mathfrak B(f)$.
\textup{(b)} $\mathcal K(f)\subseteq\cl S(f)$, a closed linear subspace of
$\{g:g(\xi)=0\}$.
\textup{(c)} If $a\in c_{00}$, then $\mathcal K(f)\subseteq\Li_N\Cset(f'_N)$
for the canonical truncations; in any case $\cl\Cert^{\rm sh}(f)\subseteq
\Li_n\Cset(f_n)$ for every sequence $f_n\to f$.
\end{proposition}

\begin{proof}
(a) is a summary of the quoted results ($\Cert(f)\subseteq\Cert^{\rm
sh}(f)$).  (b) Conditions (C1), (C2) imply that $g_c\in S(f)$, for
shifted certificates as well, and $\mathfrak B(f)\subseteq S(f)$ by
definition.  (c) Theorem~\ref{thm:shifted} for $\Cert^{\rm sh}$,
Theorem~\ref{thm:transfer} for $\mathfrak B(f)$, and closedness and
convexity of lower limits of convex sets.
\end{proof}

\begin{corollary}[Linear obstruction]\label{cor:linear}
$\Cset(f)\setminus\cl S(f)\subseteq\Def(f)$.  If $F$ and every $Q_m$ are
finite, then $S(f)$ is finite dimensional, and a mate $g$ lies outside
$\mathcal K(f)$ as soon as $g\notin\spn\{e_j^*:j\in F\}+\spn\{u_{k,m}:k\in
Q_m\}+\spn\{R_m^*w_m\}$.
\end{corollary}

When some block has infinitely many strict non-peaks whose vectors are
spread in $\ker\xi$, the closure of $S(f)$ may be all of $\ker\xi$ and the
obstruction is void.  It bites at first rows with few non-peaks.  We now
construct such a first row with an infinite-dimensional fibre.

\subsection{An admissible operator with an exact resonance}
Let $\nu_j:=\|U^*e_j^*\|\to0$ and $J_0:=\{j\ge2:\nu_j\le\frac12\}$
(cofinite).  Put $\bar\alpha:=1/(1+\nu_1)$, $a:=\bar\alpha e_1^*$ (so $q^*(a)=1$),
$e:=U^*e_1^*/\nu_1$, and note $(Ue)_1=\nu_1$.  Fix a bijection
$\ell:\N\times\N\to\N$ with $\ell(1,1)=1$ and $\ell(n,m)<\ell(n',m)$ for $n<n'$, pairwise
disjoint infinite sets $S_l\subseteq J_0$ ($l\in\N$), and put
$h^{(l)}:=\sum_{s\in S_l}2^{-s}e_s^*$, $c_l:=2^{-(l-1)^2}$ (so $c_1=1$ and
$\sum_{l'\ge l}c_{l'}\le2c_l$).  Let $l_0:=\ell(2,1)$, $K_0:=S_{l_0}$,
\[
 z:=e_1+1_{K_0}\in\ell_\infty,\quad \zh:=z+Ue,\quad
 z^{(n)}:=e_1+1_{K_0\cap[1,n]},\quad x^{(n)}:=z^{(n)}+Ue\in c_0 .
\]
Note $\zh_1=x^{(n)}_1=1+\nu_1=1/\bar\alpha$.  For $k\ge2$, $(k,m)\ne(2,1)$, put
$\varpi_{k,m}:=2^{1-m-k}c_{\ell(k,m)}/c_{\ell(1,m)}\le\frac14$ and
$r_{\ell(k,m)}:=\sqrt{\varpi_{k,m}}$, and call $(k,m)$ \emph{exceptional} if
$\varpi_{k,m}\ge\epsilon_*^2$, $\epsilon_*:=1/(72(1+\|U\|))$ (finitely many).
Let $\delta_l:=\min(2^{-l},1/(8(1+\|U\|)))$ for $l=\ell(1,m)$, $\delta_{l_0}:=1$,
$\delta_l:=\min(2^{-l},1/(16(1+\|U\|)))$ for exceptional $l$, and
$\delta_l:=\min(2^{-l},r_l/(1+\|U\|))$ otherwise.  Fix a dense sequence
$(y^{(i)})_{i\ge1}$ in $c_{00}\cap S_{q^*}$ with $y^{(1)}:=e_1^*/q^*(e_1^*)$,
let $\Lambda_i:=\{n:\supp y^{(i)}\cap K_0\cap(n,\infty)\ne\emptyset\}$
(finite), and call $i$ \emph{allowed} at a non-exceptional
$l=\ell(k,m)$ ($k\ge2$, $l\ne l_0$) if
\begin{gather*}
 \supp y^{(i)}\cap S_l=\emptyset,\qquad 3r_l(2|\Lambda_i|+3)\le\tfrac18,\\
 2c_l\le2^{-2s}c_{l'}\delta_{l'}
 \quad\text{for } s\in\supp y^{(i)}\cap S_{l'},\ l'\ne l.
\end{gather*}
Each $i$ is allowed at all large $l$, and $i=1$ is allowed everywhere
($\Lambda_1=\emptyset$ and $9r_l\le\frac18$).  Fix a sequence $(i_n)$ in
which every positive integer occurs infinitely often, and let
$i(n,m):=i_n$ if $i_n$ is allowed at $\ell(n,m)$ and $i(n,m):=1$ otherwise.
Define, with $n_l$ the $q^*$-norm of the bracket,
\begin{itemize}
\item $u_{1,m}:=(e_1^*/q^*(e_1^*)+\delta_lh^{(l)})/n_l$ ($l=\ell(1,m)$) and, for exceptional
$(k,m)$, $u_{k,m}:=(e_1^*/q^*(e_1^*)+\delta_lh^{(l)})/n_l$ ($l=\ell(k,m)$);
\item $u_{2,1}:=(-\kappa_u e_1^*+h^{(l_0)})/n_{l_0}$, with
$\kappa_u:=(\|h^{(l_0)}\|_1+\ip{U^*h^{(l_0)}}{e})/(1+\nu_1)$;
\item for the remaining $(k,m)$, with $l=\ell(k,m)$ and $y:=y^{(i(k,m))}$:
$u_{k,m}:=(y+s_l\bar\alpha e_1^*+\delta_lh^{(l)})/n_l$, where
$|s_l|\le3r_l(2|\Lambda_i|+3)$ is chosen so that $|y(x)+s_l|\ge3r_l$
for every $x$ in the finite set $\{\zh\}\cup\{x^{(n)}:n\in\Lambda_i\}$;
\end{itemize}
and $Te_{n,m}:=c_{\ell(n,m)}u_{n,m}$.  The choice of $s_l$ is possible: the
excluded values form $|\Lambda_i|+1$ intervals of length $6r_l$ inside an
interval of length $6r_l(2|\Lambda_i|+3)$; and $(\bar\alpha
e_1^*)(x)=\bar\alpha x_1=1$ for all these $x$.

\begin{lemma}[The operator $T$]\label{lem:T}
$T$ is admissible, $\Phi_m(k)=2^{-m-k}c_{\ell(k,m)}$, and:
\begin{enumerate}
\item[(P1)] $\supp u_{2,1}=\{1\}\cup K_0$, $u_{2,1}(1)<0<u_{2,1}(j)$ for
$j\in K_0$, and $u_{2,1}(\zh)=0$;
\item[(P2)] $|u_{1,m}(x)|\ge7/9$ for $x\in X_0:=\{\zh\}\cup\{x^{(n)}:n\in\N\}$
and every $m$;
\item[(P3)] $|u_{k,m}(x)|\ge\frac32r_{\ell(k,m)}$ for $x\in X_0$, every
non-exceptional $(k,m)$ with $k\ge2$, $(k,m)\ne(2,1)$, and
$|u_{k,m}(x)|\ge\frac12$ for exceptional $(k,m)$; in both cases
$|u_{k,m}(x)|>\varpi_{k,m}$.
\end{enumerate}
\end{lemma}

\begin{proof}
\emph{Norms.}  $q^*(h^{(l)})\le(1+\|U\|)\|h^{(l)}\|_1\le1+\|U\|$, so
$q^*(\delta_lh^{(l)})\le\frac18$ for $l=\ell(1,m)$, $\le\frac1{16}$ for exceptional
$l$, and $\le r_l\le\frac1{72}$ otherwise; also $q^*(s_l\bar\alpha e_1^*)=
|s_l|\le\frac18$.  Hence $n_l\in[\frac78,\frac98]$ in the first two cases and
$n_l\in[\frac34,\frac54]$ in the last.  (T-a): $q^*(Te_{n,m})=c_{\ell(n,m)}\le1
=c_1$.

(P1): $|\ip{U^*h^{(l_0)}}{e}|\le\sum_{s\in K_0}2^{-s}\nu_s\le\|h^{(l_0)}\|_1/2$, so
$\kappa_u>0$, and $(-\kappa_u e_1^*+h^{(l_0)})(\zh)=-\kappa_u(1+\nu_1)+\|h^{(l_0)}\|_1+
\ip{U^*h^{(l_0)}}{e}=0$ because $z=1$ on $K_0$.

(P2), (P3): every $x\in X_0$ satisfies $x_1=1/\bar\alpha$, $x=Ue$ on
$\bigcup_{l\ne l_0}S_l$ and $|x_s|\le\|U\|$ there.  Hence
$(e_1^*/q^*(e_1^*))(x)=1$ and $|\delta_lh^{(l)}(x)|\le\delta_l\|U\|$; this gives
$|u_{1,m}(x)|\ge(1-\frac18)/\frac98=\frac79$ and, for exceptional $(k,m)$,
$\ge(1-\frac1{16})/\frac{17}{16}>\frac12\ge2\varpi_{k,m}$.  For the remaining
$(k,m)$ and $n\notin\Lambda_i$, $x^{(n)}-\zh=-1_{K_0\cap(n,\infty)}$ is
annihilated by $y$, by $e_1^*$ and by $h^{(l)}$, so the choice of $s_l$ gives
$|(y+s_l\bar\alpha e_1^*)(x)|\ge3r_l$ for all $x\in X_0$, and
$|u_{k,m}(x)|\ge(3r_l-r_l)/\frac54=\frac85r_l>r_l\ge\varpi_{k,m}$.

(T-d): fix $m$ and $i$.  Infinitely many $n$ have $i(n,m)=i$, and along
them $\|u_{n,m}-y^{(i)}\|_1\to0$, because $r_l,\delta_l\to0$,
$|s_l|\le3r_l(2|\Lambda_i|+3)\to0$ and $n_l\to1$ ($l=\ell(n,m)$).  Hence
every $y^{(i)}$ is a limit point of $(u_{n,m})_n$.

(T-b), (T-c): let $x\in\ell_1(\N\times\N)$, indexed by $l$, with
$\mathcal Z:=\sum_lx_lc_lu_l\in c_{00}$, and suppose $x_{l'}\ne0$.  Write
$u_l=(y_l+\delta_lh^{(l)})/n_l$ with $y_l\in c_{00}$; the $y_l$ of the first
two items and the multiples of $e_1^*$ live on $\{1\}$, and
$1\notin\bigcup S_l$.  For $s\in S_{l'}$ only $h^{(l')}$ among the $h^{(l)}$ is
nonzero at $s$, and $y_{l'}(s)=0$ (allowedness).  If $s\in\supp y_l$ for
some $l$, let $L(s)$ be the least such $l$; allowedness at $L(s)$ gives
$2c_{L(s)}\le2^{-2s}c_{l'}\delta_{l'}$, and with $|y_l(s)|\le1$,
$n_l\ge\frac34$,
\[
 |\mathcal Z(s)|\ge\frac{|x_{l'}|c_{l'}\delta_{l'}2^{-s}}{n_{l'}}-\frac43
 \|x\|_\infty\sum_{l\ge L(s)}c_l\ge c_{l'}\delta_{l'}2^{-s}
 \Big(\frac{|x_{l'}|}{n_{l'}}-\frac43\|x\|_\infty2^{-s}\Big)>0
\]
for all large $s\in S_{l'}$.  As $S_{l'}$ is infinite, $\mathcal Z\notin c_{00}$,
a contradiction.  Hence $x=0$: $T$ is injective and $Y\cap c_{00}=\{0\}$.
\end{proof}

The construction is flexible (any finite set of prescribed vectors can be
inserted, and the targets may be the same in all blocks, as in Mart\'in's
proof); only (T-a)--(T-d) and (P1)--(P3) are used below.

\begin{lemma}[The first row]\label{lem:firstrow}
With $T$ as in Lemma~\ref{lem:T} and any finite $I\ni1$, put
$\xi:=\zh/p^{**}(\zh)$ and $f:=a+L^*J_V(L^{**}\xi)$.  Then $f\in S_{p^*}$ is
not norm attaining, $\xi$ is its normer, $q_0=1/p^{**}(\zh)$, its forced data
are $a=\bar\alpha e_1^*$, $w=J_V(L^{**}\xi)$ and $z=e_1+1_{K_0}$, $F=\{1\}$, and
$K=K_0$ is infinite.  Moreover $(2,1)$ is a strict non-peak with $w_1(2)=0$,
every other coordinate $(k,m)$, $m\in I$, is a non-degenerate peak, and
$S(f)=\R u_{2,1}$.
\end{lemma}

\begin{proof}
$\zh\in B_{\ell_\infty}+U(B_H)=B_{q^{**}}$ and $a(\zh)=\bar\alpha\zh_1=1=q^*(a)$,
so $q^{**}(\zh)=1$ and $p^{**}(\zh)=1+\|L^{**}\zh\|$.  Then
$f(\xi)=(1+\|L^{**}\zh\|)/p^{**}(\zh)=1$ and $p^*(f)\le\max(q^*(a),\|J_V\|)
=1$ by Lemma~\ref{lem:dualball}.  Proposition~\ref{prop:forced} gives the
forced data, and $\xi\notin c_0$ shows that $f$ does not attain its norm
(a norming point in $X$ would be the unique normer).

Blocks: $\zeta_m(k)=\lambda_{k,m}q_0u_{k,m}(\zh)$, and
$\sigma_m\le\|\zeta_m\|_1\le q_0m2^{-m}$.  By Lemma~\ref{lem:threshold},
off $P_m$ we have $|\zeta_m(k)|=\sigma_m\Phi_m(k)^2|w_m(k)|/C_m\le\sigma_m
\Phi_m(k)$, i.e.\ $|u_{k,m}(\xi)|\le\sigma_m/m$, and
$|u_{k,m}(\xi)|>\vartheta_m\Phi_m(k)$ implies $k\in P_m$ with positive
margin.  By (P2), $|u_{1,m}(\xi)|\ge\frac79q_0>q_02^{-m}\ge\sigma_m/m$, so
$1\in P_m$; hence $C_m\ge\Phi_m(1)M_m$ and
$\vartheta_m\Phi_m(k)\le\sigma_m\Phi_m(k)/(m\Phi_m(1))\le q_0\varpi_{k,m}$.  By
(P3), every $(k,m)\ne(2,1)$ with $k\ge2$ is a peak with positive margin; and
$u_{1,m}(\xi)$ exceeds $q_0/2\ge\vartheta_m\Phi_m(1)$.  Finally
$u_{2,1}(\xi)=0$ by (P1), so $\zeta_1(2)=0$, $2\notin P_1$ and
$w_1(2)=0$.  Thus $Q_1=\{2\}$, $Q_m=\emptyset$ for $m\ne1$, and
$S(f)=(\ell_1(\{1\})\cap\xi^\perp)+\R y_{2,1}=\R\lambda_{2,1}u_{2,1}$
because $\xi_1\ne0$ and $w_1(2)=0$.
\end{proof}

\begin{proposition}[Two-piece mates]\label{prop:twopiece}
Let $u:=u_{2,1}$, $\lambda_0:=\lambda_{2,1}=\Phi_1(2)$, $c>0$, $v:=cu$, and for
$K_1\subseteq K_0$ put $K_2:=K_0\setminus K_1$,
$\beta:=(v1_{K_1})(\zh)$ and $g_{K_1}:=v1_{K_1}-\beta a$.  There is $c_*>0$,
depending only on $\|U\|,\nu_1,M_1,C_1,\lambda_0$, such that for $0<c\le c_*$
every $g_{K_1}$ lies in $\Cset(f)$, with the admissible decompositions
\[
 f+tg=(a+tg)+L^*w\ \ (t\ge0),\qquad
 f+tg=(a+t(g-v))+L^*(w+tD)\ \ (t\le0),
\]
where $D:=(c/\lambda_0)e_2$ in block $1$ \textup{(}so $L^*D=v$\textup{)}.
\end{proposition}

\begin{proof}
We use $\|xe+y\|\le x+\ip ey+\|y\|^2/(2x)$ for $x>0$, $y\in H$, which follows
from $(\ip ey+\|y\|^2/(2x))^2\ge0$.  Let $|t|\le1$ and $c\le c_1:=
\min(\bar\alpha/4,1/(4(1+\|U\|)))$; then $|\beta|\le\|v\|_1\|\zh\|_\infty\le
c(1+\|U\|)\le\frac14$.  Put $\nu:=\|U^*a\|=\bar\alpha\nu_1$.

\emph{Side $t\ge0$.}  $a+tg=(1-t\beta)\bar\alpha e_1^*+tv1_{K_1}$, with
$1-t\beta\ge\frac34$ and $v>0$ on $K_1$.  So
$q^*(a+tg)\le(1-t\beta)(\bar\alpha+\nu)+t(\|v1_{K_1}\|_1+\ip e{h_1})+\frac23t^2
\|h_1\|^2/\nu$ with $h_1:=U^*(v1_{K_1})$.  Here $\bar\alpha+\nu=1$ and
$\|v1_{K_1}\|_1+\ip e{h_1}=(v1_{K_1})(z+Ue)=\beta$ since $z=1$ on $K_1$.
Hence $q^*(a+tg)\le1+\frac23t^2\|h_1\|^2/\nu$, and the block part is $w$.

\emph{Side $t=-s\le0$.}  $g-v=-\beta a-v_1e_1^*-v1_{K_2}$ with
$v_1:=v(1)<0$, $|v_1|\le c$, so $a+t(g-v)=(\bar\alpha(1+s\beta)+sv_1)e_1^*+
sv1_{K_2}$ with first coefficient $\ge\bar\alpha/2$.  As before,
$q^*(a+t(g-v))\le(1+s\beta)+sv_1/\bar\alpha+s(v1_{K_2})(\zh)+s^2\|h_2\|^2/\nu$
with $h_2:=U^*(v1_{K_2})$, and the first-order bracket equals
$s\,v(\zh)=sc\,u(\zh)=0$ (P1).  In block $1$, $W:=w_1+tD$ differs from
$w_1$ only at $2$, where $|W(2)|=sc/\lambda_0\le M_1$ if $c\le M_1\lambda_0$;
the peaks keep the value $M_1$, so $\|W\|_\infty=M_1$, and
$\|D_1W\|^2=C_1^2+s^2c^2$ because $w_1(2)=0$ and $\Phi_1(2)/\lambda_0=1$.
Thus $N_1(W)\le1+s^2c^2/(2C_1)$.

\emph{Conclusion.}  For $|t|\le1$, $p^*(f+tg)\le1+t^2\max(\|U\|^2c^2/\nu,
c^2/(2C_1))\le1+\frac38t^2\le s(t)$ if moreover $c^2\|U\|^2\le3\nu/8$ and
$c^2\le\frac34C_1$.  For $|t|\ge1$, $p^*(f+tg)\le1+|t|q^*(g)\le1+|t|\cdot
2c(1+\|U\|)\le1+|t|(\sqrt2-1)\le s(t)$ if $c\le0.2/(1+\|U\|)$, since
$(s(t)-1)/|t|$ increases.  Take $c_*$ the minimum of these bounds.
\end{proof}

\begin{theorem}[Nonempty defect]\label{thm:defect}
For the admissible operator $T$ of Lemma~\ref{lem:T} and every finite
$I\ni1$, the first row $f$ of Lemma~\ref{lem:firstrow} satisfies
$\mathcal K(f)\subseteq\R u_{2,1}$.  Let $c_*$ be the constant of
Proposition~\ref{prop:twopiece} and $0<c\le c_*$.  Then the mates
$g_{K_1}=g_{K_1}(c)$ of Proposition~\ref{prop:twopiece} satisfy
$g_{K_1}\in\Cset(f)$ for every $K_1\subseteq K_0$, and if
$\emptyset\ne K_1\ne K_0$, then $g_{K_1}\notin\R u_{2,1}$, so
$g_{K_1}\in\Def(f)$.  The mates $g_{\{j\}}=c\,u_{2,1}(j)
(e_j^*-\zh_j\bar\alpha e_1^*)$, $j\in K_0$, are finitely supported, and the
defect spans an infinite-dimensional space.
\end{theorem}

\begin{proof}
By Lemma~\ref{lem:firstrow} and Proposition~\ref{prop:intrinsic}(b),
$\mathcal K(f)\subseteq\R u_{2,1}$, and $g_{K_1}\in\Cset(f)$ for $0<c\le c_*$
by Proposition~\ref{prop:twopiece}.  If $g_{K_1}=\mu u_{2,1}$, comparing
coordinates $j\in K_0$ (where $a$ vanishes and $u_{2,1}(j)>0$) gives
$c1_{K_1}(j)=\mu$ for all $j\in K_0$, which forces $K_1\in\{\emptyset,K_0\}$.
For $j\in K_0$ we have $\beta=c\,u_{2,1}(j)\zh_j$, which gives the formula
for $g_{\{j\}}$; since $K_0$ is infinite, $\{j\}\ne K_0$, so every
$g_{\{j\}}$ lies in $\Def(f)$, and these vectors are linearly independent
because their $e_j^*$-coordinates $c\,u_{2,1}(j)\ne0$ sit at distinct
$j\notin\{1\}$.
\end{proof}

\begin{remark}\label{rem:twopiece}
(a) The mechanism is an \emph{exact resonance}: the block vector $u_{2,1}$
of a strict non-peak with $w=0$ is, as a base functional, supported on
$F\cup K_0$ with the contact signs on $K_0$, and $u_{2,1}(\zh)=0$; hence the
block carrier and the contact vector represent the same functional, and
splitting the contacts between the two signs of $t$ produces switching
mates.  For $K_1\in\{\emptyset,K_0\}$, $g_{K_1}$ is a two-sided certificate.
(b) If $g$ is supported in $\{1\}\cup K_0$, $g(\xi)=0$ and
$\sup_{j\in K_0}|g_j/u_{2,1}(j)|\le c\le c_*$, then the decompositions
$f+tg=(a+t(g\pm cu_{2,1}))+L^*(w\mp t(c/\lambda_0)e_2)$ for $\pm t\ge0$
(block part in block $1$) are estimated exactly as in the proof of
Proposition~\ref{prop:twopiece}, the base parts using the contacts with
their signs; so, possibly after decreasing $c_*$ by a fixed factor,
$\Cset(f)$ contains an infinite-dimensional slab.  Conversely, every mate is
supported in $\{1\}\cup K_0$ with $g_j/u_{2,1}(j)$ bounded on $K_0$
(Corollary~\ref{cor:exampleRecovered}(b)).
(c) Since $K_0$ is infinite, $f$ is not C-tame, and
Theorem~\ref{thm:tamefibre} does not apply at $f$; the infinite contact set
is unavoidable (Proposition~\ref{prop:exact}).  But $f$ satisfies the
block-tameness condition \textup{(BT)} of Section~\ref{sec:engineered},
which allows infinite contact sets (Corollary~\ref{cor:exampleRecovered}).
\end{remark}

\subsection{Non-recovery along canonical truncations}
The defect of Theorem~\ref{thm:defect} is defined relative to the known
intrinsic classes; a priori its elements could still be recovered along
every sequence.  For the canonical truncations this is not the case.

\begin{proposition}[Non-recovery]\label{prop:nonrecovery}
Let $T$, $I$ and $f$ be as in Theorem~\ref{thm:defect}, and let
$f_n:=\nabla p(x^{(n)})$, $x^{(n)}=z^{(n)}+Ue$, be the canonical truncations
of $f$.  Then:
\begin{enumerate}
\item[(a)] $f_n\in\NA\cap S_{p^*}$ and $f_n\to f$;
\item[(b)] for large $n$ the normer $x'_n:=x^{(n)}/p(x^{(n)})$ of $f_n$ is
C-tame, with $Q_{n,1}=\{2\}$, $Q_{n,m}=\emptyset$ for $m\ne1$, and no
degenerate peaks;
\item[(c)] $\Cset(f_n)\subseteq\R y^{(n)}$, where
$y^{(n)}:=\lambda_0u_{2,1}-(\Phi_1(2)^2w^{(n)}_1(2)/C^{(n)}_1)R_1^*w^{(n)}_1
\to\lambda_0u_{2,1}$;
\item[(d)] $\Li_n\Cset(f_n)\subseteq\R u_{2,1}$.  In particular, for
$\emptyset\ne K_1\ne K_0$ and $\rho\in(0,1]$, $\rho g_{K_1}\notin
\Li_n\Cset(f_n)$: the mates $g_{K_1}$ are not recovered along the canonical
truncations.
\end{enumerate}
\end{proposition}

\begin{proof}
(a) $z^{(n)}\in c_{00}$, $\|z^{(n)}\|_\infty\le1$ and $z^{(n)}_1=1=\sgn a_1$,
so $\nabla q(x^{(n)})=a$ and $f_n\in\NA\cap S_{p^*}$ by
Proposition~\ref{prop:smooth}(c); $z^{(n)}\to z$ coordinatewise, and
Proposition~\ref{prop:approximants} gives $f_n\to f$.

(b) Let $c'_n:=q(x'_n)=1/p(x^{(n)})\to q_0>0$.  Since $|u_{k,m}(x^{(n)})|
\le q(x^{(n)})=1$, $\sigma'_{n,m}\le c'_nm2^{-m}$.  Exactly as in the proof
of Lemma~\ref{lem:firstrow}, using (P2) and (P3) at $x^{(n)}\in X_0$: $1$ is a
peak of every block, $\vartheta'_{n,m}\Phi_m(k)\le c'_n\varpi_{k,m}$, and every
$(k,m)\ne(2,1)$ is a peak with margin $\mu'_{k,m}\ge c'_nr_{\ell(k,m)}$ if
$(k,m)$ is non-exceptional with $k\ge2$, and $\mu'_{k,m}\ge c'_n/4$
otherwise.  As $r_{\ell(k,m)}=(2\Phi_m(k)/c_{\ell(1,m)})^{1/2}\ge(2\Phi_m(k))^{1/2}$
and $\Phi_m(k+1)\le\Phi_m(k)/2$, for $s<c'_n/4$ we get
$\sum_{\mu'_{k,m}<s}\Phi_m(k)\le\sum_{\Phi_m(k)<s^2/(2c'^2_n)}\Phi_m(k)\le
s^2/c'^2_n=o(s)$: (MS) holds.  Next,
$u_{2,1}(x^{(n)})=u_{2,1}(\zh)-\sum_{j\in K_0,j>n}u_{2,1}(j)\to0$, while
$\vartheta'_{n,1}\Phi_1(2)\to\vartheta_1\Phi_1(2)>0$
(Proposition~\ref{prop:continuity}); so for large $n$,
$|u_{2,1}(x'_n)|<\vartheta'_{n,1}\Phi_1(2)$ and $(2,1)$ is a strict
non-peak.  Finally $a\in c_{00}$ and the contact set of $x'_n$ is
$K_0\cap[1,n]$, which is finite.

(c) By Theorem~\ref{thm:tamefibre} at $\eta=x'_n$,
$\Cset(f_n)\subseteq S(f_n)=(\ell_1(\{1\})\cap(x'_n)^\perp)+\R y^{(n)}=\R
y^{(n)}$, since $x'_n(1)\ne0$.  By Proposition~\ref{prop:continuity},
$w^{(n)}_1(2)\to w_1(2)=0$ and $C^{(n)}_1\to C_1>0$, so $y^{(n)}\to
\lambda_0u_{2,1}$.

(d) If $\gamma_ny^{(n)}\to g$, then $(\gamma_n)$ is bounded because
$\|y^{(n)}\|\to\lambda_0\|u_{2,1}\|>0$; along a subsequence $\gamma_n\to\gamma$
and $g=\gamma\lambda_0u_{2,1}$.  The last claim follows from
Theorem~\ref{thm:defect}.
\end{proof}

Thus, for this $T$, the mates recovered along \emph{every} sequence are
contained in the certificate line $\R u_{2,1}$: universal recovery fails, and
a recovery of the switching mates must use norm-attaining approximants whose
certificate spans grow, for instance by putting base mass on the contact
windows.  This is done in Section~\ref{sec:engineered}
(Remark~\ref{rem:engineered}).  By Remark~\ref{rem:tameNA}, along
any sequence of norm-attaining approximants with C-tame normers, recovery is
a linear-algebra design problem: $\Li_n\Cset(f_n)\subseteq\Li_nS(f_n)$.

\begin{remark}[Engineered approximants]\label{rem:engineered}
Proposition~\ref{prop:approximants} shows that, outside a finite window
containing $\supp a'_n$, the base coordinates $z'_n$ of a norm-attaining
approximant are free.  Section~\ref{sec:engineered} uses this freedom:
engineered approximants put masses on finitely many contacts and truncate
$z$ far out, and transfer peaks rebalance levels between the base and the
blocks.  It is proved there that every mate of the first row $f$ of
Theorem~\ref{thm:defect} is recovered (Corollary~\ref{cor:exampleRecovered});
more generally every \textup{(BT)} point lies in $\Rec$
(Corollary~\ref{cor:BTrecovered}), and two-piece mates are recovered at
every $f$ with finite base support under the conditions of
Theorem~\ref{thm:engineered}.  Thus the defect of Theorem~\ref{thm:defect}
is a defect of the intrinsic mechanisms only, and
Proposition~\ref{prop:nonrecovery} is a statement about the canonical
truncations, not about recoverability.
\end{remark}

\subsection{Exact resonances}

\begin{proposition}[Exact resonances need infinite one-sided sets]\label{prop:exact}
Let $f\in S_{p^*}$, $\Omega\in V^*$ with $D:=L^*\Omega\ne0$, and suppose that
for a sequence $s_n\to0^+$,
\[
 q^*(a+s_nD)\le s(2s_n)\quad\text{and}\quad N_m(w_m-s_n\Omega_m)\le s(2s_n)
 \quad(m\in I).
\]
Then $\sum_{j\notin F}(|D_j|-z_jD_j)=0$: off $F$, $D$ is supported in the
contact set $K$ with $z_jD_j=|D_j|$.  Since $D\in Y\setminus\{0\}$ is not
finitely supported, $F\cup K$ is infinite.  In particular there are no such
transfers when $a\in c_{00}$ and $K$ is finite \textup{(}e.g.\ at every
norm-attaining $f$\textup{)}.
\end{proposition}

\begin{proof}
By Lemma~\ref{lem:base} (dropping nonnegative terms),
$q^*(a+sD)\ge1+sD(\zh)+s\sum_{j\notin F}(|D_j|-z_jD_j)$ for $s>0$, and
$N_m(w_m-s\Omega_m)\ge\ip{w_m-s\Omega_m}{\zeta_m}/\sigma_m=1-s\ip{\Omega_m}
{\zeta_m}/\sigma_m$.  With $s=s_n$, multiply the first inequality by $q_0$
and the $m$-th by $\sigma_m$ and add: since $q_0D(\zh)=D(\xi)=\ip{\Omega}
{L^{**}\xi}=\sum_m\ip{\Omega_m}{\zeta_m}$ and $q_0+\sum_m\sigma_m=1$,
\[
 q_0s_n\sum_{j\notin F}(|D_j|-z_jD_j)\le s(2s_n)-1\le2s_n^2 .
\]
Divide by $s_n$ and let $n\to\infty$.  Each term $|D_j|-z_jD_j$ is
nonnegative and vanishes only if $D_j=0$ or $|z_j|=1$ with $z_jD_j=|D_j|$.
\end{proof}

\begin{corollary}[Two-sided linear decompositions]\label{cor:twosided}
Let $g\in\Cset(f)$ and suppose there are $\tau>0$ and
$(B_\pm,\Omega_\pm)\in\ell_1\times V^*$ with $g=B_\pm+L^*\Omega_\pm$,
$\max(q^*(a+tB_+),\|w+t\Omega_+\|)\le s(t)$ for $0\le t\le\tau$ and
$\max(q^*(a+tB_-),\|w+t\Omega_-\|)\le s(t)$ for $-\tau\le t\le0$.  If
$F\cup K$ is finite, then $B_+=B_-$, $\Omega_+=\Omega_-$, and
$g\in\cl\Cert(f)$.
\end{corollary}

\begin{proof}
For $0<t\le\tau$ average the decompositions of $f+tg$ and $f-tg$:
$f=(a+\frac t2D)+L^*(w-\frac t2\Omega_D)$ with $D:=B_+-B_-$,
$\Omega_D:=\Omega_--\Omega_+$, $L^*\Omega_D=D$, and both component norms are
at most $s(t)$ by convexity.  If $D\ne0$, Proposition~\ref{prop:exact} with
$s_n=t_n/2$ contradicts the finiteness of $F\cup K$.  Hence $D=0$, and
$\Omega_+=\Omega_-$ by injectivity of $L^*$.  The common decomposition makes
$g$ locally split, and Theorem~\ref{thm:locsplit} applies.
\end{proof}

\begin{remark}[Resonance types; the sign rule]\label{rem:resonancetypes}
The transfer in the proof of Corollary~\ref{cor:twosided} exists for every
mate and every pair of admissible decompositions at scales $\pm t$; on the
contacts, for instance, the cheaply used parts of $B_+$ and of $B_-$ have
opposite signs, so they add up in $D=B_+-B_-$.  (A general ``matching at
scale'' statement for all one-sided resources is only available as a
sketch and is not used.)  A mate outside $\mathcal
K(f)$ therefore needs \emph{resonances}: cheap transfers with large
one-sided mass at arbitrarily small scales.  Proposition~\ref{prop:exact}
shows that exact (scale-free) resonances require infinitely many contacts
(or near-flip coordinates of $F$ when $a\notin c_{00}$);
Theorem~\ref{thm:defect} realizes one.  Approximate resonances run through
one-sided block carriers: weak peaks and near-threshold strict non-peaks
at depth comparable to the scale.  For these we only record an elementary
fact: if $u_{k,m}(\xi)\ne0$, then $\sgn w_m(k)=\sgn u_{k,m}(\xi)$ (on $P_m$
by Lemma~\ref{lem:threshold}, off $P_m$ because $w_m(k)=C_m\zeta_m(k)/
(\Phi_m(k)^2\sigma_m)$).  Consequently a one-sided usage of $(k,m)$ pushes in
the direction $-\sgn u_{k,m}(\xi)$, and two near-threshold carriers whose
vectors differ by $o(\Phi)$ in the $\xi$-direction (for instance Mart\'in's
cross-block near duplicates with small perturbations) serve the same side
of $t=0$ and cannot switch with each other.  This is a heuristic guide; no
classification of approximate resonances is proved here.  When the gaps of
the carriers are not summable, Theorem~\ref{thm:carriers} puts the
corresponding mates in $\cl\Cert(f)$; the summable case and the weak-peak
case are open (Remark~\ref{rem:summablegaps}).
\end{remark}

\section{Engineered approximants}\label{sec:engineered}

Throughout this section the block set $I$ is finite and $T$ is an arbitrary
admissible operator.  We prove that the switching mates of
Section~\ref{sec:resonance} are recovered, and much more: every first row
satisfying a block-tameness condition that allows arbitrary contact sets
lies in $\Rec$ (Corollary~\ref{cor:BTrecovered}).  The approximants are
norm attaining and are \emph{engineered}: unlike the canonical truncations
of Section~\ref{sec:second}, they put base mass on finitely many contacts.
The mechanism that makes this work is \emph{first-order rebalancing}
through the transfer peaks of Theorem~\ref{thm:transfer}: a linear
mismatch between the base part and the block parts of a decomposition is a
pure shift of levels, and transfer peaks move levels between the base and
the blocks at a relative cost that can be fixed in advance.

\subsection{Levels, excesses and rebalancing}
Let $f\in S_{p^*}$ have forced data as in Section~\ref{sec:setting}; the
results of this subsection apply both to non-attaining $f$ (normer
$\xi\in S_{p^{**}}$) and to norm-attaining $f'=\nabla p(x')$ (normer
$x'\in S_p$, forced data at $x'$, with $\zh'=x'/q(x')$).  For $A\in\ell_1$
and $W_m\in\ell_\infty$ put
\begin{gather*}
 \lin_b(A):=(A-a)(\zh),\qquad G_b(A):=q^*(A)-A(\zh)=q^*(A)-1-\lin_b(A),\\
 \lin_m(W_m):=\frac{\ip{W_m-w_m}{\zeta_m}}{\sigma_m},\\
 G_m(W_m):=N_m(W_m)-\frac{\ip{W_m}{\zeta_m}}{\sigma_m}
 =N_m(W_m)-1-\lin_m(W_m).
\end{gather*}
The \emph{first-order terms} $\lin$ are linear and the \emph{excesses} $G$
are nonnegative: $q^{**}(\zh)=1$ and $|\ip{W}{\zeta_m}|\le
N_m(W)|\zeta_m|_m$.

\begin{lemma}[Bookkeeping and exact base excess]\label{lem:bookkeeping}
\textup{(a)} If $f+h=A+\sum_mR_m^*W_m$ with $h\in X^*$, then
$q_0\lin_b(A)+\sum_m\sigma_m\lin_m(W_m)=h(\xi)$.

\textup{(b)} For $B\in\ell_1$,
$G_b(a+B)=\Exc(B)+\nu\Psi(U^*B/\nu)$, where
$\Exc(B):=\sum_j\big(|a_j+B_j|-|a_j|-z_jB_j\big)\ge0$ and $\Psi$ is the
function of Lemma~\ref{lem:base}.  Moreover $\Exc(B+B')\le\Exc(B)+
2\|B'\|_1$, and if $|B_j|\le|a_j|$ for $j\in F$, then
$\Exc(B)=\sum_{j\notin F}(|B_j|-z_jB_j)$.

\textup{(c)} For every block, $G_m(W)=\sum_{k\in P_m}|\alpha_m(k)|
\big(\|W\|_\infty-\sgn(w_m(k))W(k)\big)+\big(\|D_mW\|_2-\ip{D_mW}{D_mw_m}/C_m
\big)$, and both terms are nonnegative.
\end{lemma}

\begin{proof}
(a) $q_0\lin_b(A)=(A-a)(\xi)$ and $\sigma_m\lin_m(W_m)=\ip{W_m-w_m}
{R_m^{**}\xi}$; their sum is $(A+L^*W-a-L^*w)(\xi)=h(\xi)$.
(b) This is \eqref{eq:baseidentity} with $t=1$: on $F$, $z_j=\sgn a_j$ and
$|a_j+B_j|-|a_j|-z_jB_j=2(-\sgn(a_j)B_j-|a_j|)_+$; off $F$ the summand is
$|B_j|-z_jB_j$.  Each summand is $\ge0$ because $|z_j|\le1$, and it changes
by at most $2|B'_j|$ when $B_j$ is replaced by $B_j+B'_j$.  If $|B_j|\le|a_j|$
the summand at $j\in F$ vanishes.
(c) By Lemma~\ref{lem:threshold}, $\zeta_m/\sigma_m=\alpha_m+D_m^2w_m/C_m$
with $\|\alpha_m\|_1=1$, $\alpha_m$ supported in $P_m$ with the signs of
$w_m$; hence $\ip{W}{\zeta_m}/\sigma_m=\sum_{P_m}|\alpha_m(k)|\sgn(w_m(k))
W(k)+\ip{D_mW}{D_mw_m}/C_m$, and $\|W\|_\infty=\sum_{P_m}|\alpha_m(k)|
\|W\|_\infty$.  The second term is $\ge0$ by Cauchy--Schwarz.
\end{proof}

\begin{lemma}[Transfer vectors]\label{lem:TV}
Fix a block and drop its index.  Let $\gamma\in(0,M)$,
$L:=\{k:|w(k)|\ge\gamma\}$, let $k_*\in L$ with $w(k_*)=M$, let
$\Lambda\ge0$, and put
\[
 y^{\downarrow}:=\sgn(w)1_L+\Lambda e_{k_*},\qquad
 y^{\uparrow}:=\sgn(w)1_L-\Lambda e_{k_*}.
\]
Let $W\in\ell_\infty$ with $\|W-w\|_\infty\le\eta\le(M-\gamma)/4$.
\begin{enumerate}
\item[(a)] If $\varepsilon\ge0$, $\varepsilon(2+\Lambda)\le\gamma-\eta$ and
$\varepsilon\le(M-\gamma)/4$, then $\|W-\varepsilon y^\downarrow\|_\infty
\le\|W\|_\infty-\varepsilon$.
\item[(b)] If $\varepsilon\le0$ and $|\varepsilon|(2+\Lambda)\le\gamma-\eta$,
then $\|W-\varepsilon y^\uparrow\|_\infty\le\|W\|_\infty-\varepsilon$.
\item[(c)] If $\|D(W-w)\|_2\le C/2$, then for all $y\in\ell_\infty$ and
$\varepsilon\in\R$,
\[
 \|D(W-\varepsilon y)\|_2\le\|DW\|_2-\varepsilon\frac{\ip{Dw}{Dy}}C
 +\frac{4|\varepsilon|\,\|D(W-w)\|_2\|Dy\|_2+\varepsilon^2\|Dy\|_2^2}{C}.
\]
\end{enumerate}
\end{lemma}

\begin{proof}
We may assume $\varepsilon\ne0$; then $\eta<\gamma$.  For $k\in L$,
$|W(k)-w(k)|\le\eta<\gamma\le|w(k)|$, so $\sgn W(k)=\sgn w(k)$ and
$|W(k)|\ge\gamma-\eta$; also $\|W\|_\infty\ge W(k_*)\ge M-\eta$.

(a) For $k\in L\setminus\{k_*\}$, $|W(k)|\ge\gamma-\eta\ge\varepsilon$, so
$|W(k)-\varepsilon\sgn w(k)|=|W(k)|-\varepsilon$.  At $k_*$,
$W(k_*)-\varepsilon(1+\Lambda)\le\|W\|_\infty-\varepsilon$, and
$W(k_*)-\varepsilon(1+\Lambda)\ge-(\|W\|_\infty-\varepsilon)$ because
$W(k_*)+\|W\|_\infty\ge2(\gamma-\eta)\ge\varepsilon(2+\Lambda)$.  For $k\notin
L$, $|W(k)|<\gamma+\eta\le M-\eta-\varepsilon\le\|W\|_\infty-\varepsilon$,
since $2\eta+\varepsilon\le\frac34(M-\gamma)$.

(b) Put $e:=-\varepsilon>0$.  For $k\in L\setminus\{k_*\}$,
$|W(k)+e\sgn w(k)|=|W(k)|+e$.  At $k_*$, $W(k_*)+e(1-\Lambda)\le
\|W\|_\infty+e$, and $W(k_*)+e(1-\Lambda)+\|W\|_\infty+e\ge2(\gamma-\eta)
-e(\Lambda-2)\ge0$.  For $k\notin L$, $|W(k)|\le\|W\|_\infty+e$.

(c) Put $\mathbf a:=DW$, $\mathbf b:=Dw$ (so $\|\mathbf b\|=C$ and
$\|\mathbf a\|\ge C/2$).  From $\sqrt{A^2+x}\le A+x/(2A)$,
$\|\mathbf a-\varepsilon Dy\|\le\|\mathbf a\|-\varepsilon\ip{\mathbf a}{Dy}/
\|\mathbf a\|+\varepsilon^2\|Dy\|^2/(2\|\mathbf a\|)$, and
\[
 \Big|\frac{\ip{\mathbf a}{Dy}}{\|\mathbf a\|}-\frac{\ip{\mathbf b}{Dy}}
 {\|\mathbf b\|}\Big|\le\frac{\|\mathbf a-\mathbf b\|\|Dy\|}{\|\mathbf a\|}
 +\frac{\|Dy\|\,\big|\|\mathbf b\|-\|\mathbf a\|\big|}{\|\mathbf a\|}
 \le\frac{4\|\mathbf a-\mathbf b\|\|Dy\|}{C}.\qedhere
\]
\end{proof}

\begin{proposition}[Rebalancing]\label{prop:rebalancing}
Let $f\in S_{p^*}$, $h\in X^*$ with $h(\xi)=0$, and
$f+h=A+\sum_{m\in I}R_m^*W_m$ with $A\in\ell_1$, $W_m\in\ell_\infty$.  For
each $m\in I$ let $y_m\in\ell_\infty$ and $\varepsilon_m\in\R$ satisfy
\begin{equation}\label{eq:R1}
 \|W_m-\varepsilon_my_m\|_\infty\le\|W_m\|_\infty-\varepsilon_m,\qquad
 \|D_m(W_m-w_m)\|_2\le C_m/2 .
\end{equation}
Put $Y_m:=\ip{y_m}{\zeta_m}$, $e_{y,m}:=1+\ip{D_mw_m}{D_my_m}/C_m$, the
\emph{inefficiency}
\[
 \iota_m:=\frac{|Y_m-\sigma_me_{y,m}|}{q_0}+q^*\Big(R_m^*y_m-\frac{Y_m}{q_0}a
 \Big),
\]
and $\lambda:=\sum_m\varepsilon_mY_m/q_0$; assume $\lambda\ge-1$.  Let
$\hat G_b\ge G_b(A)$ and $\hat G_m\ge G_m(W_m)$, put
$\hat\Gamma:=q_0\hat G_b+\sum_m\sigma_m\hat G_m$, and assume that $e_{y,m}>0$
and
\[
 \varepsilon_m=\big(\lin_m(W_m)+\hat G_m-\hat\Gamma\big)/e_{y,m}\qquad(m\in I).
\]
Then $p^*(f+h)\le1+\hat\Gamma+E$, where
\begin{multline*}
 E:=\sum_m|\varepsilon_m|\iota_m+|\lambda|\big(|q^*(A)-1|+q^*(A-a)\big)\\
 +\max_m\frac{4|\varepsilon_m|\,\|D_m(W_m-w_m)\|_2\|D_my_m\|_2
 +\varepsilon_m^2\|D_my_m\|_2^2}{C_m}.
\end{multline*}
\end{proposition}

\begin{proof}
Put $A':=A+\sum_m\varepsilon_mR_m^*y_m$ and $W'_m:=W_m-\varepsilon_my_m$; then
$A'+\sum_mR_m^*W'_m=f+h$, and Lemma~\ref{lem:dualball} gives
$p^*(f+h)\le\max(q^*(A'),\max_mN_m(W'_m))$.

\emph{Blocks.}  By \eqref{eq:R1} and Lemma~\ref{lem:TV}(c),
$N_m(W'_m)\le N_m(W_m)-\varepsilon_me_{y,m}+E$, and
$N_m(W_m)=1+\lin_m(W_m)+G_m(W_m)\le1+\lin_m(W_m)+\hat G_m$; by the choice of
$\varepsilon_m$ the right side is at most $1+\hat\Gamma+E$.

\emph{Base.}  Since $(R_m^*y_m)(\xi)=Y_m$, we have
$R_m^*y_m=(Y_m/q_0)a+e_m$ with $e_m:=R_m^*y_m-(Y_m/q_0)a$, so
$A'=(1+\lambda)a+B+\sum_m\varepsilon_me_m$ with $B:=A-a$.  For
$\lambda\ge0$, $q^*((1+\lambda)a+B)\le q^*(a+B)+\lambda$; for
$-1\le\lambda<0$, $(1+\lambda)a+B=(1+\lambda)(a+B)-\lambda B$ gives
$q^*((1+\lambda)a+B)\le q^*(a+B)+\lambda+|\lambda|(|q^*(a+B)-1|+q^*(B))$.
Writing $\lambda=\sum_m\varepsilon_m\sigma_me_{y,m}/q_0+\sum_m\varepsilon_m
(Y_m-\sigma_me_{y,m})/q_0$, we get
\[
 q^*(A')\le q^*(A)+\sum_m\frac{\varepsilon_m\sigma_me_{y,m}}{q_0}+E .
\]
By the choice of $\varepsilon_m$ and Lemma~\ref{lem:bookkeeping}(a) (with
$h(\xi)=0$),
\[
 \sum_m\frac{\varepsilon_m\sigma_me_{y,m}}{q_0}
 =\frac{-q_0\lin_b(A)+(\hat\Gamma-q_0\hat G_b)-(1-q_0)\hat\Gamma}{q_0}
 =-\lin_b(A)-\hat G_b+\hat\Gamma,
\]
using $\sum_m\sigma_m=1-q_0$.  Since $q^*(A)\le1+\lin_b(A)+\hat G_b$, we
obtain $q^*(A')\le1+\hat\Gamma+E$.
\end{proof}

Thus, up to the error $E$, every piece of the rebalanced decomposition sits
at the \emph{mass-weighted average} $\hat\Gamma$ of the estimated excesses,
whatever the first-order terms $\lin_m$ are; the price of moving a level
$\varepsilon_m$ is $|\varepsilon_m|\iota_m$.  The next lemma provides
transfer vectors of arbitrarily small inefficiency; it is
Step~2 of the proof of Theorem~\ref{thm:transfer}, for both signs.

\begin{lemma}[Transfer data]\label{lem:transferdata}
Let $f\in S_{p^*}$, $m\in I$ and $\gamma\in(0,M_m)$ with $|w_m(k)|\ne\gamma$
for all $k$; put $L:=\{k:|w_m(k)|\ge\gamma\}$.  There is $\Lambda_0<\infty$,
depending only on $f,m,\gamma$, such that for every $\eta_1\in(0,\frac12]$
and $k_0\in\N$ there are, for $\varsigma\in\{\downarrow,\uparrow\}$,
indices $k_*^\varsigma\ge k_0$ and numbers $\Lambda^\varsigma\ge0$ with the
following properties, where $y^\varsigma$ is the vector of
Lemma~\ref{lem:TV} built with $(\gamma,k_*^\varsigma,\Lambda^\varsigma)$ and
$Y^\varsigma,e^\varsigma_y,\iota^\varsigma$ are as in
Proposition~\ref{prop:rebalancing}:
\begin{enumerate}
\item[(a)] $k_*^\varsigma\in P_m$, $w_m(k_*^\varsigma)=M_m$, and its margin
$\mu^\varsigma_*:=\mu_{k_*^\varsigma,m}$ is positive;
\item[(b)] $Y^\varsigma=\sigma_me^\varsigma_y\pm\Lambda^\varsigma
\lambda_{k_*^\varsigma,m}\mu^\varsigma_*$ \textup{(}$+$ for $\downarrow$, $-$
for $\uparrow$\textup{)}, and $\Lambda^\varsigma\lambda_{k^\varsigma_*,m}\le
\Lambda_0$;
\item[(c)] $\iota^\varsigma\le\eta_1$;
\item[(d)] $|e^\varsigma_y-e^0_y|\le\eta_1$, where
$e^0_y:=1+\sum_{k\in L}\Phi_m(k)^2|w_m(k)|/C_m\ge1$; and
$\|D_my^\varsigma\|_2\le\|\Phi_m\|_2+\Lambda_0/m$.
\end{enumerate}
\end{lemma}

\begin{proof}
Drop $m$.  Put $\hat\pi:=R^*(\sgn(w)1_L)\in\ell_1$ and
$\tau_*:=(\hat\pi(\xi)/q_0)a-\hat\pi$, so $\tau_*(\xi)=0$, and
$\Lambda_0:=2(q^*(\tau_*)+1)$.  For $\delta_0\in(0,1]$ put
$\mathcal T^\downarrow:=\tau_*+\delta_0(q^*(\tau_*)+1)a$ and $\mathcal
T^\uparrow:=-\tau_*+\delta_0(q^*(\tau_*)+1)a$.  Then $\mathcal
T^\varsigma(\xi)=\delta_0(q^*(\tau_*)+1)q_0>0$ and $\ell^\varsigma:=q^*(\mathcal
T^\varsigma)\in(0,\Lambda_0]$.  Recall that $|v(\xi)|=q_0|v(\zh)|\le q_0q^*(v)$
for $v\in\ell_1$, because $q^{**}(\zh)=1$.

Fix $\varsigma$ and drop it.  Let $\delta_1\in(0,\mathcal
T(\xi)/(2\ell q_0))$.  By (T-d) the tail $\{u_k:k\ge k_0\}$ is dense in
$S_{q^*}$, and $\Phi(k)\to0$; choose $k_*\ge k_0$ with $q^*(u_{k_*}-\mathcal
T/\ell)\le\delta_1$, $\vartheta\Phi(k_*)<\mathcal T(\xi)/(2\ell)$ and
$\ell\Phi(k_*)M/(mC)\le\eta_1$.  Then $u_{k_*}(\xi)\ge\mathcal
T(\xi)/\ell-\delta_1q_0\ge\mathcal T(\xi)/(2\ell)>\vartheta\Phi(k_*)$, so
by Lemma~\ref{lem:threshold} $k_*$ is a peak with $w(k_*)=M$ and positive
margin $\mu_*$: (a).  Put $\Lambda:=\ell/\lambda_{k_*}$.

(b) By Lemma~\ref{lem:threshold}, $\ip{y}{\zeta}=\sigma(\ip y\alpha+
\ip{Dy}{Dw}/C)$.  As $P\subseteq L$ and $\alpha$ lives on $P$ with the signs
of $w$, $\ip{\sgn(w)1_L}{\alpha}=\|\alpha\|_1=1$, and $\alpha(k_*)\ge0$; so
$Y=\sigma e_y\pm\sigma\Lambda\alpha(k_*)=\sigma e_y\pm\Lambda\lambda_{k_*}\mu_*$
by \eqref{eq:margin}.  Also $\Lambda\lambda_{k_*}=\ell\le\Lambda_0$.

(c) For $\downarrow$, $R^*y=\hat\pi+\ell u_{k_*}=\hat\pi+\mathcal T+\ell
(u_{k_*}-\mathcal T/\ell)=\big(\hat\pi(\xi)/q_0+\delta_0(q^*(\tau_*)+1)\big)a+
\ell(u_{k_*}-\mathcal T/\ell)$; for $\uparrow$, $R^*y=\hat\pi-\ell u_{k_*}$
and the same computation applies with $-\tau_*$.  In both cases
$R^*y-(Y/q_0)a=\varrho-(\varrho(\xi)/q_0)a$ with
$\varrho:=\pm\ell(u_{k_*}-\mathcal T/\ell)$, whose $q^*$-norm is at most
$2q^*(\varrho)\le2\ell\delta_1$.  Moreover $\ell\mu_*\le\ell u_{k_*}(\xi)\le
\mathcal T(\xi)+\ell\delta_1q_0$, so
$\iota\le\delta_0(q^*(\tau_*)+1)+3\Lambda_0\delta_1$.  Choosing first
$\delta_0$ and then $\delta_1$ small gives $\iota\le\eta_1$.

(d) $e_y-e^0_y=\pm\Lambda\Phi(k_*)^2M/C=\pm\ell\Phi(k_*)M/(mC)$, since
$\lambda_{k_*}=m\Phi(k_*)$; and $\|Dy\|_2\le\|\Phi\|_2+\Lambda\Phi(k_*)=
\|\Phi\|_2+\ell/m$.
\end{proof}

\begin{lemma}[Persistence of transfer data]\label{lem:persistence}
Let $f_n\to f$ in $S_{p^*}$, and let $(\gamma,k_*^\varsigma,\Lambda^\varsigma)$
be transfer data of $f$ in block $m$ as in Lemma~\ref{lem:transferdata}.  Put
$L_n:=\{k:|w_{n,m}(k)|\ge\gamma\}$ and $y_n^\varsigma:=\sgn(w_{n,m})1_{L_n}
\pm\Lambda^\varsigma e_{k^\varsigma_*}$.  Then for $n$ large: $k_*^\varsigma$
is a peak of $w_{n,m}$ with $w_{n,m}(k_*^\varsigma)=M_{n,m}$ and positive
margin, and $P_{n,m}\subseteq L_n$; and, computing the quantities of
Proposition~\ref{prop:rebalancing} at $f_n$,
$Y^\varsigma_n/q^{(n)}_0\to Y^\varsigma/q_0$, $e^\varsigma_{y,n}\to
e^\varsigma_y$, $\iota^\varsigma_n\to\iota^\varsigma$, and
$D_my^\varsigma_n\to D_my^\varsigma$ in $\ell_2$.
\end{lemma}

\begin{proof}
By Proposition~\ref{prop:continuity}, $\xi_n\to\xi$ weak$^*$, $w_{n,m}\to
w_m$ coordinatewise, $D_mw_{n,m}\to D_mw_m$, and the scalar data and $a_n$
converge.  Since $u_{k_*,m}\in\ell_1$, $u_{k_*,m}(\xi_n)\to u_{k_*,m}(\xi)$,
and $\vartheta_{n,m}\to\vartheta_m$; so the strict threshold inequality of
Lemma~\ref{lem:threshold} at $k_*$ persists, with the positive sign.  Peaks
of $w_{n,m}$ have modulus $M_{n,m}\to M_m>\gamma$.  As $|w_m(k)|\ne\gamma$,
$y^\varsigma_n(k)\to y^\varsigma(k)$ for every $k$, with $|y_n^\varsigma(k)|\le1$
off $k_*$; dominated convergence with the summable majorants
$\lambda_{k,m}\|u_{k,m}\|_1$ and $\Phi_m(k)$ gives
$R_m^*y_n^\varsigma\to R_m^*y^\varsigma$ in $\ell_1$ and
$D_my^\varsigma_n\to D_my^\varsigma$ in $\ell_2$.  Hence
$Y_n=(R_m^*y_n)(\xi_n)\to Y$, and all the listed quantities converge.
\end{proof}

\begin{lemma}[Assembly]\label{lem:assembly}
Let $f,f'\in S_{p^*}$, $g\in\Cset(f)$, $\rho\in(0,1)$, $g'\in X^*$,
$0<\delta\le1$ and $0<T_0\le1$ with $T_0^2\le4\delta$.  If
$p^*(f'+\tau g')\le1+\frac{\tau^2}2(1-\delta)$ for $|\tau|\le T_0$,
$p^*(f'-f)\le(1-\rho^2)T_0^2/6$ and $p^*(g'-\rho g)\le(1-\rho^2)T_0/6$, then
$g'\in\Cset(f')$.  If moreover $f'$ attains its norm, then $(f',g')$
attains its norm and $\|(f',g')-(f,\rho g)\|\le p^*(f'-f)+p^*(g'-\rho g)$.
\end{lemma}

\begin{proof}
For $|\tau|\le T_0$: $\frac{\tau^2}2\delta\ge\frac{\tau^4}8$, so the bound is
at most $1+\frac{\tau^2}2-\frac{\tau^4}8\le s(\tau)$.  For $|\tau|\ge T_0$,
$p^*(f'+\tau g')\le s(\rho\tau)+(1-\rho^2)(T_0^2+|\tau|T_0)/6$, and the last
term is at most $(1-\rho^2)\min(\tau^2,|\tau|)/3$; conclude with
Lemma~\ref{lem:slack}(a).  The last assertions are
Proposition~\ref{prop:reduction}(d) and $\|(f_1,g_1)\|\le p^*(f_1)+p^*(g_1)$.
\end{proof}

\subsection{Two-piece data and the exact one-sided coefficient}

\begin{definition}[Side-admissible data, two-piece data]\label{def:twopiece}
Let $f\in S_{p^*}$ with $F$ finite, and $g\in X^*$.  A pair $(b,\omega)$
with $b\in\ell_1$ and $\omega=(\omega_m)_{m\in I}$, $\omega_m\in c_{00}$,
$\supp\omega_m\subseteq Q_m$, \emph{represents} $g$ if
$g=b+\sum_mR_m^*(\omega_m-d_m(\omega_m)w_m)$, where
$d_m(\omega):=\ip{D_mw_m}{D_m\omega}/C_m$.  It is \emph{side-$+$ admissible}
\textup{(}resp.\ \emph{side-$-$ admissible}\textup{)} if moreover $b_j=0$
for $j\notin F\cup K$ and $z_jb_j\ge0$ \textup{(}resp.\ $z_jb_j\le0$\textup{)}
for $j\in K$.  \emph{Two-piece data} for $g$ are a side-$+$ admissible pair
$(b^+,\omega^+)$ and a side-$-$ admissible pair $(b^-,\omega^-)$; put
$\Delta d_m:=d_m(\omega^-_m)-d_m(\omega^+_m)$,
\[
 \Gamma_w(b,\omega):=q_0h(b)+\sum_m\sigma_mH_m(\omega_m),\qquad
 \kappa_w:=\max\big(\Gamma_w(b^+,\omega^+),\Gamma_w(b^-,\omega^-)\big).
\]
The data are \emph{$d$-neutral} if $\Delta d_m=0$ for all $m$.  For $g$ with
$g(\xi)=0$ the \emph{one-sided coefficients} are
$\gamma^\pm(g):=\inf\{\Gamma_w(b,\omega):(b,\omega)\text{ side-$\pm$
admissible for }g\}\in[0,\infty]$.
\end{definition}

If $(b,\omega)$ represents $g$, then $b(\xi)=g(\xi)$, because
$\ip{\omega_m-d_m(\omega_m)w_m}{\zeta_m}=0$ (Lemma~\ref{lem:algebra}).  A
pair with $b$ supported in $F$ is admissible on both sides; it is then a
balanced finite certificate in the sense of Definition~\ref{def:tame}, and
$\Gamma_w(b,\omega)=\Gamma_w(g)$.  The mates of
Proposition~\ref{prop:twopiece} carry the two-piece data
$(b^+,\omega^+)=(g_{K_1},0)$ and $(b^-,\omega^-)=(g_{K_1}-cu_{2,1},
(c/\lambda_0)e_2\text{ in block }1)$, which are $d$-neutral because
$w_1(2)=0$.  Since $h$ and $H_m$ are positive semidefinite quadratic forms,
$\Gamma_w$ is convex in $(b,\omega)$.

\begin{definition}[Block-tame points]\label{def:BT}
$f\in S_{p^*}$ satisfies \textup{(BT)} if $F$ is finite, every $Q_m$ is
finite, there are no degenerate peaks, and \textup{(MS)} holds at $\xi$.
Nothing is assumed about the contact set $K$, which may be infinite.
\end{definition}

Thus the points of Theorem~\ref{thm:tame} are the \textup{(BT)} points with
$K$ finite, and the first row $f$ of Lemma~\ref{lem:firstrow} is a
\textup{(BT)} point with $K=K_0$ infinite
(Corollary~\ref{cor:exampleRecovered} below).

\begin{proposition}[One-sided upper bound]\label{prop:onesidedupper}
Let $f\in S_{p^*}$ with $F$ finite, $g\in X^*$ with $g(\xi)=0$, and let
$(b,\omega)$ be side-$+$ admissible for $g$.  Then
\[
 \limsup_{t\to0^+}\frac{2(p^*(f+tg)-1)}{t^2}\le\Gamma_w(b,\omega),
\]
and the same holds for $t\to0^-$ and side-$-$ admissible pairs.  In
particular, for every balanced finite certificate $g\in S(f)$ with
$a\in c_{00}$, $\limsup_{t\to0}2(p^*(f+tg)-1)/t^2\le\Gamma_w(g)$.
\end{proposition}

\begin{proof}
Let $t>0$; the case $t<0$ is the same with $-g$, $-b$ and $-\omega$, which
turns side $-$ into side $+$.  We have the exact decomposition
$f+tg=A+\sum_mR_m^*W_m$, $A:=a+tb$, $W_m:=w_m+t(\omega_m-d_mw_m)$,
$d_m:=d_m(\omega_m)$.

\emph{First-order terms and excesses.}  $\lin_b(A)=tb(\zh)=0$, since
$b(\xi)=g(\xi)=0$, and $\lin_m(W_m)=0$ by Lemma~\ref{lem:algebra}.  For
$t\|b\|_\infty\le\min_F|a_j|$ there are no flips on $F$; off $F$, $b$
vanishes on $J$ and $z_jb_j=|b_j|$ on $K$; so $\Exc(tb)=0$ by
Lemma~\ref{lem:bookkeeping}(b), and by Lemma~\ref{lem:base},
$G_b(A)=\nu\Psi(tU^*b/\nu)\le\frac{t^2}2h(b)(1+2t\|U^*b\|/\nu)$ for
$t\|U^*b\|\le\nu/2$.  By Lemma~\ref{lem:block}(d), for $t$ below the
coordinatewise radius of $\omega_m$,
$G_m(W_m)=N_m(W_m)-1\le\frac{t^2}2H_m(\omega_m)(1+2t|d_m|M_m/C_m)$.  Hence
there is $K_1$ with $G_b(A)\le\hat G_b:=\frac{t^2}2h(b)(1+K_1t)$ and
$G_m(W_m)\le\hat G_m:=\frac{t^2}2H_m(\omega_m)(1+K_1t)$ for small $t$, and
$\hat\Gamma=\frac{t^2}2\Gamma_w(b,\omega)(1+K_1t)$.

\emph{Rebalancing.}  Fix $\eta_1\in(0,\frac14]$.  In every block choose
$\gamma_m\in(0,M_m)$ with $|w_m(k)|\ne\gamma_m$ for all $k$ and transfer
data as in Lemma~\ref{lem:transferdata}; then $e^\varsigma_{y,m}\ge
\frac12$, $\iota^\varsigma_m\le\eta_1$, and $\|D_my^\varsigma_m\|_2\le K_2$,
$|Y^\varsigma_m|/q_0\le K_2$ with $K_2$ independent of $\eta_1$.  Put
$\varepsilon_m:=(\hat G_m-\hat\Gamma)/e_{y,m}$, where $y_m:=y^\downarrow_m$
if $\hat G_m\ge\hat\Gamma$ and $y_m:=y^\uparrow_m$ otherwise.  Then
$|\varepsilon_m|\le K_3t^2$ with $K_3$ independent of $\eta_1$.  Since
$\|W_m-w_m\|_\infty+\|D_m(W_m-w_m)\|_2\le K_4t$, the hypotheses of
Lemma~\ref{lem:TV}(a),(b),(c) hold for small $t$ (the radii depend on
$\Lambda^\varsigma_m$, which is fixed once $\eta_1$ is), so \eqref{eq:R1}
holds and $|\lambda|\le|I|K_2K_3t^2\le1$.  Proposition~\ref{prop:rebalancing}
gives
\[
 p^*(f+tg)\le1+\frac{t^2}2\Gamma_w(b,\omega)(1+K_1t)+|I|K_3\eta_1t^2+O(t^3),
\]
because $|q^*(A)-1|+q^*(A-a)=O(t)$.  Hence the $\limsup$ is at most
$\Gamma_w(b,\omega)+2|I|K_3\eta_1$, and $\eta_1$ is arbitrary.  The last
sentence follows because such a $g$ has a pair $(b,\omega)$ with $\supp
b\subseteq F$.
\end{proof}

The last sentence proves the statement left as a sketch in
Remark~\ref{rem:transfer}(c).

\begin{theorem}[Exact one-sided coefficient]\label{thm:onesided}
Let $f$ satisfy \textup{(BT)} and let $g\in X^*$ with $g(\xi)=0$.  Then:
\begin{enumerate}
\item[(a)] $\displaystyle\lim_{t\to0^\pm}\frac{2(p^*(f+tg)-1)}{t^2}=
\gamma^\pm(g)$ in $[0,\infty]$;
\item[(b)] if $\gamma^\pm(g)<\infty$, the infimum defining it is attained;
\item[(c)] every $g\in\Cset(f)$ satisfies $\gamma^\pm(g)\le1$, and therefore
has two-piece data with $\kappa_w=\max(\gamma^+(g),\gamma^-(g))\le1$.
\end{enumerate}
\end{theorem}

\begin{proof}
By Proposition~\ref{prop:onesidedupper}, $\limsup_{t\to0^+}\le\gamma^+(g)$.
We prove $\liminf_{t\to0^+}\ge\gamma^+(g)$ and (b) for side $+$; side $-$
is symmetric (replace $t$ by $-t$ and the cone $C_+$ below by
$C_-:=-C_+$).

\emph{Step 1: test points.}  Let
$C_+:=\{\chi_c\in c_{00}:\chi_c=0\text{ on }F,\ z_j\chi_{c,j}\le0\text{ for }
j\in K\}$, a convex cone (inward moves of contacts, arbitrary moves of the
free coordinates).  Let $k\in H$ with $\ip ke=0$ and $\chi_c\in C_+$, and put
$\chi:=Uk+\chi_c\in c_0$.  As in the base test of the proof of
Proposition~\ref{prop:lowerbound} (with $k$ in place of $k_\perp$, which is
no longer restricted to $E_F$), let $\kappa_2:=\nu\|k\|^2/(2q_0)$, choose
$k_2\in H$ with $(Uk_2)_j=-\kappa_2z_j$ for $j\in F$ and
$\ip{e}{k_2}=\kappa_2-\|k\|^2/(2q_0)$ (possible by the Gram system there,
which only involves $F$), put $\chi_2:=(Uk_2)1_{\N\setminus F}$,
$\eta_t:=\xi+t\chi+t^2\chi_2$, $\mathbf h_t:=q_0e+tk+t^2k_2$ and
$Z_t:=\eta_t-U\mathbf h_t$.  Then $Z_{t,j}=(q_0+t^2\kappa_2)z_j$ for $j\in F$
and $Z_{t,j}=q_0z_j+t\chi_{c,j}$ for $j\notin F$: the vector $k$ does not
move any coordinate of $Z_t$.  Since $\chi_c$ is finitely supported, moves
contacts inward and $|z_j|<1$ on its support off $F\cup K$, we have
$\|Z_t\|_\infty\le q_0+t^2\kappa_2$ for small $t$; and $\|\mathbf h_t\|=
q_0+t^2\kappa_2+O(t^3)$ because $k\perp e$.  As
$B_{q^{**}}=B_{\ell_\infty}+U(B_H)$, $q^{**}(\eta_t)\le q_0+t^2\kappa_2+
O(t^3)$, while $a(\eta_t-\xi)=t\ip{U^*a}{k}=0$.  Hence the base excess
satisfies $E_q(t\chi+t^2\chi_2)\le\frac{t^2}2\frac\nu{q_0}\|k\|^2+O(t^3)$.

\emph{Step 2: the test inequality.}  For each block let
$\mathfrak d_m(\delta):=(S_{P_m}(\delta),(\delta_k)_{k\in Q_m})\in
\R\times\R^{Q_m}$, where $S_{P_m}(\delta)=\sum_{k\in P_m}\sgn(w_m(k))\delta_k$.
The block test of the proof of Proposition~\ref{prop:lowerbound} (which uses
only that $Q_m$ is finite, that there are no degenerate peaks, and
\textup{(MS)}) gives $E_m(tR_m\chi+t^2R_m\chi_2)\le\frac{t^2}2\mathcal
Q_m(\mathfrak d_m(R_m\chi))+o(t^2)$, where $\mathcal Q_m$ is the Hessian
of the function $G$ there, a quadratic form on $\R\times\R^{Q_m}$.  The
formulas of the block suprema in that proof show that $\mathcal Q_m$ is
positive semidefinite and vanishes at $\mathfrak d_m(\zeta_m)$, and that
for every $\omega\in\R^{Q_m}$
\begin{equation}\label{eq:blocksup}
 \sup_{\mathfrak d\in\R\times\R^{Q_m}}\big(2v(\mathfrak d)-\mathcal Q_m(
 \mathfrak d)\big)=\sigma_mH_m(\omega),\qquad v:=\omega-d_m(\omega)w_m,
\end{equation}
where a vector $V\in\ell_\infty$ which is a constant multiple of
$\sgn(w_m)$ on $P_m$ is regarded as the linear form
$\mathfrak d_m(\delta)\mapsto\ip V\delta$ on the data (recall
$P_m\cup Q_m=\N$, and every $\mathfrak d\in\R\times\R^{Q_m}$ is
$\mathfrak d_m(\delta)$ for some $\delta\in\ell_1$).  Every linear form
on $\R\times\R^{Q_m}$ is $\mathfrak d\mapsto\ip V\delta$ with
$V=c\,\sgn(w_m)1_{P_m}+\tilde\omega$, $\tilde\omega\in\R^{Q_m}$, and then
$V=\omega+(c/M_m)w_m$ with $\omega:=\tilde\omega-(c/M_m)w_m1_{Q_m}$; since
$\ip{\omega}{\zeta_m}=\sigma_md_m(\omega)$, $V$ is of the form $2v$ in
\eqref{eq:blocksup} iff $\ip V{\zeta_m}=0$.  If $\ip V{\zeta_m}\ne0$, the
supremum of $\ip V\delta-\mathcal Q_m$ is $+\infty$ (take
$\delta=s\zeta_m$).  As in \eqref{eq:testratio}, $p^*(f+tg)\ge
(f+tg)(\eta_t)/p^{**}(\eta_t)=1+t^2g(\chi)-\Delta(\eta_t)+O(t^3)$, so
\[
 \liminf_{t\to0^+}\frac{2(p^*(f+tg)-1)}{t^2}\ge\Phi(k,\chi_c):=2g(\chi)
 -\frac\nu{q_0}\|k\|^2-\sum_m\mathcal Q_m(\mathfrak d_m(R_m\chi)).
\]

\emph{Step 3: weak duality.}  Let $(b,\omega)$ be side-$+$ admissible for
$g$.  Then $b(\chi_c)\le0$, $2b(Uk)=2\ip{P^\perp U^*b}{k}\le q_0h(b)+
\frac\nu{q_0}\|k\|^2$, and by \eqref{eq:blocksup}
$2(\omega_m-d_mw_m)(R_m\chi)-\mathcal Q_m(\mathfrak d_m(R_m\chi))\le
\sigma_mH_m(\omega_m)$.  Adding, $\Phi(k,\chi_c)\le\Gamma_w(b,\omega)$.

\emph{Step 4: strong duality.}  Let $n:=\sum_m(|Q_m|+1)$,
$\mathfrak d(\chi):=(\mathfrak d_m(R_m\chi))_m\in\R^n$,
$\mathcal Q(y):=\sum_m\mathcal Q_m(y_m)$ (convex and finite on $\R^n$), and
\[
 \Psi(y):=\sup\Big\{2g(Uk+\chi_c)-\frac\nu{q_0}\|k\|^2:\ k\perp e,\ \chi_c\in
 C_+,\ \mathfrak d(Uk+\chi_c)=y\Big\}
\]
($\sup\emptyset:=-\infty$).
$\Psi$ is concave (the constraint set is convex and the objective jointly
concave), $\Psi(0)\ge0$, and $\sup\Phi=\sup_y(\Psi(y)-\mathcal Q(y))$.  If
$\Psi(y_0)=+\infty$ for some $y_0$, then $\sup\Phi=+\infty$, so by Step~3 no
side-$+$ admissible pair exists, $\gamma^+(g)=+\infty$, and Step~2 gives
the claim.  Otherwise $\Psi$ is proper, $\operatorname{ri}(\operatorname{dom}
\mathcal Q)=\R^n$ meets $\operatorname{ri}(\operatorname{dom}\Psi)$, and
Fenchel's duality theorem \cite[Theorem~31.1]{Rockafellar} gives
\[
 \sup_y\big(\Psi(y)-\mathcal Q(y)\big)=\min_{\lambda\in\R^n}\big(\mathcal
 Q^*(\lambda)-\Psi_*(\lambda)\big),
\]
the minimum being attained (both sides may be $+\infty$), where $\mathcal
Q^*(\lambda)=\sup_y(\ip\lambda y-\mathcal Q(y))$ and $\Psi_*(\lambda)=
\inf_y(\ip\lambda y-\Psi(y))$.  Write $\ip{\lambda}{\mathfrak d(\chi)}=
\sum_m\ip{V_m}{R_m\chi}=\varphi_\lambda(\chi)$ with $V_m$ as in Step~2 and
$\varphi_\lambda:=\sum_mR_m^*V_m\in\ell_1$.  Then, with
$\beta:=g-\varphi_\lambda/2$,
\[
 -\Psi_*(\lambda)=\sup_{k\perp e}\big(2\ip{P^\perp U^*\beta}{k}-\tfrac\nu{q_0}
 \|k\|^2\big)+\sup_{\chi_c\in C_+}2\beta(\chi_c),
\]
the first supremum is $(q_0/\nu)\|P^\perp U^*\beta\|^2=q_0h(\beta)$, and the
second is $0$ if $\beta_j=0$ on $J$ and $z_j\beta_j\ge0$ on $K$, and
$+\infty$ otherwise.  By Step~2, $\mathcal Q^*(\lambda)=\sum_m\sigma_m
H_m(\omega_m)$ if $V_m=2(\omega_m-d_m(\omega_m)w_m)$ with $\omega_m\in
\R^{Q_m}$ for every $m$, and $+\infty$ otherwise.  In the finite case
$\varphi_\lambda/2=\sum_mR_m^*(\omega_m-d_m(\omega_m)w_m)$, so
$\mathcal Q^*(\lambda)-\Psi_*(\lambda)$ is finite iff $(\beta,\omega)$ is a
side-$+$ admissible pair for $g$, and then equals $\Gamma_w(\beta,\omega)$.
Conversely every side-$+$ admissible pair arises from some $\lambda$.
Hence $\sup\Phi=\gamma^+(g)$, and a minimizing $\lambda$ yields a pair
attaining $\gamma^+(g)$ when it is finite.  With Step~2 this proves (a)
and (b).

(c) If $g\in\Cset(f)$, then $g(\xi)=0$ and $p^*(f+tg)\le s(t)\le1+t^2/2$, so
both limits in (a) are at most $1$; the minimizers of (b) are two-piece data
with $\kappa_w\le1$.
\end{proof}

\begin{remark}\label{rem:onesided}
(a) If $K$ is finite, two pairs representing $g$ whose base parts are
supported in $F\cup K$ coincide (Lemma~\ref{lem:rigidity}(a)); a pair that
is admissible on both sides has $b=0$ on $K$.  Hence
$\max(\gamma^+(g),\gamma^-(g))<\infty$ iff $g\in S(f)$, and then
$\gamma^+(g)=\gamma^-(g)=\Gamma_w(g)$; so at \textup{(BT)} points with $K$ finite,
Theorem~\ref{thm:onesided} contains Proposition~\ref{prop:lowerbound} and
gives again $\Cset(f)\subseteq S(f)$.  The
proof above needs no decoupling, so the cofiniteness of $J$
(Remark~\ref{rem:kinks}(a)) is not needed.
(b) When $K$ is infinite, $\gamma^\pm(g)$ can be smaller than $\Gamma_w$ of
a balanced finite certificate representing $g$: the optimal side
decompositions re-split through contacts with the cheap sign.  This is the
phenomenon observed numerically in finite models in
Remark~\ref{rem:kinks}(a) (numerical evidence only).
\end{remark}

\subsection{The engineered approximants}
Throughout this subsection $f\in S_{p^*}$ has $F$ finite, $g\in\Cset(f)$
carries two-piece data $(b^\pm,\omega^\pm)$, and $\rho\in(0,1)$.  Put
\[
 b^\theta:=\tfrac12(b^++b^-),\quad \omega^\theta:=\tfrac12(\omega^++\omega^-),
 \quad v:=b^+-b^-,\quad \Omega_m:=\supp\omega^+_m\cup\supp\omega^-_m .
\]
Then $(b^\theta,\omega^\theta)$ represents $g$, $v$ is supported in $F\cup
K$ with $z_jv_j=|v_j|$ on $K$, and $b^\pm=b^\theta\pm v/2$.

\begin{definition}[Engineered approximants]\label{def:engineered}
For a \emph{window} $N\ge\max F$, a \emph{scale} $s_1\in(0,1]$ and a
\emph{cut-off} $N''>N$ put $m_j:=4\rho s_1|b^\theta_j|$ for $j\in
K\cap[1,N]$ and
\begin{gather*}
 a'':=a+\sum_{j\in K\cap[1,N]}m_jz_je_j^*,\qquad a':=a''/q^*(a''),\qquad
 e':=U^*a'/\|U^*a'\|,\\
 z'_j:=z_j\ (j\le N''),\qquad z'_j:=0\ (j>N''),\qquad \hat x':=z'+Ue',\\
 x':=\hat x'/p(\hat x'),\qquad f':=\nabla p(x').
\end{gather*}
A \emph{construction sequence} is a sequence of parameters
$(N_i,s_{1,i},N''_i)$ with $N_i\to\infty$ and $s_{1,i}\to0$; a statement
holds \emph{at late stages} if it holds for all large $i$.  We write
$t_{k,m}:=\|u_{k,m}1_{(N'',\infty)}\|_1$ and
$\tau(N''):=\sum_{m\in I}\sum_k\lambda_{k,m}t_{k,m}$; by dominated convergence
$\tau(N'')\to0$ as $N''\to\infty$.  We shall always assume
\begin{equation}\label{eq:N1}
 \tau(N''_i)\le s_{1,i}^2 .
\end{equation}
\end{definition}

\begin{lemma}[The approximants]\label{lem:approxfacts}
\textup{(a)} $f'\in\NA((X,p),\R)\cap S_{p^*}$; its forced data are $a'$,
$e'$, $z'$, $\zh'=\hat x'$ and $q_0'=c':=q(x')=1/p(\hat x')$;
$\supp a'=F\cup\{j\in K\cap[1,N]:b^\theta_j\ne0\}$, with $\sgn a'_j=z_j$
there.  If $4\rho s_1\|b^\theta\|_1(1+\|U\|)\le1$, then $q^*(a'')\le2$, so
$|a'_j|\ge|a_j|/2$ on $F$ and $|a'_j|\ge m_j/2$ on window contacts.

\textup{(b)} Along every construction sequence $f'\to f$, and all forced data
of $f'$ converge to those of $f$ as in Proposition~\ref{prop:continuity}.

\textup{(c)} With $K_e:=8\rho\|U\|\|b^\theta\|_1/\nu$, and at late stages
for the bounds involving $|R_m\hat x'|_m$:
\begin{enumerate}
\item[(E1)] $\|a''-a\|_1\le4\rho s_1\|b^\theta\|_1$ and $\|e'-e\|\le K_es_1$;
\item[(E2)] $|u(\hat x')-u(\zh)|\le K_es_1q^*(u)+\|u1_{(N'',\infty)}\|_1$ for
$u\in\ell_1$;
\item[(E3)] $\big||R_m\hat x'|_m-|R_m^{**}\zh|_m\big|\le\sum_k\lambda_{k,m}
|u_{k,m}(\hat x'-\zh)|\le K_es_1+\tau(N'')$, and
$|1/c'-1/q_0|\le|I|(K_es_1+\tau(N''))$;
\item[(E4)] for every $\omega\in c_{00}$ vanishing on $P_m\cup P'_m$ we
have $|d'_m(\omega)-d_m(\omega)|\le K(\omega)(s_1+\tau(N''))$, where
$d'_m(\omega):=\ip{D_mw'_m}{D_m\omega}/C'_m$ and the constant $K(\omega)$
depends only on $f,g,\rho,\omega$;
\item[(E5)] the \emph{Bregman terms} $\mathrm{Bx}_m:=\ip{w'_m-w_m}{R_m\hat x'}$
satisfy
\[
 0\le\mathrm{Bx}_m\le K_es_1\sum_k\lambda_{k,m}|w'_m(k)-w_m(k)|+2\tau(N'');
\]
hence $\mathrm{Bx}_m/s_1\to0$ along every construction
sequence satisfying \eqref{eq:N1}.
\end{enumerate}
\end{lemma}

\begin{proof}
(a) $a''\in c_{00}$, $z'\in c_{00}$, $\|z'\|_\infty\le1$, and $z'=\sgn a'$ on
$\supp a'$ (on $F$ because $z=\sgn a$ there, on window contacts because
$a''_j=m_jz_j$).  By Proposition~\ref{prop:smooth}(c), $q(\hat x')=1$,
$\nabla q(\hat x')=a'$, and $f'=\nabla p(\hat x')=\nabla p(x')$ attains its
norm at $x'$, which is its unique normer; Proposition~\ref{prop:forced} at
$x'$ gives the forced data.  Finally $q^*(a'')\le1+(1+\|U\|)\|a''-a\|_1$.
(b) $a'\to a$ in $\ell_1$ and $z'\to z$ coordinatewise; apply
Proposition~\ref{prop:approximants} and Proposition~\ref{prop:continuity}.
(E1) $\|e'-e\|=\|U^*a''/\|U^*a''\|-U^*a/\|U^*a\|\|\le2\|U^*(a''-a)\|/\nu$.
(E2) $\hat x'-\zh=-z1_{(N'',\infty)}+U(e'-e)$ and $\|U^*u\|\le q^*(u)$.
(E3) $w_m$ norms $R_m^{**}\zh$ and $w'_m$ norms $R_m\hat x'$, so
$\ip{w_m}{R_m(\hat x'-\zh)}\le|R_m\hat x'|_m-|R^{**}_m\zh|_m\le
\ip{w'_m}{R_m(\hat x'-\zh)}$; both sides are bounded by
$\sum_k\lambda_{k,m}|u_{k,m}(\hat x'-\zh)|$, and $\sum_k\lambda_{k,m}\le1$.
Moreover $1/c'=p(\hat x')=1+\sum_m|R_m\hat x'|_m$ and $1/q_0=p^{**}(\zh)=
1+\sum_m|R^{**}_m\zh|_m$.
(E4) Off $P'_m$, Lemma~\ref{lem:threshold} at $x'$ gives
$\Phi_m(k)^2w'_m(k)/C'_m=(R_m\hat x')(k)/|R_m\hat x'|_m$, so
$d'_m(\omega)=\omega(R_m\hat x')/|R_m\hat x'|_m$, and likewise
$d_m(\omega)=\omega(R^{**}_m\zh)/|R^{**}_m\zh|_m$; use (E2), (E3).
(E5) $\ip{w'_m}{R_m\hat x'}=|R_m\hat x'|_m\ge\ip{w_m}{R_m\hat x'}$ and
$\ip{w'_m-w_m}{R_m^{**}\zh}\le0$; hence $0\le\mathrm{Bx}_m\le
\ip{w'_m-w_m}{R_m(\hat x'-\zh)}$, which is bounded by (E2).  Since
$w'_m\to w_m$ coordinatewise and $|w'_m-w_m|\le2$,
$\sum_k\lambda_{k,m}|w'_m(k)-w_m(k)|\to0$.
\end{proof}

\begin{lemma}[Lipschitz control of the block data]\label{lem:F1}
Along a construction sequence satisfying \eqref{eq:N1} there is $K_F$,
depending only on $f,g,\rho$, such that at late stages, for all $m$ and $k$,
\[
 |C'_m-C_m|=|M'_m-M_m|\le K_Fs_1,\qquad
 \Phi_m(k)|w'_m(k)-w_m(k)|\le K_F(s_1+t_{k,m}),
\]
and consequently $\|D_m(w'_m-w_m)\|_2^2\le K_F'(1+\log_2(1/s_1))s_1^2$, with
$K'_F$ depending only on $K_F$.
\end{lemma}

\begin{proof}
Fix $m$ and drop it.  Put $v_k:=m|u_k(\zh)|/|R^{**}\zh|$ and
$v'_k:=m|u_k(\hat x')|/|R\hat x'|$.  By Lemma~\ref{lem:threshold} (peaks
are the $k$ with $|\zeta(k)|/|\zeta|\ge\Phi_k^2M/C$, and off $P$,
$w(k)=C\zeta(k)/(\Phi_k^2|\zeta|)$), we have the \emph{clamp formula}
$\Phi_kw(k)=\sgn(u_k(\zh))\min(\Phi_kM,Cv_k)$, and the same at $x'$.
Squaring and summing, and using $M=1-C$, $C$ is a root of
\[
 F(c;v):=\sum_k\min\Big(\Phi_k\frac{1-c}c,v_k\Big)^2=1,
\]
and $C'$ is a root of $F(\cdot;v')=1$.  Each $F(\cdot;v)$ is nonincreasing on
$(0,1)$.  Since $\|\alpha\|_1=1$, there is $k_\natural\in P$ with
$\alpha(k_\natural)\ne0$, i.e.\ $Cv_{k_\natural}>\Phi_{k_\natural}M$; choose
$r>0$ with $[C-r,C+r]\subseteq(0,1)$ and $v_{k_\natural}\ge
\Phi_{k_\natural}(1-c)/c+r$ for $c\in[C-r,C+r]$.  At late stages
$|C'-C|\le r$ and $|v'_{k_\natural}-v_{k_\natural}|\le r$, so on
$[C-r,C+r]$ the $k_\natural$-term of $F(\cdot;v')$ equals
$\Phi_{k_\natural}^2((1-c)/c)^2$, whose derivative is at most
$-c_*:=-2\Phi_{k_\natural}^2(1-C-r)/(C+r)^3$; the other terms are
nonincreasing.  Hence $c_*|C'-C|\le|F(C;v')-F(C';v')|=|F(C;v')-F(C;v)|$, and
since $|\min(x,s)^2-\min(x,s')^2|\le2x|s-s'|$,
\[
 |C'-C|\le\frac{2(1-C)}{Cc_*}\sum_k\Phi_k|v'_k-v_k| .
\]
By (E2), (E3) and $|u_k(\zh)|\le1$, at late stages $|v'_k-v_k|\le
K(s_1+t_k+\tau(N''))$ with $K$ depending only on $f,g,\rho$; as
$\sum_k\Phi_k\le1$ and $\sum_k\Phi_kt_k\le\tau(N'')\le s_1^2$, this gives
$|C'-C|\le K_Fs_1$.  If $\sgn u_k(\hat x')=\sgn u_k(\zh)$ or one of them
vanishes, the clamp formula gives $|\Phi_kw'(k)-\Phi_kw(k)|\le
|C'-C|(\Phi_k+v_k)+C'|v'_k-v_k|$, and $v_k\le m/|R^{**}\zh|$.  If the signs
are opposite, both $|u_k(\hat x')|$ and $|u_k(\zh)|$ are at most
$|u_k(\hat x')-u_k(\zh)|$, and $|\Phi_kw'(k)|+|\Phi_kw(k)|\le C'v'_k+Cv_k$.
In both cases the bound follows from (E2).  Finally, using
$\min(x+y,c)^2\le\min(x,c)^2+2cy$ and $\Phi_k\le2^{-k}$,
\begin{align*}
 \|D(w'-w)\|_2^2&\le\sum_k\min(K_Fs_1,2\Phi_k)^2+4K_F\sum_k\Phi_kt_k\\
 &\le K_F^2s_1^2\lceil\log_2(1/s_1)\rceil+\tfrac43s_1^2+4K_Fs_1^2 .
 \qedhere
\end{align*}
\end{proof}

\begin{lemma}[Anchor]\label{lem:anchor}
Fix a block and drop its index.  Let $w,w'\in\ell_\infty$ with
$N(w)=N(w')=1$ be the norming functionals of nonzero block vectors, with
data $M,C,P$ and $M',C',P'$, and put $\gamma:=C'-C$; assume $C\le\frac12$
and $|\gamma|\le C/4$.  Let $S_1:=\{k\in P\cap P':\sgn w(k)=\sgn w'(k)\}$,
$S_2:=\{k\notin P\cup P':|2w'(k)-w(k)|\le M'-|\gamma|\}$,
$A:=\N\setminus(S_1\cup S_2)$ \textup{(}the \emph{scrambled set}\textup{)},
and
\[
 y(k):=2w'(k)-w(k)\ (k\in S_1\cup S_2),\qquad
 y(k):=(1-\gamma_+/M')w'(k)\ (k\in A).
\]
Then $N(y)\le1+\mathcal X/(2(C+2\gamma))$ with
$\mathcal X:=4\|D(w'-w)\|_2^2+4\sum_{k\in A}\Phi(k)^2$.  Moreover
$w'-w=(y-w')+Z_A$ with $Z_A:=(w'-w)1_A+(\gamma_+/M')w'1_A$,
$\|y-w'\|_\infty\le3$, and if $w'$ norms $R x'$ with $\sigma':=|Rx'|>0$,
then
\[
 N(y)-\frac{\ip{y}{Rx'}}{\sigma'}\le\frac{\mathcal X}{2(C+2\gamma)}+
 \frac1{\sigma'}\sum_{k\in A}|(Rx')(k)|\big(|w'(k)-w(k)|+\gamma_+\big).
\]
\end{lemma}

\begin{proof}
\emph{Sup norm.}  Note $M'-M=-\gamma$, $M\ge\frac12$ and $|\gamma|\le\frac18$.
On $S_1$, $|y(k)|=2M'-M=M'-\gamma$; on $S_2$, $|y(k)|\le M'-|\gamma|$; on
$A$, $|y(k)|\le M'-\gamma_+$.  So $\|y\|_\infty\le M'-\gamma$.

\emph{Hilbert part.}  Put $Y:=y-w'=(w'-w)1_{S_1\cup S_2}-(\gamma_+/M')w'1_A$
and $R_A:=\ip{Dw'}{D(w'-w)1_A}+(\gamma_+/M')\|Dw'1_A\|^2$.  Then
$\ip{Dw'}{DY}=\ip{Dw'}{D(w'-w)}-R_A$ and $2\ip{Dw'}{D(w'-w)}=C'^2-C^2+
\|D(w'-w)\|^2$, so
\begin{gather*}
 \|Dy\|^2=C'^2+2\ip{Dw'}{DY}+\|DY\|^2=2C'^2-C^2+\mathcal X_0,\\
 \mathcal X_0:=\|D(w'-w)\|^2+\|DY\|^2-2R_A .
\end{gather*}
Since $1_{S_1\cup S_2}$ and $1_A$ are disjoint and $\gamma_+/M'\le1$,
$\|DY\|^2\le\|D(w'-w)\|^2+\|Dw'1_A\|^2$ and $-2R_A\le\|Dw'1_A\|^2+
\|D(w'-w)\|^2$; with $|w'|\le1$ this gives $\mathcal X_0\le\mathcal X$.  As
$2C'^2-C^2=(C+2\gamma)^2-2\gamma^2$ and $C+2\gamma\ge C/2>0$,
$\|Dy\|\le C+2\gamma+\mathcal X/(2(C+2\gamma))$.  Adding,
$N(y)\le M'+C+\gamma+\mathcal X/(2(C+2\gamma))=1+\mathcal X/(2(C+2\gamma))$.

\emph{The rest.}  The identity for $w'-w$ is immediate, and
$\|y-w'\|_\infty\le\|w'-w\|_\infty+1\le3$.  Finally $\ip{y}{Rx'}=\sigma'+
\ip{Y}{Rx'}$ and $\ip{Y}{Rx'}=\ip{w'-w}{Rx'}-\ip{(w'-w)1_A}{Rx'}-
(\gamma_+/M')\ip{w'1_A}{Rx'}$, where the first term is $\ge0$ because $w'$
norms $Rx'$ and $N(w)=1$.
\end{proof}

\begin{definition}[Scrambling condition]\label{def:SC}
For a block $m$ and $s>0$ let
\[
 \mathrm{Scr}_m(s):=\sum_{\substack{k\in P_m\\ \mu_{k,m}\le s}}
 \min(\Phi_m(k),s)+\sum_{\substack{k\in Q_m\\ \Phi_m(k)\gap_m(k)\le s\ \text{or}
 \ \gap_m(k)\le s}}\min(\Phi_m(k),s).
\]
A set $I_-\subseteq I$ of blocks satisfies the \emph{scrambling condition}
\textup{(SC)} if for every $K\ge1$ there is a sequence $s_i\downarrow0$
\textup{(}common to all $m\in I_-$\textup{)} with
$\max_{m\in I_-}\mathrm{Scr}_m(Ks_i)/s_i\to0$.
\end{definition}

The common sequence matters: each $\mathrm{Scr}_m$ may be $o(s)$ along a
sequence of scales while the good scales of two blocks are disjoint.
\textup{(SC)} excludes degenerate peaks in the blocks of $I_-$ (each
contributes $\min(\Phi_m(k),Ks)=Ks$ for small $s$).  If $f$ satisfies
\textup{(BT)}, then every block satisfies $\mathrm{Scr}_m(Ks)=o(s)$ as
$s\to0$ for every $K$: for small $s$ the finitely many strict non-peaks do not
contribute, and the peak part is at most
$\sum\{\Phi_m(k):\mu_{k,m}<2Ks\}=o(s)$ by \textup{(MS)}.

\begin{lemma}[Scrambling estimate]\label{lem:scrambling}
There is $K_*\ge1$, depending only on $f,g,\rho$, with the following
property.  Let $I_-$ satisfy \textup{(SC)}, let $(s_i)$ be the sequence of
\textup{(SC)} for $K:=2K_*$, and consider a construction sequence with
$s_{1,i}=s_i$ and $N''_i$ so large that \eqref{eq:N1} and
\begin{equation}\label{eq:N2}
 \sum_{m\in I}\sum_k\Phi_m(k)\,1\big[t_{k,m}>s_{1,i}\big]\le s_{1,i}^2
\end{equation}
hold.  Then for $m\in I_-$, with $A_m$ the scrambled set of
Lemma~\ref{lem:anchor} for $(w_m,w'_m)$,
\begin{multline*}
 \mathfrak S_m:=\|D_m(w'_m-w_m)\|_2^2\\
 +\sum_{k\in A_m}\Big(\Phi_m(k)^2+
 \lambda_{k,m}\big(|w'_m(k)-w_m(k)|+(C'_m-C_m)_+\big)\Big)=o(s_1).
\end{multline*}
\end{lemma}

Condition \eqref{eq:N2} can be met by choosing $N''$ after $s_1$: for fixed
$s_1$ the left side tends to $0$ as $N''\to\infty$ by dominated convergence.

\begin{proof}
Fix $m\in I_-$ and drop it.  The first term is $o(s_1)$ by
Lemma~\ref{lem:F1}.  At late stages $|C'-C|\le C/4$, so
Lemma~\ref{lem:anchor} applies ($C\le\|\Phi\|_2<\frac12$).

\emph{Peaks.}  Put $\hat\vartheta:=|R^{**}\zh|M/(mC)$ and $\hat\vartheta':=
|R\hat x'|M'/(mC')$, so that $\mu_k=q_0(|u_k(\zh)|-\hat\vartheta\Phi_k)$, and
$k$ is a peak of $w'$ of sign $\sgn u_k(\hat x')$ as soon as
$|u_k(\hat x')|>\hat\vartheta'\Phi_k$ (Lemma~\ref{lem:threshold}).  By (E3)
and Lemma~\ref{lem:F1}, $|\hat\vartheta'-\hat\vartheta|\le K_\vartheta s_1$
at late stages.  If $k\in P$ and $\mu_k>K_*(s_1+t_k)$ with $K_*\ge
K_e+K_\vartheta+1$, then by (E2) $|u_k(\hat x')-u_k(\zh)|<|u_k(\zh)|$ and
$|u_k(\hat x')|-\hat\vartheta'\Phi_k>0$; so $k\in P'$ with the same sign,
i.e.\ $k\in S_1$.

\emph{Strict non-peaks.}  If $k\in Q$, $\Phi_k\gap_k>K_*(s_1+t_k)$ and
$\gap_k>K_*s_1$ with $K_*\ge4K_F$, then by Lemma~\ref{lem:F1}
$|w'(k)-w(k)|<\gap_k/4$ and $|\gamma|=|M'-M|<\gap_k/4$.  Hence $|w'(k)|<M'$,
so $k\notin P'$, and $|2w'(k)-w(k)|\le|w(k)|+2|w'(k)-w(k)|<M-\gap_k/2\le
M'-|\gamma|$; i.e.\ $k\in S_2$.

Put $K_*:=\max(K_e+K_\vartheta+1,4K_F)$.  Let $A^{(1)}:=\{k\in A:t_k\le
s_1\}$ and $A^{(2)}:=A\setminus A^{(1)}$.  By the two steps, every $k\in
A^{(1)}$ is a peak with $\mu_k\le2K_*s_1$ or a strict non-peak with
$\Phi_k\gap_k\le2K_*s_1$ or $\gap_k\le2K_*s_1$; so
$\sum_{A^{(1)}}\min(\Phi_k,2K_*s_1)\le\mathrm{Scr}(2K_*s_1)=o(s_1)$.  In
particular, at late stages no $k\in A^{(1)}$ has $\Phi_k\ge2K_*s_1$, and
$\sum_{A^{(1)}}\Phi_k=o(s_1)$.  By \eqref{eq:N2}, $\sum_{A^{(2)}}\Phi_k\le
s_1^2$.  Since $\Phi_k\le1$, $|w'-w|\le2$, $\gamma_+\le1$ and
$\lambda_k=m\Phi_k$, the sum over $A$ in $\mathfrak S$ is at most
$4m\sum_A\Phi_k=o(s_1)$.
\end{proof}

\subsection{Engineering without steering}

\begin{theorem}[Engineered recovery]\label{thm:engineered}
Let $I$ be finite, let $f\in S_{p^*}$ have finite base support $F$, let
$g\in\Cset(f)$ carry two-piece data $(b^\pm,\omega^\pm)$, and let
$\rho\in(0,1)$ with $\rho^2\kappa_w<1$.  Assume that
$I_-:=\{m:\Delta d_m<0\}$ satisfies the scrambling condition \textup{(SC)}
\textup{(}no assumption if $I_-=\emptyset$\textup{)}.  Then there are
norm-attaining $f'_i\to f$ \textup{(}engineered approximants\textup{)} and
$g'_i\in\Cset(f'_i)$ with $g'_i\to\rho g$; in particular
$(f,\rho g)\in\cl\NA((c_0,p),\ell_2^2)$.  If moreover $\kappa_w\le1$, this
holds for every $\rho<1$, and $(f,g)\in\cl\NA((c_0,p),\ell_2^2)$.
No hypothesis is made on the contact set $K$, on the number of active
blocks, on the sets $Q_m$ \textup{(}for $m\notin I_-$\textup{)}, or on
rates of $T$.
\end{theorem}

\begin{proof}
We use the notation of Definition~\ref{def:engineered}; data of $f'$ carry
primes, $\sigma'_m=|R_mx'|_m$, $c'=q(x')$, and $d'_m$ is as in (E4).  Put
$d^\sigma_m:=d_m(\omega^\sigma_m)$ for $\sigma\in\{+,-,\theta\}$, so
$d^\theta_m=\frac12(d^+_m+d^-_m)$.

\emph{Step 0: fixed data.}  Let $\delta:=(1-\rho^2\kappa_w)/2$ and
$\Gamma_\sigma:=\Gamma_w(b^\sigma,\omega^\sigma)$; by convexity
$\Gamma_\theta\le\kappa_w$, so $\rho^2\Gamma_\sigma\le1-2\delta$ for all
three $\sigma$.  Let $\gamma_0>0$ be the least gap $\gap_m(k)$, $k\in
\Omega_m$, $m\in I$, and $a_{\min}:=\min_F|a_j|$.  In every block fix
$\gamma_m\in(M_m/2,M_m)$ with $|w_m(k)|\ne\gamma_m$ for all $k$.  By
Lemma~\ref{lem:approxfacts}, Lemma~\ref{lem:F1}, Lemma~\ref{lem:scrambling}
there are $T_1\in(0,1]$ and $K_0\ge1$, depending only on $f,g,\rho$, such
that the conditions marked $(T_1)$ below hold for $T_0\le T_1$, and such
that the quantities of Step~3 satisfy, at late stages and for
$0<|\tau|\le T_1$, $|\lin'_m|\le K_0|\tau|s_1$ and $\hat G_b,\hat G_m\le
K_0\tau^2$ (for $\sigma=\pm$ we use $|\tau|>s_1$ to bound the terms
$|\tau|\,o(s_1)$ by $\tau^2$).  Put
\[
 \eta_1:=\min\Big(\frac14,\frac\delta{192|I|K_0}\Big),
\]
and choose transfer data $(\gamma_m,k^\varsigma_{*,m},\Lambda^\varsigma_m)$ of
$f$ in every block by Lemma~\ref{lem:transferdata} with this $\eta_1$.  By
Lemma~\ref{lem:persistence} the transfer vectors $y'^\varsigma_m$ built at
$f'$ satisfy, at late stages, $e'^\varsigma_{y,m}\ge\frac12$,
$\iota'^\varsigma_m\le2\eta_1$, $\|D_my'^\varsigma_m\|_2\le K_y$ and
$|Y'^\varsigma_m|/c'\le K_Y$, where $K_y,K_Y$ do not depend on $\eta_1$
(Lemma~\ref{lem:transferdata}(b),(d)).  Finally choose $T_0\in(0,T_1]$ with
$T_0^2\le4\delta$ and so small that the finitely many conditions marked
$(T_0)$ below hold; they involve $f,g,\rho$, $\eta_1$ and the transfer
data.

\emph{Construction.}  If $I_-\ne\emptyset$ let $(s_i)$ be the sequence of
\textup{(SC)} for $K:=2K_*$ (Lemma~\ref{lem:scrambling}); otherwise let
$s_i\downarrow0$ be arbitrary.  Consider the engineered approximants with
$N_i\to\infty$, $s_{1,i}:=s_i$, and $N''_i$ chosen after $s_{1,i}$ so that
\eqref{eq:N1}, \eqref{eq:N2} and
\begin{equation}\label{eq:N3}
 \rho V_{>N''}\le\delta s_1/16,\qquad V_{>N''}:=\sum_{j\in K,\,j>N''}|v_j|,
\end{equation}
hold.  All statements below are at late stages.  By
Lemma~\ref{lem:approxfacts}(b), $w'_m\to w_m$ coordinatewise, $C'_m\to C_m$,
$M'_m\to M_m$, $\sigma'_m\to\sigma_m$, $c'\to q_0$, $\nu'\to\nu$, $e'\to e$,
and $R^*_mw'_m\to R_m^*w_m$ in $\ell_1$; in particular every $k\in\Omega_m$
is a strict non-peak of $w'_m$ with $\gap'_m(k)\ge\gamma_0/2$, the
coordinatewise radius of $\omega^\sigma_m$ at $f'$ is at least half the one
at $f$, and $d'_m(\omega)\to d_m(\omega)$, $H'_m(\omega)\to H_m(\omega)$ for
$\omega\in\{\omega^\sigma_m\}$.

\emph{Step 1: the target.}  Let
\[
 g'':=b^\theta1_{[1,N]}+\sum_mR_m^*\big(\omega^\theta_m-d'_m(\omega^\theta_m)
 w'_m\big),\qquad c:=g''(\hat x'),\qquad g':=\rho(g''-ca').
\]
Then $g'(\hat x')=0$.  Since $(b^\theta,\omega^\theta)$ represents $g$,
$g''-g=-b^\theta1_{(N,\infty)}+\sum_mR_m^*(d^\theta_mw_m-d'_m(\omega^\theta_m)
w'_m)\to0$, and $c=(g''-g)(\hat x')+g(\hat x')\to g(\zh)=0$ because $\hat
x'\to\zh$ weak$^*$ boundedly.  So $g'\to\rho g$.

\emph{Step 2: three exact decompositions.}  Let
$\beta^\theta:=b^\theta1_{[1,N]}$, $\beta^\pm:=b^\pm1_{[1,N]}\pm\frac12v
1_{(N,\infty)}$, $B^\sigma:=\rho(\beta^\sigma-ca')$, and
\[
 \Omega^\theta_m:=\rho\big(\omega^\theta_m-d'_m(\omega^\theta_m)w'_m\big),\qquad
 \Omega^\pm_m:=\rho\big(\omega^\pm_m-d'_m(\omega^\theta_m)w'_m-(d^\pm_m-
 d^\theta_m)w_m\big).
\]
Since $\beta^\pm-\beta^\theta=\pm\frac12v$, $\omega^\pm-\omega^\theta=
\mp\frac12(\omega^--\omega^+)$, $d^\pm-d^\theta=\mp\frac12\Delta d$ and
$v=\sum_mR_m^*\big((\omega^-_m-\omega^+_m)-\Delta d_mw_m\big)$, we get
$g'=B^\sigma+\sum_mR^*_m\Omega^\sigma_m$ for $\sigma\in\{+,-,\theta\}$, and
so $f'+\tau g'=(a'+\tau B^\sigma)+\sum_mR^*_m(w'_m+\tau\Omega^\sigma_m)$ for
every $\tau$.  We use $\sigma=\theta$ for $0<|\tau|\le s_1$ and
$\sigma=\sgn\tau$ for $s_1<|\tau|\le T_0$.

For $\sigma=\pm$ rewrite, with $\kappa^\sigma_m:=\rho\,(d'_m-d_m)(\omega^\sigma_m
-\omega^\theta_m)$,
\[
 \Omega^\sigma_m=\rho\big(\omega^\sigma_m-d'_m(\omega^\sigma_m)w'_m\big)+
 \kappa^\sigma_mw'_m+\rho(d^\sigma_m-d^\theta_m)(w'_m-w_m).
\]
For $\sigma=\sgn\tau$ we have $\tau\rho(d^\sigma_m-d^\theta_m)=-r_m\sgn
(\Delta d_m)$ with $r_m:=\frac12|\tau|\rho|\Delta d_m|$.  Put
$W_{0,m}:=w'_m+\tau\rho(\omega^\sigma_m-d'_m(\omega^\sigma_m)w'_m)+
\tau\kappa^\sigma_mw'_m$.  If $\Delta d_m\ge0$, let $Y_m:=w_m$ and
$W^\natural_m:=w'_m+\tau\Omega^\sigma_m=W_{0,m}-r_mw'_m+r_mY_m$.  If
$\Delta d_m<0$, let $Y_m:=y_{A,m}$ and $Z_{A,m}$ be given by
Lemma~\ref{lem:anchor} for $(w_m,w'_m)$, and put $W^\natural_m:=w'_m+\tau
\Omega^\sigma_m-r_mZ_{A,m}=W_{0,m}-r_mw'_m+r_mY_m$; the remainder
$r_mR^*_mZ_{A,m}$ is moved to the base.  The decomposition
\[
 f'+\tau g'=A_\tau+\sum_mR^*_mW^\natural_m,\qquad A_\tau:=a'+\tau B^\sigma
 +\mathcal Z,\quad\mathcal Z:=\sum_{m\in I_-}r_mR^*_mZ_{A,m},
\]
is exact.  For $\sigma=\theta$ put $W^\natural_m:=w'_m+\tau\Omega^\theta_m$,
$A_\tau:=a'+\tau B^\theta$.  In all cases $\|W^\natural_m-w'_m\|_\infty+
\|D_m(W^\natural_m-w'_m)\|_2\le K_W|\tau|$ with $K_W$ depending only on
$f,g,\rho$, because $\|Y_m-w'_m\|_\infty\le3$.

\emph{Step 3: first-order terms and excesses.}
\emph{Blocks.}  Since $\omega^\sigma_m$ vanishes on $P'_m$,
$\lin'_m(w'_m+\tau\rho(\omega^\sigma_m-d'_m(\omega^\sigma_m)w'_m))=0$
(Lemma~\ref{lem:algebra} at $f'$).  Hence $\lin'_m=0$ for $\sigma=\theta$, and
for $\sigma=\pm$
\[
 \lin'_m(W^\natural_m)=\tau\kappa^\sigma_m+r_m\ip{Y_m-w'_m}{R_mx'}/\sigma'_m .
\]
By (E4) and \eqref{eq:N1}, $|\kappa^\sigma_m|\le Ks_1$.  For $Y_m=w_m$,
$\ip{Y_m-w'_m}{R_mx'}=-c'\mathrm{Bx}_m$; for $Y_m=y_{A,m}$ its modulus is at
most $c'\mathrm{Bx}_m+c'\sum_{A_m}\lambda_{k,m}(|w'_m-w_m|(k)+(C'_m-C_m)_+)$
(proof of Lemma~\ref{lem:anchor} and $|(R_mx')(k)|\le\lambda_{k,m}c'$).  By
(E5) and Lemma~\ref{lem:scrambling} both are $o(s_1)$; so
$|\lin'_m|\le K_0|\tau|s_1$.  For the excess write
$(1-r_m)V_m:=W_{0,m}-r_mw'_m=(1-r_m+\tau\kappa^\sigma_m)\big(w'_m+t'
(\omega^\sigma_m-d'_m(\omega^\sigma_m)w'_m)\big)$ with
$t':=\tau\rho/(1-r_m+\tau\kappa^\sigma_m)$.  By convexity and homogeneity of
$N_m$ and Lemma~\ref{lem:block}(d) at $f'$ (condition $(T_1)$: $2\rho T_0$ is
at most half the radius of every $\omega^\sigma_m$ at $f$, and
$\rho T_0\max_m|\Delta d_m|+2T_0Ks_1\le\frac14$),
\[
 N_m(W^\natural_m)\le(1-r_m+\tau\kappa^\sigma_m)+\frac{\rho^2\tau^2}2
 H'_m(\omega^\sigma_m)(1+K_HT_0)+r_mN_m(Y_m),
\]
and therefore, using $\ip{w'_m}{R_mx'}=\sigma'_m$,
\[
 G'_m(W^\natural_m)\le\hat G_m:=\frac{\rho^2\tau^2}2H'_m(\omega^\sigma_m)
 (1+K_HT_0)+r_m\mathfrak g_m,\qquad
 \mathfrak g_m:=N_m(Y_m)-\frac{\ip{Y_m}{R_mx'}}{\sigma'_m}.
\]
For $Y_m=w_m$, $\mathfrak g_m=c'\mathrm{Bx}_m/\sigma'_m=o(s_1)$ by (E5);
for $Y_m=y_{A,m}$, $\mathfrak g_m\le K\mathfrak S_m=o(s_1)$ by
Lemmas~\ref{lem:anchor} and~\ref{lem:scrambling}.  For $\sigma=\theta$ the
same computation with $r_m=\kappa=0$ gives $\hat G_m:=\frac{\rho^2\tau^2}2
H'_m(\omega^\theta_m)(1+K_HT_0)$.

\emph{Base.}  By Lemma~\ref{lem:bookkeeping}(b) at $f'$,
$G'_b(A_\tau)=\Exc'(\tau B^\sigma+\mathcal Z)+\nu'\Psi'(U^*(\tau B^\sigma+
\mathcal Z)/\nu')$ and $\Exc'(\tau B^\sigma+\mathcal Z)\le\Exc'(\tau B^\sigma)
+2\|\mathcal Z\|_1$.  We claim $\Exc'(\tau B^\sigma)\le\frac12\rho|\tau|
V_{>N''}$, and $=0$ for $\sigma=\theta$.  On $\supp a'$,
$a'_j+\tau B^\sigma_j=(1-\tau\rho c)a'_j+\tau\rho\beta^\sigma_j$ with
$|\tau\rho c|\le\frac12$.  On $F$, $|a'_j|\ge a_{\min}/2$ and
$|\tau\rho\beta^\sigma_j|\le a_{\min}/8$ (condition $(T_1)$: $\rho T_0
(\|b^+\|_\infty+\|b^-\|_\infty)\le a_{\min}/8$): no flip.  On a window
contact with mass: for $\sigma=\theta$, $|\tau\rho\beta^\theta_j|\le
s_1\rho|b^\theta_j|=m_j/4\le|a'_j|/2$; for $\sigma=+$, $\tau>0$ and
$z_jb^+_j\ge0$, so both terms have the sign $z_j=\sgn a'_j$ or vanish; for
$\sigma=-$, $\tau<0$ and $z_jb^-_j\le0$, with the same conclusion.  Off
$\supp a'$, $B^\sigma_j=\rho\beta^\sigma_j$, and the summand of $\Exc'$ is
$|\tau\rho\beta^\sigma_j|-z'_j\tau\rho\beta^\sigma_j$: it vanishes off
$F\cup K$ (where $b^\pm$ and $v$ vanish), on window contacts without mass
(there $\beta^\theta_j=0$ and $\tau\beta^\pm_j=\tau b^\pm_j$ has the sign
$z_j=z'_j$ on the designated side), and on contacts in $(N,N'']$ (there
$\tau\beta^\pm_j=\pm\frac12\tau v_j$ has the sign $z_j$ because $z_jv_j\ge0$,
and $\beta^\theta_j=0$); on contacts $j>N''$, $z'_j=0$ and the summand is at
most $\frac12\rho|\tau||v_j|$, and it is $0$ for $\sigma=\theta$.  This proves
the claim.  Further $\|\mathcal Z\|_1\le\sum_{m\in I_-}r_m\sum_{A_m}
\lambda_{k,m}(|w'_m-w_m|(k)+(C'_m-C_m)_+)=|\tau|o(s_1)$ by
Lemma~\ref{lem:scrambling}, and $P'^\perp U^*a'=0$, so by
Lemma~\ref{lem:base}
\[
 G'_b(A_\tau)\le\hat G_b:=\frac{\rho^2\tau^2}2h'(\beta^\sigma)(1+K_hT_0)
 +\tfrac12\rho|\tau|V_{>N''}+|\tau|\,o(s_1),
\]
where $h'(\beta):=\|P'^\perp U^*\beta\|^2/\nu'\to h(b^\sigma)$ for
$\beta=\beta^\sigma$ (as $\beta^\sigma-b^\sigma=-b^\theta1_{(N,\infty)}\to0$
in $\ell_1$); for $\sigma=\theta$ the last two terms are absent.

\emph{Step 4: rebalancing.}  Apply Proposition~\ref{prop:rebalancing} at
$f'$ with $h:=\tau g'$ (so $h(x')=0$), the decomposition of Step~2, the
bounds $\hat G_b,\hat G_m$ of Step~3, and $y_m:=y'^\downarrow_m$ or
$y'^\uparrow_m$ according to the sign of $\lin'_m+\hat G_m-\hat\Gamma$.
Then $|\varepsilon_m|\le2(|\lin'_m|+\hat G_m+\hat\Gamma)\le6K_0\tau^2$ when
$|\tau|>s_1$, and $\le4K_0\tau^2$ for $\sigma=\theta$.  Condition
\eqref{eq:R1} follows from Lemma~\ref{lem:TV} with $\eta:=K_W|\tau|$
(condition $(T_0)$: $K_WT_0\le\min_m(M_m-\gamma_m)/8$, $6K_0T_0^2(2+
\Lambda^\varsigma_m)\le\gamma_m/2$, $6K_0T_0^2\le\min_m(M_m-\gamma_m)/8$ and
$K_WT_0\le\min_mC_m/4$), and $|\lambda|\le6|I|K_0K_YT_0^2\le1$ by
condition $(T_0)$.  The error $E$ of Proposition~\ref{prop:rebalancing}
satisfies
\[
 E\le12|I|K_0\eta_1\tau^2+6|I|K_0K_Y|\tau|^3K_A+\frac{(48K_0K_WK_y|\tau|^3
 +36K_0^2K_y^2\tau^4)}{\min_mC_m/2},
\]
where $K_A$ bounds $(|q^*(A_\tau)-1|+q^*(A_\tau-a'))/|\tau|$.  By the choice
of $\eta_1$ and condition $(T_0)$ on the last two terms, $E\le\delta\tau^2/8$.

\emph{Step 5: the levels.}  By Step~3,
$\hat\Gamma=c'\hat G_b+\sum_m\sigma'_m\hat G_m\le\frac{\rho^2\tau^2}2
\big(c'h'(\beta^\sigma)+\sum_m\sigma'_mH'_m(\omega^\sigma_m)\big)(1+K'T_0)
+\frac12\rho|\tau|V_{>N''}+|\tau|o(s_1)$.  At late stages the bracket is at
most $\Gamma_\sigma+\delta/8$; by \eqref{eq:N3} and $|\tau|>s_1$ the last two
terms are at most $\delta\tau^2/16$ (and they are absent for
$\sigma=\theta$).  With condition $(T_0)$: $K'T_0(1+\rho^2\kappa_w)\le\delta/8$,
\[
 p^*(f'+\tau g')\le1+\hat\Gamma+E\le1+\frac{\tau^2}2\Big(\rho^2\kappa_w+
 \frac\delta8+\frac\delta8+\frac\delta8+\frac\delta4\Big)\le1+\frac{\tau^2}2
 (1-\delta)
\]
for $0<|\tau|\le T_0$, since $\rho^2\kappa_w=1-2\delta$.

\emph{Step 6: conclusion.}  At late stages $p^*(f'-f)\le(1-\rho^2)T_0^2/6$
and $p^*(g'-\rho g)\le(1-\rho^2)T_0/6$.  Lemma~\ref{lem:assembly} gives
$g'\in\Cset(f')$, and $f'$ attains its norm.  If $\kappa_w\le1$, then
$\rho^2\kappa_w<1$ for every $\rho<1$, and $(f,g)=\lim_{\rho\to1}(f,\rho g)
\in\cl\NA$.
\end{proof}

\begin{remark}[On the proof]\label{rem:engineeredproof}
(a) Earlier engineered constructions had to make the two side
decompositions represent the \emph{same} functional at the new normer
\textup{(}``steering''\textup{)}: far contacts with flipped signs and negative
masses, a tuning mass fixed by the intermediate value theorem, a spanning
hypothesis for several blocks.  None of this is needed here.  The linear
mismatches between the pieces are of order $|\tau|s_1$, and
Proposition~\ref{prop:rebalancing} removes them at the cost
$|\tau|s_1\iota$, the inefficiency $\iota$ being fixed before $s_1$; the
first-order identity of Lemma~\ref{lem:bookkeeping}(a) guarantees that a
linear mismatch is a pure shift of levels.
(b) The only genuinely nonlinear first-order costs are the kinks beyond the
cut-off, the Bregman terms $\mathrm{Bx}_m$ (always $o(s_1)$) and, in blocks
with $\Delta d_m<0$, the anchor costs controlled by \textup{(SC)}.
(c) The approximants are not the canonical truncations; this is consistent
with Proposition~\ref{prop:nonrecovery}.
\end{remark}

\subsection{Consequences}

\begin{corollary}[Weighted engineered recovery]\label{cor:weightedengineered}
Let $I$ be finite, $f\in S_{p^*}$ with $F$ finite, and let $g\in\Cset(f)$
carry $d$-neutral two-piece data with any number of active blocks.  If
$\rho^2\kappa_w<1$, then $(f,\rho g)\in\cl\NA((c_0,p),\ell_2^2)$; if
$\kappa_w\le1$, then $(f,g)\in\cl\NA((c_0,p),\ell_2^2)$.
\end{corollary}

\begin{proof}
$I_-=\emptyset$; apply Theorem~\ref{thm:engineered}.
\end{proof}

Since $q_0+\sum_m\sigma_m=1$, the coefficient $\kappa_w$ is at most the
max-form coefficient
\[
 \kappa_{\max}:=\max_\pm\max\big(h(b^\pm),\max_mH_m(\omega^\pm_m)\big),
\]
and it can be much smaller.  Corollary~\ref{cor:weightedengineered} therefore contains the
engineered recovery of $d$-neutral two-piece mates of the companion notes
(which used the max-form coefficient, and for several active blocks a
spanning hypothesis), with a self-contained proof.

\begin{corollary}[Non-negative $d$-mismatch]\label{cor:D1}
Let $I$ be finite, $f\in S_{p^*}$ with $F$ finite, and let $g\in\Cset(f)$
carry two-piece data with $\Delta d_m\ge0$ for every $m$ and $\kappa_w\le1$.
Then $(f,g)\in\cl\NA((c_0,p),\ell_2^2)$.  No hypothesis is made on $Q_m$
\textup{(}infinitely many strict non-peaks, near-threshold coordinates and
degenerate peaks are allowed\textup{)}, on $K$, or on $T$.
\end{corollary}

\begin{proof}
$I_-=\emptyset$; apply Theorem~\ref{thm:engineered}.  The fine structure of
the blocks enters the proof only through (E4), (E5) and
Lemma~\ref{lem:F1}, which hold at every $f$ with $F$ finite.
\end{proof}

\begin{corollary}[Block-tame points are recoverable]\label{cor:BTrecovered}
Let $I$ be finite.  Every $f\in S_{p^*}$ satisfying \textup{(BT)} belongs to
$\Rec$; that is, $(f,g)\in\cl\NA((c_0,p),\ell_2^2)$ for every
$g\in\Cset(f)$, whatever the contact set $K$.
\end{corollary}

\begin{proof}
By Theorem~\ref{thm:onesided}(c), every $g\in\Cset(f)$ carries two-piece
data with $\kappa_w\le1$.  By the remark after Definition~\ref{def:SC},
\textup{(BT)} implies $\mathrm{Scr}_m(Ks)=o(s)$ as $s\to0$ for every block
and every $K$, so every set of blocks satisfies \textup{(SC)}.
Theorem~\ref{thm:engineered} gives $(f,\rho g)\in\cl\NA$ for every
$\rho<1$, hence $(f,g)\in\cl\NA$.
\end{proof}

Corollary~\ref{cor:BTrecovered} contains Theorem~\ref{thm:tame}
(which needed $K$ finite) and Proposition~\ref{prop:intrinsic}(c) at
\textup{(BT)} points.

\begin{corollary}[The example of Section~\ref{sec:resonance}]\label{cor:exampleRecovered}
Let $T$, $I\ni1$ and $f$ be as in Lemma~\ref{lem:firstrow}.  Then $f$
satisfies \textup{(BT)}, and:
\begin{enumerate}
\item[(a)] $f\in\Rec$: every mate of $f$, in particular every element of the
defect $\Def(f)$ of Theorem~\ref{thm:defect}, is recovered along engineered
norm-attaining approximants;
\item[(b)] every $g\in\Cset(f)$ is supported in $\{1\}\cup K_0$, and
$\mu^+\le g_j/u_{2,1}(j)\le\mu^-$ for $j\in K_0$, where
$(\mu^\pm/\lambda_0)e_2$ \textup{(}in block $1$\textup{)} are the block parts
of optimal two-piece data of $g$.
\end{enumerate}
\end{corollary}

\begin{proof}
\textup{(BT)}: $F=\{1\}$, $Q_1=\{2\}$ and $Q_m=\emptyset$ for $m\ne1$
(Lemma~\ref{lem:firstrow}).  By the proof of Lemma~\ref{lem:firstrow},
$\vartheta_m\Phi_m(k)\le q_0\varpi_{k,m}$ for $k\ge2$ and
$\vartheta_m\Phi_m(1)\le q_0/2$.  With (P2), (P3) this gives
$\mu_{1,m}\ge q_0(\frac79-\frac12)\ge q_0/4$, $\mu_{k,m}\ge q_0(\frac12-
\varpi_{k,m})\ge q_0/4$ for exceptional $(k,m)$, and $\mu_{k,m}\ge
q_0(\frac32r_l-r_l^2)\ge q_0r_l$ ($l=\ell(k,m)$, $r_l\le\frac12$) for the
other peaks; so there are no degenerate peaks.  For $s<q_0/4$,
$\mu_{k,m}<s$ forces a non-exceptional $(k,m)$ with $q_0r_l<s$, i.e.\
$2\Phi_m(k)\le\varpi_{k,m}<s^2/q_0^2$.  Since $\Phi_m(k+1)\le\Phi_m(k)/2$, the
sum of $\Phi_m(k)$ over these $k$ is at most $s^2/q_0^2=o(s)$: \textup{(MS)}
holds.  (a) is Corollary~\ref{cor:BTrecovered}.  (b) By
Theorem~\ref{thm:onesided}(c), $g=b^\pm+R_1^*(\omega^\pm_1-d_1(\omega^\pm_1)
w_1)$ with $\omega^\pm_1=(\mu^\pm/\lambda_0)e_2$; since $w_1(2)=0$,
$d_1(\omega^\pm_1)=0$ and $g=b^\pm+\mu^\pm u_{2,1}$.  The vectors $b^\pm$ are
supported in $F\cup K=\{1\}\cup K_0$, with $b^+_j\ge0\ge b^-_j$ on $K_0$
(where $z_j=1$), and $\supp u_{2,1}=\{1\}\cup K_0$ with $u_{2,1}>0$ on $K_0$
(P1).
\end{proof}

Part (b) is the converse statement of Remark~\ref{rem:twopiece}(b), and
(a) shows that the defect of Theorem~\ref{thm:defect} is a defect of the
intrinsic mechanisms of Sections~\ref{sec:certificates}--\ref{sec:second}
only.  The explicit mates $g_{K_1}$ are also covered directly by
Corollary~\ref{cor:weightedengineered}, through their $d$-neutral data
described after Definition~\ref{def:twopiece} (for $\rho$ with
$\rho^2\kappa_w<1$).

\begin{remark}[Negative $d$-mismatch without \textup{(SC)}]\label{rem:withoutSC}
(a) Without \textup{(SC)}, the proof of Theorem~\ref{thm:engineered} still
gives $(f,\rho g)\in\cl\NA$ whenever
\[
 \rho^2\kappa_w+\rho\sum_{m\in I_-}|\Delta d_m|K^{\rm scr}_m<1,\qquad
 K^{\rm scr}_m:=\limsup_i\frac{\mathfrak g_m+2\|R^*_mZ_{A,m}\|_1}{s_{1,i}},
\]
along a construction sequence for which all $K^{\rm scr}_m$ are finite: in
Step~5 the anchor costs then contribute at most
$\frac{\tau^2}2\rho\sum_{I_-}|\Delta d_m|K^{\rm scr}_m(1+o(1))$ for
$|\tau|>s_1$.  Finiteness of $K^{\rm scr}_m$ is not automatic: the truncated
mass of near-threshold peaks can be of order $s\log(1/s)$ for admissible
$T$; it holds, for instance, under \textup{(MS)} together with gaps of the
strict non-peaks bounded below off a finite set.  So without \textup{(SC)}
the possible failure of recovery for $\Delta d<0$ is confined to $\rho$
close to $1$ only in such cases.
(b) \emph{Sketch, not used.}  One can construct admissible operators for
which \textup{(SC)} fails at a first row (prescribe block vectors in
$\zh^\perp$ at a positive proportion of the indices of one block, so that
these become strict non-peaks with $w_m(k)=0$ at all fine scales).  Then the
estimate of Theorem~\ref{thm:engineered} does not apply to mates with
$\Delta d<0$ in that block for $\rho$ close to $1$.  Whether these mates are
recovered along other approximants is open; that the approximants of
Definition~\ref{def:engineered} fail to recover them is \emph{not} proved.
\end{remark}

\section{A designed admissible operator}\label{sec:designed}

Mart\'in's construction works with every admissible operator $T$
(Remark~\ref{rem:lemmaB}).  In this section we \emph{design} $T$ and prove
that, for this $T$ and every finite block set, every first row in a large
class $\mathcal R_0$ is recoverable; consequently density for this $T$ is
equivalent to a single statement about approximating mates by mates at
first rows in $\mathcal R_0$ (Lemma~Z, Problem~\ref{prob:lemmaZ}).  The
mechanism is \emph{signature pinning}: every block vector carries a private
signature on its own set of coordinates, so that at a first row whose
signature sets keep some room, the block coefficients of any two admissible
decompositions of a mate are fixed by the mate itself, up to $O(t)$.

Throughout this section $N\ge1$ is fixed, $I=\{1,\dots,N\}$ and $p=p_N$.
The symbols $S_l,h_l,\delta_l,c_l,n_l,y_l,u_l$ below are local to this
section and are unrelated to the construction of
Section~\ref{sec:resonance}.

\subsection{The signature-ladder design}

\begin{definition}[Signature-ladder design]\label{def:SLD}
\textup{(D0)} Fix a bijection $\iota:\N\times\N\to\N$ with
$\iota(k,m)<\iota(k',m)$ for $k<k'$, and write $(k(l),m(l)):=\iota^{-1}(l)$.
Fix a coordinate $j_0$, pairwise disjoint infinite sets
$S_l\subseteq\N\setminus\{j_0\}$ $(l\in\N)$, and put
\[
 h_l:=\sum_{s\in S_l}2^{-s}e_s^*\ (\ge0),\qquad \delta_l:=\min\big\{2^{-l},
 (4(1+\|U\|)\|h_l\|_1)^{-1}\big\},
\]
so that $q^*(\delta_lh_l)\le(1+\|U\|)\delta_l\|h_l\|_1\le\frac14$.  Fix
\emph{targets} $y^{(i)}\in c_{00}\cap S_{q^*}$ $(i\ge1)$, dense in
$S_{q^*}$, with $y^{(1)}:=e_{j_0}^*/q^*(e_{j_0}^*)$, and a sequence
$(i_r)_{r\ge1}$ in which every positive integer occurs infinitely often.

\textup{(D1)} Recursively in $l=1,2,\dots$: given $c_l\in(0,1]$ ($c_1:=1$),
call $y\in c_{00}$ \emph{allowed at} $l$ if
\begin{enumerate}
\item[(a)] $\supp y\cap S_{l'}=\emptyset$ for every $l'\ge l$, and
\item[(b)] $2c_l\le2^{-2s}c_{l'}\delta_{l'}$ for every $l'<l$ and every
$s\in\supp y\cap S_{l'}$.
\end{enumerate}
Let $y_l:=y^{(i_{k(l)})}$ if this target is allowed at $l$, and
$y_l:=y^{(1)}$ otherwise ($y^{(1)}$ is allowed at every $l$).  Put
\begin{gather*}
 n_l:=q^*(y_l+\delta_lh_l),\quad u_l:=(y_l+\delta_lh_l)/n_l,\quad
 \delta^\circ_l:=\delta_l\|h_l\|_1/n_l,\quad
 \Lambda^\circ(l):=\prod_{l''\le l}\Big(1+\frac2{\delta^\circ_{l''}}\Big),\\
 T_{\rm hi}(l):=\min\Big\{T_{\rm lo}(l-1),\frac{2^{-l^3}}{l\Lambda^\circ(l)}
 \Big\}\ (T_{\rm lo}(0):=1),\quad n^w_l:=\big\lceil l2^{l^3}\Lambda^\circ(l)
 \big\rceil,\\
 T_{\rm lo}(l):=2^{-n^w_l}T_{\rm hi}(l),\qquad
 c_{l+1}:=\min\{c_l/4,T_{\rm lo}(l)^3\}.
\end{gather*}

\textup{(D2)} $Te_{k,m}:=c_{\iota(k,m)}u_{\iota(k,m)}$, extended linearly and
continuously to $\ell_1(\N\times\N)$.
\end{definition}

The interval $\mathcal W(l):=[T_{\rm lo}(l),T_{\rm hi}(l)]$ is the
\emph{$l$-th window}; it contains the $n^w_l$ dyadic scales
$T_{\rm hi}(l)2^{1-i}$, $1\le i\le n^w_l$.

\begin{theorem}[The design is admissible]\label{thm:SLD}
The operator $T$ of Definition~\ref{def:SLD} is admissible.  For
$l=\iota(k,m)$ we have $u_{k,m}=u_l$, $\Phi_m(k)=2^{-m-k}c_l$ and
$\lambda_l:=\lambda_{k,m}\le c_l/4$, and:
\begin{enumerate}
\item[(P1)] $y_l\in c_{00}$, $\|y_l\|_1\le1$, $n_l\in[\frac34,\frac54]$ and
$\supp y_l\cap S_{l'}=\emptyset$ for $l'\ge l$.  Consequently, for
$s\in S_l$: $u_l(s)=\delta_l2^{-s}/n_l$, $u_{l'}(s)=0$ for $l'<l$, and
$u_{l'}(s)=y_{l'}(s)/n_{l'}$ for $l'>l$.
\item[(P2)] $\sum_{l'>l}c_{l'}\le2c_{l+1}\le2T_{\rm lo}(l)^3$ and
$\sum_{l'>l}\lambda_{l'}\le T_{\rm lo}(l)^3$.
\item[(P3)] $T_{\rm hi}(l)2^{l^3}\Lambda^\circ(l)\le1/l$,
$n^w_l\ge l2^{l^3}\Lambda^\circ(l)$, and $T_{\rm hi}(l+1)\le T_{\rm lo}(l)$.
\end{enumerate}
\end{theorem}

\begin{proof}
(P1): $q^*(y_l)=1$, $\|y_l\|_1\le q^*(y_l)$ and $q^*(\delta_lh_l)\le\frac14$
give the bounds; the support statement is (a); for $s\in S_l$ the
signatures $h_{l'}$, $l'\ne l$, vanish at $s$ (disjointness), and the targets
$y_{l'}$ with $l'\le l$ vanish on $S_l$ by (a).  (P2): $c_{l'+1}\le
c_{l'}/4$ and $\lambda_{l'}\le c_{l'}/4$.  (P3) is immediate from (D1).

\emph{(T-a).} $q^*(Te_{k,m})=c_l\le1=c_1$.

\emph{(T-d).}  Fix $m$ and $i$.  The finite set $\supp y^{(i)}$ meets only
finitely many $S_{l'}$, and $c_l\to0$, so $y^{(i)}$ is allowed at every
large $l$.  Since $i$ occurs infinitely often in $(i_r)$, $y_l=y^{(i)}$ for
infinitely many $l=\iota(k,m)$, and along these $\delta_l\to0$ and
$n_l\to1$, so $q^*(u_l-y^{(i)})\to0$.  Hence every $y^{(i)}$ is a limit of
$u_{k,m}$, $k\to\infty$, and $\{u_{k,m}:k\in\N\}$ is dense in $S_{q^*}$.

\emph{(T-b), (T-c).}  Let $x\in\ell_1(\N\times\N)$, indexed by $l$, with
$\mathcal Z:=T x=\sum_lx_lc_lu_l\in c_{00}$, and suppose $x_{l'}\ne0$.  For
$s\in S_{l'}$, by (P1),
\[
 \mathcal Z(s)=\frac{x_{l'}c_{l'}\delta_{l'}2^{-s}}{n_{l'}}+\sum_{l>l',\
 s\in\supp y_l}\frac{x_lc_ly_l(s)}{n_l}.
\]
If the last sum is not empty, let $L(s)$ be its least index; by (b) at
$L(s)$, $2c_{L(s)}\le2^{-2s}c_{l'}\delta_{l'}$, and the sum is at most
$\frac43\|x\|_\infty\sum_{l\ge L(s)}c_l\le\frac{16}9\|x\|_\infty c_{L(s)}\le
\frac89\|x\|_\infty2^{-2s}c_{l'}\delta_{l'}$.  Hence $|\mathcal Z(s)|\ge
c_{l'}\delta_{l'}2^{-s}\big(|x_{l'}|/n_{l'}-\frac89\|x\|_\infty2^{-s}\big)>0$
for all large $s\in S_{l'}$, contradicting $\mathcal Z\in c_{00}$.  So
$x=0$: $T$ is injective and $T(\ell_1(\N\times\N))\cap c_{00}=\{0\}$.
\end{proof}

Only (T-a)--(T-d), (P1), (P2) and (P3) are used below.  No rate of
approximation of the targets and no independence property of the targets
is used; the same target sequence serves all blocks, so Mart\'in's
cross-block near duplicates are compatible with the design.

\subsection{The class $\mathcal R_0$}
Let $\mathcal L_N:=\{l:m(l)\le N\}$, the ladder indices of the blocks
present in $p_N$.  For $f\in S_{p^*}$ and $\gamma\in(0,1)$ let
$J_\gamma:=\{j\notin F:|z_j|\le1-\gamma\}$ (coordinates with room
$\ge\gamma$).

\begin{definition}[Signature room]\label{def:R0}
$f\in S_{p^*}$ belongs to $\mathcal R_0$ if $F$ is finite and
\begin{enumerate}
\item[(SR)] there are $\gamma\in(0,1)$ and $\vartheta_0\in(0,1]$ with
$\|h_l1_{S_l\cap J_\gamma}\|_1\ge\vartheta_0^l\|h_l\|_1$ for every
$l\in\mathcal L_N$.
\end{enumerate}
For such $f$ put $\delta'_l:=\delta_l\|h_l1_{S_l\cap J_\gamma}\|_1/n_l\ge
\vartheta_0^l\delta^\circ_l$, $\Lambda_f(l):=\prod_{l''\le l,\,l''\in
\mathcal L_N}(1+2/\delta'_{l''})$.
\end{definition}

Since $1+2/\delta'_l\le\vartheta_0^{-l}(1+2/\delta^\circ_l)$ and
$\sum_{l''\le l}l''\le l^2$, we have $\Lambda_f(l)\le\vartheta_0^{-l^2}
\Lambda^\circ(l)$.

\begin{lemma}\label{lem:R0}
\textup{(a)} Every $f'\in\NA((X,p),\R)\cap S_{p^*}$ belongs to
$\mathcal R_0$.  \textup{(b)} Every \emph{base-tame} $f$ \textup{(}$F$ and
$K$ finite and $\sup_{j\notin F\cup K}|z_j|<1$\textup{)} belongs to
$\mathcal R_0$.  \textup{(c)} $\mathcal R_0$ contains first rows with
infinite contact sets and with near-contacts, with arbitrary block
structure.
\end{lemma}

\begin{proof}
(a) By Proposition~\ref{prop:smooth}(c), $F'=\supp a'$ is finite and
$z'\in c_0$, so $\{j\notin F':|z'_j|>\frac12\}$ is finite.  Choose
$\gamma\in(0,\frac12]$ with $\gamma\le1-|z'_j|$ for each of these $j$ with
$|z'_j|<1$.  Then $\N\setminus J_\gamma=F'\cup K'$ is finite.  Every $S_l$
is infinite, so $\|h_l1_{S_l\cap J_\gamma}\|_1>0$ for all $l$, and the
ratio is $1$ for all but finitely many $l$; take $\vartheta_0$ as the
minimum of $1$ and of the $l$-th roots of the finitely many other ratios.
(b) The same argument with $\gamma:=1-\sup_{j\notin F\cup K}|z_j|$.
(c) Since $\sum_{s>s_0}2^{-s}=2^{-s_0}$, the least element of $S_l$ carries
at least half of the mass of $h_l$.  Hence every $f$ with $F$ finite and
disjoint from $\bigcup_lS_l$ (e.g.\ $F=\{j_0\}$),
$|z_j|\le\frac12$ at $j=\min S_l$ for all $l$, and arbitrary $z$
elsewhere (for instance $|z_j|=1$ on the infinite set
$\bigcup_l(S_l\setminus\{\min S_l\})$) satisfies (SR) with $\gamma=\frac12$
and $\vartheta_0=\frac12$.  First rows with prescribed $(a,z)$ exist by
Remark~\ref{rem:lemmaZ}(c).
\end{proof}

The first row of Section~\ref{sec:resonance}, transplanted to this design
(its contact set being a signature set), is the typical first row outside
$\mathcal R_0$ (\emph{signature-resonant}).

\subsection{Two-sided decompositions (any admissible $T$)}
In this subsection $T$ is an arbitrary admissible operator, $I$ is finite,
$f\in S_{p^*}$ and $g\in\Cset(f)$.  For $\Theta\in\ell_\infty$ put
$H_m(\Theta):=\|P_m^\perp D_m\Theta\|_2^2/C_m$ (this agrees with
Definition~\ref{def:certificate} on vectors vanishing on $P_m$, and
$H_m(\Theta+cw_m)=H_m(\Theta)$), and for $B\in\ell_1$,
$\Theta=(\Theta_m)\in\prod_m\ell_\infty$ put $\Gamma_w(B,\Theta):=q_0h(B)+
\sum_m\sigma_mH_m(\Theta_m)$.  Then $\sqrt{\Gamma_w}$ is a seminorm.

\begin{lemma}[Two-sided decompositions]\label{lem:twosided}
For every $t>0$ there are $B_\pm\in\ell_1$ and $\Theta_\pm\in\prod_m
\ell_\infty$ with $g=B_++L^*\Theta_+=B_-+L^*\Theta_-$ and
\[
 q^*(a+tB_+),\ N_m(w_m+t\Theta_{+,m}),\ q^*(a-tB_-),\ N_m(w_m-t\Theta_{-,m})
 \ \le s(t)\qquad(m\in I).
\]
\end{lemma}

\begin{proof}
By Lemma~\ref{lem:dualball}, $f+tg=A+L^*W$ with $\max(q^*(A),\|W\|_{V^*})=
p^*(f+tg)\le s(t)$; put $B_+:=(A-a)/t$ and $\Theta_+:=(W-w)/t$.  The same
with $-t$ gives $B_-,\Theta_-$.
\end{proof}

We call $(B_\pm,\Theta_\pm)$ a \emph{two-sided decomposition at scale $t$}.

\begin{lemma}[Uniform smallness]\label{lem:smallness}
For every $\eta>0$ there is $t_\eta>0$ such that every two-sided
decomposition at a scale $t\in(0,t_\eta]$ satisfies, for both signs and
every $m$, $\|tB_\pm\|_1\le\eta$, $\|D_m(t\Theta_{\pm,m})\|_2\le\eta$ and
$\big|\|w_m\pm t\Theta_{\pm,m}\|_\infty-M_m\big|\le\eta$.
\end{lemma}

\begin{proof}
Otherwise there are $t_n\to0$ and decompositions violating one of the
inequalities, say on the $+$ side.  Put $A_n:=a+t_nB_{+,n}$ and
$W_n:=w+t_n\Theta_{+,n}$; they are bounded, so after passing to a
subsequence $A_n\to A$ weak$^*$ in $\ell_1$ and $W_{n,m}\to W_m$ weak$^*$ in
$\ell_\infty$.  As $L^*$ is weak$^*$-continuous, $A+L^*W=f$ with
$q^*(A)\le1$ and $\|W\|\le1$, so $A=a$ and $W=w$ by
Proposition~\ref{prop:forced}(b).  Since $U$ is compact, $\|U^*A_n\|\to\nu$,
hence $\limsup\|A_n\|_1\le\lim(s(t_n)-\|U^*A_n\|)=\|a\|_1$; for bounded
coordinatewise convergent sequences this gives $\|A_n-a\|_1\to0$.  $D_m$ is
weak$^*$-to-norm continuous on bounded sets, so $\|D_m(W_{n,m}-w_m)\|_2\to0$
and $\|D_mW_{n,m}\|_2\to C_m$; then $\limsup\|W_{n,m}\|_\infty\le
1-C_m=M_m\le\liminf\|W_{n,m}\|_\infty$.  This contradicts the choice of the
decompositions.
\end{proof}

\begin{lemma}[First-order terms, budget and one-sided bounds]\label{lem:budget}
Let $\eta\le\eta_*:=\frac12\min\big(\nu/(2\|U\|+2),\min_mC_m,\frac14\big)$
and $0<t\le\min(t_\eta,1)$, and let $(B_\pm,\Theta_\pm)$ be a two-sided
decomposition at scale $t$.  Then for both signs:
\begin{enumerate}
\item[(a)] $|q_0B_\pm(\zh)|\le t/2$ and $|\ip{\Theta_{\pm,m}}{\zeta_m}|\le
t/2$ for every $m$;
\item[(b)] $\sum_{j\notin F}(|B_+(j)|-z_jB_+(j))\le t/(2q_0)$ and
$\sum_{j\notin F}(|B_-(j)|+z_jB_-(j))\le t/(2q_0)$; consequently
$\sum_{j\in J_\gamma}|B_\pm(j)|\le t/(2\gamma q_0)$ for $\gamma\in(0,1)$, and
$\|B_+1_{F^c}\|_1+\|B_-1_{F^c}\|_1\le\|(B_+-B_-)1_{F^c}\|_1+t/q_0$;
\item[(c)] $\sum_{k\in P_m}|\alpha_m(k)|\big(\|W\|_\infty-\sgn(w_m(k))W(k)
\big)\le t^2/(2\sigma_m)$ for $W=w_m+t\Theta_{+,m}$ and for
$W=w_m-t\Theta_{-,m}$;
\item[(d)] $\Gamma_w(B_\pm,\Theta_\pm)\le1+\eta_\Gamma(\eta)$, where
$\eta_\Gamma(\eta):=(1+\|U\|\eta/\nu)\max_m(1+\eta/C_m)-1\to0$ as $\eta\to0$.
\end{enumerate}
\end{lemma}

\begin{proof}
Consider the $+$ side; the $-$ side is the same with $-t$, $-B_-$,
$-\Theta_-$.  Let $A:=a+tB_+$, $W_m:=w_m+t\Theta_{+,m}$.  By
Lemma~\ref{lem:bookkeeping}(a) with $h=tg$ and $g(\xi)=0$,
\begin{equation}\label{eq:budget}
 q_0G_b(A)+\sum_m\sigma_mG_m(W_m)=q_0q^*(A)+\sum_m\sigma_mN_m(W_m)-1\le
 s(t)-1\le\frac{t^2}2 .
\end{equation}
(a) $q^*(A)\ge A(\zh)$ gives $tB_+(\zh)\le s(t)-1\le t^2/2$, and
$N_m(W_m)\ge\ip{W_m}{\zeta_m}/\sigma_m$ gives $\ip{\Theta_{+,m}}{\zeta_m}\le
\sigma_mt/2$.  Since $q_0B_+(\zh)+\sum_m\ip{\Theta_{+,m}}{\zeta_m}=g(\xi)=0$
and $q_0+\sum_m\sigma_m=1$, each of these terms is also $\ge-t/2$.
(b) By Lemma~\ref{lem:bookkeeping}(b), $G_b(A)\ge\sum_{j\notin F}(|tB_+(j)|-
z_jtB_+(j))$; use \eqref{eq:budget}.  On $J_\gamma$, $|x|\mp z_jx\ge\gamma|x|$.
The last inequality follows by summing over $j\notin F$ the inequality
$|x|+|y|\le|x-y|+(|x|-zx)+(|y|+zy)$ ($|z|\le1$), whose two sides differ by
$|x-y|-z(x-y)\ge0$.
(c) Lemma~\ref{lem:bookkeeping}(c) and \eqref{eq:budget}.
(d) By Lemma~\ref{lem:base}, $G_b(A)\ge\nu\Psi(tU^*B_+/\nu)\ge\frac{t^2}2
h(B_+)/(1+\|U\|\eta/\nu)$, since $t\|U^*B_+\|\le\|U\|\eta\le\nu/2$.  By
Lemma~\ref{lem:bookkeeping}(c) and Lemma~\ref{lem:smallness},
\begin{align*}
 G_m(W_m)&\ge\|D_mW_m\|-\frac{\ip{D_mW_m}{D_mw_m}}{C_m}=
 \frac{\|P^\perp_mD_mW_m\|^2}{\|D_mW_m\|+\ip{D_mW_m}{D_mw_m}/C_m}\\
 &\ge\frac{t^2}2\,\frac{H_m(\Theta_{+,m})C_m}{C_m+\eta}.
\end{align*}
Insert both bounds in \eqref{eq:budget}.
\end{proof}

\begin{lemma}[Sup-level parametrization]\label{lem:suplevel}
In the setting of Lemma~\ref{lem:budget}, fix $m$ and drop it.  Define
$d_\pm$ by $\|w+t\Theta_+\|_\infty=(1-d_+t)M$ and $\|w-t\Theta_-\|_\infty=
(1+d_-t)M$, and put $\omega_\pm:=\Theta_\pm+d_\pm w$.  Then:
\begin{enumerate}
\item[(a)] $|d_\pm|t\le\eta/M\le\frac12$;
\item[(b)] \textup{(box)} $|(1-d_+t)w(k)+t\omega_+(k)|\le(1-d_+t)M$ and
$|(1+d_-t)w(k)-t\omega_-(k)|\le(1+d_-t)M$ for every $k$;
\item[(c)] \textup{(peaks)} for $k\in P$ and $\varsigma_k:=\sgn w(k)$,
$\varsigma_k\omega_+(k)\le0\le\varsigma_k\omega_-(k)$, and
$\sum_{k\in P}|\alpha(k)||\omega_\pm(k)|\le t/(2\sigma)$;
\item[(d)] $|d(\omega_\pm)-d_\pm|\le t/\sigma$, where
$d(\omega):=\ip{Dw}{D\omega}/C$;
\item[(e)] $|t\Theta_\pm(k)|\le3$ for all $k$, and $|d_\pm|M\le3\|\Phi\|_2/t+
t/\sigma$;
\item[(f)] \textup{(non-peaks)} for $k\in Q$ and $\varsigma_k:=\sgn w(k)$ (any
sign if $w(k)=0$), $\varsigma_k\omega_+(k)\le(1-d_+t)\gap(k)/t$ and
$\varsigma_k\omega_-(k)\ge-(1+d_-t)\gap(k)/t$.
\end{enumerate}
\end{lemma}

\begin{proof}
(a) is Lemma~\ref{lem:smallness} ($M\ge\frac12$, $\eta\le\frac18$).  (b):
$w+t\Theta_+=(1-d_+t)w+t\omega_+$ and $w-t\Theta_-=(1+d_-t)w-t\omega_-$.
(c): at a peak, $|(1-d_+t)\varsigma_kM+t\omega_+(k)|\le(1-d_+t)M$ forces
$\varsigma_kt\omega_+(k)\le0$, and then $\|w+t\Theta_+\|_\infty-\varsigma_k
(w+t\Theta_+)(k)=t|\omega_+(k)|$; insert in Lemma~\ref{lem:budget}(c).  The
$-$ side is symmetric.  (d): by Lemma~\ref{lem:threshold},
$\ip{\Theta_+}{\zeta}/\sigma=\ip{\omega_+-d_+w}{\alpha}+\ip{D(\omega_+-d_+w)}
{Dw}/C=-\sum_P|\alpha(k)||\omega_+(k)|+d(\omega_+)-d_+$, using (c) and
$M+C=1$; similarly $\ip{\Theta_-}{\zeta}/\sigma=\sum_P|\alpha(k)||
\omega_-(k)|+d(\omega_-)-d_-$.  Use Lemma~\ref{lem:budget}(a) and (c).
(e): $|t\Theta_\pm(k)|\le N(w\pm t\Theta_\pm)+|w(k)|\le s(1)+1$; and
$d(\omega_\pm)-d_\pm=\ip{Dw}{D\Theta_\pm}/C-d_\pm M$, so $|d_\pm|M\le
\|D\Theta_\pm\|_2+t/\sigma$.  (f): multiply the box inequality (b) by
$\varsigma_k$ and use $|w(k)|=M-\gap(k)$.
\end{proof}

\begin{lemma}[Base part on a finite support]\label{lem:finitebase}
If $F$ is finite, there is $C_F<\infty$ with $\|\beta\|_1\le C_F(\|P^\perp
U^*\beta\|+|\beta(\zh)|)$ for every $\beta$ supported in $F$.
\end{lemma}

\begin{proof}
On the finite-dimensional space $\ell_1(F)$ the map $\beta\mapsto
(P^\perp U^*\beta,\beta(\zh))$ is injective: if $P^\perp U^*\beta=0$, then
$U^*\beta=\mu U^*a/\nu$, so $\beta=(\mu/\nu)a$ ($U^*$ injective) and
$\beta(\zh)=\mu/\nu$.
\end{proof}

\subsection{Signature pinning}
From now on $T$ is the operator of Definition~\ref{def:SLD}, $f\in
\mathcal R_0$ with constants $\gamma,\vartheta_0$, $g\in\Cset(f)$,
$\eta\le\eta_*$ with $\eta_\Gamma(\eta)\le1$, $0<t\le\min(t_\eta,1)$, and
$(B_\pm,\Theta_\pm)$ is a two-sided decomposition of $g$ at scale $t$.  For
$l\in\mathcal L_N$ put $c^\pm_l:=\lambda_l\Theta_{\pm,m(l)}(k(l))$ (and
$c^\pm_l:=0$ for $l\notin\mathcal L_N$), so that $L^*\Theta_\pm=
\sum_lc^\pm_lu_l$, and
\[
 \Delta c_l:=c^+_l-c^-_l,\qquad \Delta B:=B_+-B_-=-\sum_l\Delta c_lu_l .
\]
For a \emph{window index} $l_*$ the indices $l\le l_*$ are \emph{coarse}
and the indices $l>l_*$ are \emph{fine}.

\begin{lemma}[Box bound]\label{lem:box}
$|c^\pm_l|\le3\lambda_l/t$ for every $l$.
\end{lemma}

\begin{proof}
Lemma~\ref{lem:suplevel}(e).
\end{proof}

\begin{lemma}[Pinning inequality]\label{lem:pinning}
For every $l\in\mathcal L_N$,
\[
 \delta'_l|\Delta c_l|\le E_l+\sum_{l'>l}\kappa_{l',l}|\Delta c_{l'}|,\qquad
 E_l:=\|\Delta B1_{S_l\cap J_\gamma}\|_1,\quad\kappa_{l',l}:=
 \|y_{l'}1_{S_l}\|_1/n_{l'},
\]
and $\sum_l\kappa_{l',l}\le\frac43$ for every $l'$.
\end{lemma}

\begin{proof}
By (P1), for $s\in S_l$,
$-\Delta B(s)=\Delta c_l\delta_l2^{-s}/n_l+\sum_{l'>l}\Delta c_{l'}y_{l'}(s)/
n_{l'}$.  Take the $\ell_1$-norm over $s\in S_l\cap J_\gamma$: the first term
has norm $|\Delta c_l|\delta'_l$.  The sets $S_l$ are disjoint and
$\|y_{l'}\|_1/n_{l'}\le\frac43$.
\end{proof}

\begin{lemma}[Triangular solution]\label{lem:triangular}
For every $l_*$,
$\displaystyle\sum_{l\le l_*}|\Delta c_l|\le\Lambda_f(l_*)\Big(\sum_{l\le
l_*}E_l+\frac43\sum_{l>l_*}|\Delta c_l|\Big)$.
\end{lemma}

\begin{proof}
For coarse $l\in\mathcal L_N$ put $D_l:=|\Delta c_l|$, $X_l:=E_l+
\sum_{l'>l_*}\kappa_{l',l}D_{l'}$ and $\Sigma_l:=\sum_{l\le l'\le l_*}D_{l'}$
($\Sigma_{l_*+1}:=0$; $D_l:=0$ for $l\notin\mathcal L_N$).  By
Lemma~\ref{lem:pinning}, $D_l\le(X_l+\frac43\Sigma_{l+1})/\delta'_l$, so
$\Sigma_l\le X_l/\delta'_l+(1+\frac4{3\delta'_l})\Sigma_{l+1}$.  Unrolling
from $l_*$ downwards, $\Sigma_1\le\sum_{l\le l_*}(X_l/\delta'_l)\prod_{l''<l}
(1+\frac4{3\delta'_{l''}})\le\Lambda_f(l_*)\sum_{l\le l_*}X_l$ (products over
$\mathcal L_N$), and $\sum_{l\le l_*}X_l\le\sum_{l\le l_*}E_l+\frac43
\sum_{l'>l_*}D_{l'}$ by Lemma~\ref{lem:pinning}.
\end{proof}

\begin{proposition}[Pinning on a window]\label{prop:pinned}
If $t\in\mathcal W(l_*)$, then with $K:=K_1\Lambda_f(l_*)$,
$K_1:=1/(\gamma q_0)+14$, and $K_d:=\max_m(1/(mC_m)+2/(\sigma_mM_m))$:
\begin{enumerate}
\item[(a)] $\sum_l|\Delta c_l|\le Kt$ and $\|\Delta B\|_1\le Kt$;
\item[(b)] $\|B_+1_{F^c}\|_1+\|B_-1_{F^c}\|_1\le Kt+t/q_0$; in particular the
total base mass of both decompositions on contacts and near-contacts is
$O(Kt)$;
\item[(c)] $|d_{+,m}-d_{-,m}|\le K_dKt$ for every $m$, with $d_{\pm,m}$ as in
Lemma~\ref{lem:suplevel}.
\end{enumerate}
\end{proposition}

\begin{proof}
(a) By Lemma~\ref{lem:budget}(b) and disjointness of the $S_l$,
$\sum_lE_l\le\sum_{j\in J_\gamma}(|B_+(j)|+|B_-(j)|)\le t/(\gamma q_0)$.  By
Lemma~\ref{lem:box} and (P2), $\sum_{l>l_*}|\Delta c_l|\le6\sum_{l>l_*}
\lambda_l/t\le6T_{\rm lo}(l_*)^3/t\le6t^2$.  By Lemma~\ref{lem:triangular},
$\sum_l|\Delta c_l|\le\Lambda_f(l_*)(t/(\gamma q_0)+8t^2)+6t^2\le Kt$.  And
$\|u_l\|_1\le q^*(u_l)=1$.  (b) Lemma~\ref{lem:budget}(b) and (a).
(c) Drop $m$; put $\Delta\Theta:=\Theta_+-\Theta_-$ and $\Delta d:=d_+-d_-$.
By Lemma~\ref{lem:suplevel}(d), $\Delta d=d(\omega_+-\omega_-)+r$ with
$|r|\le2t/\sigma$, and $\omega_+-\omega_-=\Delta\Theta+\Delta d\,w$, so
$\Delta d\,M=\ip{Dw}{D\Delta\Theta}/C+r$.  Since $\lambda_k=m\Phi_m(k)$,
$\ip{Dw}{D\Delta\Theta}=\frac1m\sum_k\Phi_m(k)w(k)\Delta c_{\iota(k,m)}$,
whose modulus is at most $\frac Mm\sum_l|\Delta c_l|$.  Divide by $M$ and
use (a) and $K\ge1$.
\end{proof}

Thus every one-sided resource used by the two decompositions (contacts and
near-contacts on the base side; on the block side peaks, near-threshold
coordinates and coordinates pushed beyond their gap, as the next proof
shows) is controlled by the two-sided differences $\Delta c_l$, $\Delta B$,
which are pinned by the signatures.  At the first row of
Section~\ref{sec:resonance} this fails: there the signature set of the
resonant carrier is the contact set.

\subsection{The window certificate}

\begin{definition}[Window certificate]\label{def:windowcert}
For $t\in\mathcal W(l_*)$ let $\kappa_t:=(B_+1_F)(\zh)$,
$b_t:=B_+1_F-\kappa_ta$, and in block $m$ let
\[
 \omega^c_m(k):=\max\Big\{-\frac{2\gap_m(k)}t,\min\Big\{\omega_{+,m}(k),
 \frac{2\gap_m(k)}t\Big\}\Big\}
\]
if $k\in Q_m$, $\iota(k,m)\le l_*$ and $\gap_m(k)\ge t^2$, and
$\omega^c_m(k):=0$ otherwise.  Put $c_t:=(b_t,(\omega^c_m)_m)$.
\end{definition}

\begin{proposition}[Window certificate]\label{prop:windowcert}
Let $t\in\mathcal W(l_*)$ with $Kt\le1$.  There are constants $K_b,K_2,K_4$,
depending only on $f$, $\gamma$ and the design, such that:
\begin{enumerate}
\item[(a)] $c_t$ is a balanced finite certificate (Definition~\ref{def:tame});
on $\supp\omega^c_m$, $\gap_m(k)\ge t^2$ and $|\omega^c_m(k)|\le
2\gap_m(k)/t$; $\|D_m\omega^c_m\|_2\le1/t$ and $|d_m(\omega^c_m)|\le1/t$;
\item[(b)] $\|b_t\|_1\le K_b$;
\item[(c)] $\|g-g_{c_t}\|_1\le K_2\Lambda_f(l_*)t$;
\item[(d)] $\sqrt{\Gamma_w(c_t)}\le\sqrt{1+\eta_\Gamma(\eta)}+K_4\Lambda_f(l_*)t$.
\end{enumerate}
\end{proposition}

\begin{proof}
(a) $b_t(\zh)=\kappa_t-\kappa_ta(\zh)=0$; the coarse set is finite; and
$\|D_m\omega^c_m\|_2\le(2/t)\|\Phi_m\|_2\le1/t$, $|d_m(\omega)|\le
\|D_m\omega\|_2$.

(b) $|\kappa_t|\le|B_+(\zh)|+\|\zh\|_\infty\|B_+1_{F^c}\|_1\le t/(2q_0)+
(1+\|U\|)(Kt+t/q_0)$ by Lemma~\ref{lem:budget}(a) and
Proposition~\ref{prop:pinned}(b).  By Lemma~\ref{lem:finitebase},
$\|B_+1_F\|_1\le C_F(\|P^\perp U^*B_+\|+\|U\|\|B_+1_{F^c}\|_1+|\kappa_t|)$,
and $\|P^\perp U^*B_+\|^2=\nu h(B_+)\le2\nu/q_0$ by
Lemma~\ref{lem:budget}(d).  Use $Kt\le1$ and $t\le1$.

(c) $g-g_{c_t}=(B_+-b_t)+\sum_mR_m^*X_m$ with $X_m:=\Theta_{+,m}-\omega^c_m+
d_m(\omega^c_m)w_m$, and $\|R^*_mX\|_1\le\sum_k\lambda_{k,m}|X(k)|$.  First,
$B_+-b_t=B_+1_{F^c}+\kappa_ta$ has norm $\le(2+\|U\|)(K+2/q_0)t$.  Fix $m$
and drop it; with $\Theta_+=\omega_+-d_+w$,
\[
 X=\big[(\omega_+-\omega^c)-(d_+-d(\omega^c))w\big]1_{\rm coarse}+
 \big[\Theta_++d(\omega^c)w\big]1_{\rm fine}.
\]
\emph{Fine part.}  $\sum_{\rm fine}\lambda_k|\Theta_+(k)|=\sum_{l>l_*}
|c^+_l|\le3t^2$ and $|d(\omega^c)|\sum_{\rm fine}\lambda_k\le t^2$, by
Lemma~\ref{lem:box}, (P2) and $t\ge T_{\rm lo}(l_*)$.

\emph{Coarse part.}  Let $\rho_k:=|\omega_+(k)-\omega^c(k)|$.  We claim
$\lambda_k\rho_k\le|\Delta c_k|+\lambda_k|\Delta d|M+2t\lambda_k$, where
$\Delta c_k:=\lambda_k(\Theta_+(k)-\Theta_-(k))$, so that
$\lambda_k|\omega_+(k)-\omega_-(k)|\le|\Delta c_k|+\lambda_k|\Delta d|M$.
At a peak $\omega^c(k)=0$, and $\omega_+(k),\omega_-(k)$ have opposite signs
in the weak sense (Lemma~\ref{lem:suplevel}(c)), so $\rho_k\le|\omega_+(k)-
\omega_-(k)|$.  At $k\in Q$ with $|\omega_+(k)|\le2\gap(k)/t$: $\rho_k=0$ if
$\gap(k)\ge t^2$, and $\rho_k\le2\gap(k)/t<2t$ otherwise.  At $k\in Q$ with
$|\omega_+(k)|>2\gap(k)/t$: by Lemma~\ref{lem:suplevel}(f) and $|d_\pm t|\le
\frac12$, $\varsigma_k\omega_+(k)\le\frac32\gap(k)/t$, hence (choosing
$\varsigma_k:=-\sgn\omega_+(k)$ if $w(k)=0$) $\varsigma_k\omega_+(k)<
-2\gap(k)/t$, and $e_k:=-\varsigma_k\omega_+(k)-2\gap(k)/t>0$; also
$\varsigma_k\omega_-(k)\ge-2\gap(k)/t$, so $e_k\le\varsigma_k(\omega_-(k)-
\omega_+(k))$: the part of $\omega_+(k)$ beyond the clamp is one-sided and
bounded by the two-sided difference.  Then $\rho_k=e_k$ if $\gap(k)\ge t^2$,
and $\rho_k=2\gap(k)/t+e_k\le2t+|\omega_+(k)-\omega_-(k)|$ otherwise.  This
proves the claim.  Summing, with $\sum_k\lambda_k\le1$ and
Proposition~\ref{prop:pinned}, $\sum_{\rm coarse}\lambda_k\rho_k\le
(1+K_d)Kt+2t$.  Next, $|d_+-d(\omega^c)|\le|d_+-d(\omega_+)|+\|D(\omega_+-
\omega^c)\|_2\le t/\sigma+\sum_k\Phi_k|\omega_+(k)-\omega^c(k)|$, and
\begin{align*}
 \sum_k\Phi_k|\omega_+(k)-\omega^c(k)|&\le\frac1m\sum_{\rm coarse}
 \lambda_k\rho_k+\sum_{\rm fine}\Phi_k(|\Theta_+(k)|+|d_+|M)\\
 &\le\frac1m\sum_{\rm coarse}\lambda_k\rho_k+3t^2+\Big(\frac3t+\frac
 t\sigma\Big)t^3,
\end{align*}
by Lemma~\ref{lem:suplevel}(e) and (P2).  Multiplying by
$\sum_{\rm coarse}\lambda_k|w(k)|\le1$ and collecting terms (with $t\le1$,
$K\ge1$, $\Lambda_f\ge1$) gives (c).

(d) The pair of $c_t$ is
\[
 \big(b_t,(\omega^c_m-d_m(\omega^c_m)w_m)_m\big)=(B_+,\Theta_+)-\big(B_+-b_t,
 (X_m)_m\big),
\]
and $H_m$ vanishes on multiples of $w_m$.  As $\sqrt{\Gamma_w}$ is a
seminorm,
\[
 \sqrt{\Gamma_w(c_t)}\le\sqrt{\Gamma_w(B_+,\Theta_+)}+\sqrt{\Gamma_w(B_+-
 b_t,X)}.
\]  The first term is at most
$\sqrt{1+\eta_\Gamma(\eta)}$ (Lemma~\ref{lem:budget}(d)).  For the second,
$q_0h(x)\le\|U\|^2\|x\|_1^2/\nu$ and
$\sigma_mH_m(X_m)\le\|D_mX_m\|_2^2/C_m\le\big(\sum_k\lambda_{k,m}|X_m(k)|
\big)^2/C_m$, and both were bounded in (c) by a
constant times $\Lambda_f(l_*)t$.
\end{proof}

\subsection{Uniform transfer expansion and windowed averaging}

\begin{lemma}[Uniform transfer expansion]\label{lem:uniformtransfer}
Let $T$ be admissible, $I$ finite and $f\in S_{p^*}$ with $F$ finite.  For
every $\varepsilon_{\rm tr}>0$ and $A_0\ge1$ there are $c_1\in(0,\frac18]$ and
$t_1>0$ such that for every $t\in(0,t_1]$ and every balanced finite
certificate $c=(b,\omega)$ with
\begin{enumerate}
\item[(C-a)] $\|b\|_1\le A_0$ and $\Gamma_w(c)\le2$, and
\item[(C-b)] $\supp\omega_m\subseteq\{k\in Q_m:\gap_m(k)\ge t^2\}$ and
$|\omega_m(k)|\le2\gap_m(k)/t$ for all $m,k$,
\end{enumerate}
we have $p^*(f+sg_c)\le1+\frac{s^2}2(\Gamma_w(c)+\varepsilon_{\rm tr})$ for
all $|s|\le c_1t$.
\end{lemma}

\begin{proof}
Fix $\gamma_m\in(M_m/2,M_m)$ with $|w_m(k)|\ne\gamma_m$ for all $k$.  Let
$|s|\le c_1t$, $t\le t_1$, and consider $f+sg_c=A+\sum_mR^*_mW_m$ with
$A:=a+sb$, $W_m:=w_m+s(\omega_m-d_mw_m)$, $d_m:=d_m(\omega_m)$; the
first-order terms vanish (Lemma~\ref{lem:algebra}, $b(\xi)=0$).
\emph{Base.}  If $c_1t_1A_0\le\min(a_{\min},\nu/(2\|U\|+2))$, there are no
flips on $F$ and $\supp b\subseteq F$, so by
Lemmas~\ref{lem:bookkeeping}(b) and~\ref{lem:base},
$G_b(A)\le\hat G_b:=\frac{s^2}2h(b)(1+K_A|s|)$, $K_A:=2\|U\|A_0/\nu$.
\emph{Blocks.}  $|d_m|\le\|D_m\omega_m\|_2\le(2/t)\|\Phi_m\|_2\le1/t$, so the
coordinatewise radius of $\omega_m$ is at least $\min(t/4,C_mt/2)\ge c_1t$
if $c_1\le\min_mC_m/2$, and Lemma~\ref{lem:block}(d) gives $G_m(W_m)\le\hat
G_m:=\frac{s^2}2H_m(\omega_m)(1+2c_1/C_{\min})$, $C_{\min}:=\min_mC_m$.
Hence $\hat\Gamma\le\frac{s^2}2\Gamma_w(c)(1+K_A|s|+2c_1/C_{\min})$.
\emph{Rebalancing.}  Since $h(b)\le2/q_0$ and $H_m(\omega_m)\le2/\sigma_m$,
$|\varepsilon_m|:=|\hat G_m-\hat\Gamma|/e_{y,m}\le K_3s^2$ with $K_3$
depending only on $f$ (as $e_{y,m}\ge\frac12$).  Choose
$\eta_1:=\min(\frac14,\varepsilon_{\rm tr}/(8|I|K_3))$ and transfer data
(Lemma~\ref{lem:transferdata}); $\|D_my_m\|_2\le K_y$ and $|Y_m|/q_0\le K_Y$
with $K_y,K_Y$ independent of $\eta_1$.  Now $\|W_m-w_m\|_\infty\le
|s|(2/t+1/t)\le3c_1$ and $\|D_m(W_m-w_m)\|_2\le2c_1$; so if $c_1\le
\min_m\min((M_m-\gamma_m)/12,C_m/4)$ and $t_1$ is so small that $K_3c_1^2t_1^2
(2+\Lambda^\varsigma_m)\le\gamma_m/2$ and $K_3c_1^2t_1^2\le(M_m-\gamma_m)/4$,
Lemma~\ref{lem:TV} gives \eqref{eq:R1}.  Proposition~\ref{prop:rebalancing}
yields
\[
 p^*(f+sg_c)\le1+\hat\Gamma+|I|K_3\eta_1s^2+2|I|K_3K_Y(1+\|U\|)A_0|s|^3+
 \frac{8K_3K_yc_1s^2+K_3^2K_y^2s^4}{C_{\min}} .
\]
Choose first $\eta_1$ as above, then $c_1$ with
$(4+16K_3K_y)c_1/C_{\min}\le\varepsilon_{\rm tr}/4$, then $t_1$ so small that
the terms of order $|s|^3,s^4$ (with $|s|\le c_1t_1$) and $\Gamma_w(c)K_A|s|$
contribute at most $\frac{s^2}2\cdot\frac{\varepsilon_{\rm tr}}4$, besides the
conditions above.  Since $\Gamma_w(c)\le2$, the bound is at most
$1+\frac{s^2}2(\Gamma_w(c)+\varepsilon_{\rm tr})$.
\end{proof}

\begin{theorem}[Windowed averaging]\label{thm:windowed}
Let $I$ be finite, $f\in S_{p^*}$ with $F$ finite, $g\in\Cset(f)$,
$\rho\in(0,1)$ and $\eta_0>0$ with $\rho^2(1+\eta_0)\le(1+\rho^2)/2$.
Suppose there are $c_1\in(0,1]$ and triples $(t^{(j)},n_j,K_j)$ with
$t^{(j)}\to0$, $n_j\ge24\rho^2K_j/(c_1(1-\rho^2))$ and
$K_jt^{(j)}/n_j\to0$, such that for every $j$ and every
$t\in\{t^{(j)}2^{1-i}:1\le i\le n_j\}$ there is a balanced finite
certificate $c_t$ at $f$ with
\begin{enumerate}
\item[(i)] $p^*(f+sg_{c_t})\le1+\frac{s^2}2(1+\eta_0)$ for $|s|\le c_1t$;
\item[(ii)] $p^*(g-g_{c_t})\le K_jt$;
\item[(iii)] $\Gamma_w(c_t)\le1+\eta_0$.
\end{enumerate}
Then $(f,\rho g)\in\cl\NA((c_0,p),\ell_2^2)$.  More precisely, the averaged
certificates $c$ below satisfy $\rho g_c\in\Cset(f)$, $\Gamma_w(\rho c)\le1$
and $\rho g_c\to\rho g$.
\end{theorem}

\begin{proof}
Fix $s_0\in(0,1]$ with $s_0^2\le1-\rho^2$, and $j$ so large that
$2\rho K_jt^{(j)}/n_j\le(1-\rho^2)s_0/3$.  Write $n:=n_j$, $K:=K_j$,
$t_i:=t^{(j)}2^{1-i}$ and $c:=\frac1n\sum_ic_{t_i}$, a balanced finite
certificate with $g_c=\frac1n\sum_ig_{c_{t_i}}$.  For each $i$ and real $s$:
(A) if $\rho|s|\le c_1t_i$, then $p^*(f+s\rho g_{c_{t_i}})\le1+
\frac{\rho^2s^2}2(1+\eta_0)$; (B) always $p^*(f+s\rho g_{c_{t_i}})\le
p^*(f+s\rho g)+\rho|s|Kt_i\le s(\rho s)+\rho|s|Kt_i$.  By convexity,
$p^*(f+s\rho g_c)\le\frac1n\sum_ip^*(f+s\rho g_{c_{t_i}})$.

\emph{Case $|s|\le s_0$.}  Let $I_s:=\{i:c_1t_i<\rho|s|\}$; as the $t_i$ are
dyadic, $\sum_{I_s}t_i<2\rho|s|/c_1$.  Using (A) off $I_s$ and (B) on $I_s$,
$p^*(f+s\rho g_c)\le\max\{1+\frac{\rho^2s^2}2(1+\eta_0),s(\rho s)\}+Qs^2$
with $Q:=2\rho^2K/(c_1n)\le(1-\rho^2)/12$.  Now
$1+\frac{\rho^2s^2}2(1+\eta_0)+Qs^2\le1+\frac{s^2}2(1-\frac{1-\rho^2}3)\le
1+\frac{s^2}2-\frac{s^4}8\le s(s)$ because $s^2\le1-\rho^2$, and
$s(\rho s)+Qs^2\le s(s)$ by Lemma~\ref{lem:slack}(a).

\emph{Case $|s|\ge s_0$.}  Using (B) for all $i$, $p^*(f+s\rho g_c)\le
s(\rho s)+2\rho Kt^{(j)}|s|/n\le s(\rho s)+(1-\rho^2)s_0|s|/3\le s(s)$, by
Lemma~\ref{lem:slack}(a) and $\min(s^2,|s|)\ge s_0|s|$.

Hence $\rho g_c\in\Cset(f)$.  By convexity of $\Gamma_w$ in the data,
$\Gamma_w(\rho c)\le\rho^2(1+\eta_0)\le1$, and $p^*(\rho g_c-\rho g)\le
2\rho Kt^{(j)}/n\to0$.  Since $a\in c_{00}$, Theorem~\ref{thm:transfer}
gives $(f,\rho'\rho g_c)\in\cl\NA$ for all $\rho'<1$, hence
$(f,\rho g_c)\in\cl\NA$; let $j\to\infty$.
\end{proof}

\subsection{Recovery on $\mathcal R_0$ and the reduction to Lemma Z}

\begin{theorem}[Recovery on $\mathcal R_0$]\label{thm:R0}
Let $T$ be the operator of Definition~\ref{def:SLD} and $N\ge1$.  Then
$\mathcal R_0\subseteq\Rec$: for every $f\in\mathcal R_0$ and every
$g\in\Cset(f)$, $(f,g)\in\cl\NA((c_0,p_N),\ell_2^2)$.  No condition on the
blocks of $f$ is needed \textup{(}infinitely many strict non-peaks,
near-threshold coordinates, weak or degenerate peaks, failure of
\textup{(MS)}, several active blocks and infinite contact sets are
allowed\textup{)}.
\end{theorem}

\begin{proof}
Fix $f\in\mathcal R_0$ (constants $\gamma,\vartheta_0$), $g\in\Cset(f)$ and
$\rho\in(0,1)$.  Choose $\eta_0>0$ with $\rho^2(1+\eta_0)\le(1+\rho^2)/2$ and
put
\[
 \kappa_0:=\tfrac12\big(\sqrt{1+\eta_0/2}-1\big),\qquad
 \varepsilon_{\rm tr}:=\eta_0/2,\qquad A_0:=K_b+1,
\]
with $K_b$ from Proposition~\ref{prop:windowcert}; let $c_1,t_1$ be given by
Lemma~\ref{lem:uniformtransfer}.  Choose $\eta\le\eta_*$ with
$\eta_\Gamma(\eta)\le1$ and $\sqrt{1+\eta_\Gamma(\eta)}\le1+\kappa_0$, and put
$t_\star:=\min(t_\eta,t_1,1)$.  For a window index $l$ put $t^{(l)}:=
T_{\rm hi}(l)$, $n_l:=n^w_l$ and $K_l:=(1+\|U\|)K_2\Lambda_f(l)$.  The scales
$t_i=T_{\rm hi}(l)2^{1-i}$, $1\le i\le n^w_l$, lie in $\mathcal W(l)$.  By
(P3) and $\Lambda_f(l)\le\vartheta_0^{-l^2}\Lambda^\circ(l)$,
\[
 \Lambda_f(l)T_{\rm hi}(l)\le\vartheta_0^{-l^2}2^{-l^3}/l\to0\qquad(l\to\infty).
\]
Hence for all large $l$ and every $t\in\mathcal W(l)$: $t\le t_\star$,
$K_1\Lambda_f(l)t\le1$ and $K_4\Lambda_f(l)t\le\kappa_0$.  Take a two-sided
decomposition at scale $t$ (Lemma~\ref{lem:twosided}) and its window
certificate $c_t$.  By Proposition~\ref{prop:windowcert},
$\|b_t\|_1\le A_0$, $\sqrt{\Gamma_w(c_t)}\le1+2\kappa_0=\sqrt{1+\eta_0/2}$,
i.e.\ $\Gamma_w(c_t)\le1+\eta_0/2\le2$, and (C-b) holds; by
Lemma~\ref{lem:uniformtransfer}, $p^*(f+sg_{c_t})\le1+\frac{s^2}2(1+\eta_0)$
for $|s|\le c_1t$; and $p^*(g-g_{c_t})\le q^*(g-g_{c_t})\le(1+\|U\|)
\|g-g_{c_t}\|_1\le K_lt$.  Finally, by (P3),
\[
 \frac{n^w_l}{K_l}\ge\frac{l2^{l^3}\vartheta_0^{l^2}}{(1+\|U\|)K_2}\to\infty,
 \qquad\frac{K_lt^{(l)}}{n^w_l}\le K_lT_{\rm hi}(l)\le(1+\|U\|)K_2
 \frac{\vartheta_0^{-l^2}2^{-l^3}}l\to0,
\]
so the hypotheses of Theorem~\ref{thm:windowed} hold along large $l$, and
$(f,\rho g)\in\cl\NA$.  As $\rho<1$ was arbitrary, $(f,g)\in\cl\NA$.
\end{proof}

The super-exponential factor $2^{l^3}$ of the design absorbs the
$f$-dependent constant $\vartheta_0^{-l^2}$ and every constant depending on
$f$, $g$ and $\rho$; this is the only place where the order of quantifiers
($T$ before $f$) matters.

\begin{corollary}[No intrinsic defect on $\mathcal R_0$]\label{cor:nodefectR0}
For the operator of Definition~\ref{def:SLD} and $f\in\mathcal R_0$,
$\Cset(f)$ is the norm closure of the set of $g_c\in\Cset(f)$ with $c$ a
balanced finite certificate and $\Gamma_w(c)\le1$.  In particular
$\Def(f)=\emptyset$.
\end{corollary}

\begin{proof}
By the proof of Theorem~\ref{thm:R0} and Theorem~\ref{thm:windowed}, every
$\rho g$ ($g\in\Cset(f)$, $\rho<1$) is a limit of such $g_c$; let
$\rho\to1$.  These $g_c$ lie in $\mathfrak B(f)\subseteq\mathcal K(f)$.
\end{proof}

\begin{theorem}[Reduction]\label{thm:reductionZ}
For the operator of Definition~\ref{def:SLD} and $N\ge1$, the set
$\NA((c_0,p_N),\ell_2^2)$ is dense if and only if the following statement
\textup{(Lemma Z)} holds: for every $f\in S_{p_N^*}$, $g\in\Cset(f)$,
$\rho<1$ and $\varepsilon>0$ there is $f'\in\mathcal R_0$ with
$p_N^*(f'-f)<\varepsilon$ and $\dist(\rho g,\Cset(f'))<\varepsilon$.
\end{theorem}

\begin{proof}
If Lemma~Z holds, take $f'\in\mathcal R_0$ and $g'\in\Cset(f')$ with
$p^*(f'-f)<\varepsilon$ and $p^*(g'-\rho g)<\varepsilon$.  By
Theorem~\ref{thm:R0}, $(f',g')\in\cl\NA$, and $\|(f',g')-(f,\rho g)\|<
2\varepsilon$.  So $(f,\rho g)\in\cl\NA$ for all $f,g,\rho$, and density
follows from Proposition~\ref{prop:reduction}(b).  Conversely, if
$\NA((c_0,p_N),\ell_2^2)$ is dense, Lemma~\ref{lem:pair} gives
norm-attaining $f_n\to f$ with $\rho g\in\Li_n\Cset(f_n)$, and
$f_n\in\mathcal R_0$ by Lemma~\ref{lem:R0}(a).
\end{proof}

By Lemma~\ref{lem:martintail}, Lemma~Z for infinitely many $N$ would give
the density of $\NA((c_0,p),\ell_2^2)$ for Mart\'in's norm $p$ built with
the designed operator.

\begin{problem}[Lemma Z]\label{prob:lemmaZ}
Let $T$ be the operator of Definition~\ref{def:SLD} and $N\ge1$.  Is it true
that for every $f\in S_{p_N^*}$, every $g\in\Cset(f)$, every $\rho<1$ and
every $\varepsilon>0$ there is $f'\in\mathcal R_0$ \textup{(}not necessarily
norm attaining\textup{)} with $p_N^*(f'-f)<\varepsilon$ and
$\dist_{p_N^*}(\rho g,\Cset(f'))<\varepsilon$?
\end{problem}

\begin{remark}[What is known about Lemma Z]\label{rem:lemmaZ}
(a) Lemma~Z holds at every $f\in\mathcal R_0$ (take $f'=f$), and for every
mate covered by Theorem~\ref{thm:engineered} or
Corollary~\ref{cor:BTrecovered} (the engineered approximants are norm
attaining, hence in $\mathcal R_0$); in particular at every \textup{(BT)}
point.

(b) $f\notin\mathcal R_0$ iff $F$ is infinite or, for all $\gamma,\vartheta_0$,
some signature set satisfies $\|h_l1_{S_l\cap J_\gamma}\|_1<\vartheta_0^l
\|h_l\|_1$: the one-sided base resources (support, contacts,
near-contacts) swallow the signature of some carrier, or of a fast sequence
of carriers (\emph{signature-resonant} first rows).  If the signature of the
carrier $l_0$ is swallowed ($\delta'_{l_0}=0$), then $\Delta c_{l_0}$ is not
pinned, and every coarser carrier $l<l_0$ whose signature set is met by the
target $y_{l_0}$ becomes slaved to the free coefficient $\Delta c_{l_0}$
through the pinning inequality; so more than the swallowed carriers may
lose their pinning.

(c) The approximants in Lemma~Z need not attain their norms.  Every pair
$(a',z')$ with $q^*(a')=1$, $z'\in B_{\ell_\infty}$ and $z'=\sgn a'$ on
$\supp a'$ is the forced data of some $f'\in S_{p^*}$: put
$e':=U^*a'/\|U^*a'\|$, $\zh':=z'+Ue'$, $q'_0:=(1+\sum_m|R^{**}_m\zh'|_m)^{-1}$,
$\xi':=q_0'\zh'$ and $f':=a'+L^*J_V(L^{**}\xi')$.  Indeed $q^{**}(\zh')=
a'(\zh')=1$, so $f'(\xi')=1=p^{**}(\xi')$ and $p^*(f')\le1$ by
Lemma~\ref{lem:dualball}; Proposition~\ref{prop:forced} identifies the
forced data.  If $a'\to a$ in $\ell_1$ and $z'\to z$ coordinatewise, then
$f'\to f$ (as in the proof of Proposition~\ref{prop:approximants}).  So one
may lower $|z|$ on far parts of the swallowed signature sets; whether the
resulting $f'$ satisfy the mate condition of Lemma~Z is exactly the open
question.

(d) \emph{Heuristic, not used.}  If only finitely many carriers are unpinned,
the arguments of this section suggest that every two-sided decomposition is a
pinned part of size $O(Kt)$, plus finitely many free (and slaved)
coefficients with their base counterparts, i.e.\ an exact two-piece
structure up to $O(Kt)$ on whole windows; combining windowed averaging with
the engineering of Section~\ref{sec:engineered} is a natural route to
Lemma~Z in that case.  Lowering $|z|$ on far parts of the swallowed sets
makes the corresponding pinning constants tiny; above that scale $f'$
copies the mate structure of $f$, below it the windows of $f'$ take over,
and the transition band would have to be handled by the $\rho$-slack.
None of this is proved.

(e) \emph{Sketch, not used.}  For $F$ infinite with (SR), the proof of
Theorem~\ref{thm:R0} appears to go through after clamping $B_+$ on $F$ at
$2|a_j|/t$ (the excess beyond the no-flip capacity is one-sided and bounded
by $|\Delta B(j)|$ plus flip costs), truncating the clamped base part to a
finite set, adapting Lemma~\ref{lem:uniformtransfer} to $|b_j|\le2|a_j|/t$,
and using a version of Theorem~\ref{thm:transfer} for $a\notin c_{00}$; we
have not written out the details.
\end{remark}

\section{Status of the density problem}\label{sec:status}

\subsection{What is proved}
Let $T$ be an arbitrary admissible operator.  Density of
$\NA((c_0,p),\ell_2^2)$ is equivalent to $\Rec=S_{p^*}$
(Proposition~\ref{prop:reduction}), and by Lemma~\ref{lem:martintail} it
suffices to prove $\Rec=S_{p_N^*}$ for infinitely many $N$.  A first row
$f$ belongs to $\Rec$ under each of the following sufficient conditions:
\begin{enumerate}
\item[(S1)] $f$ attains its norm (Proposition~\ref{prop:reduction}(d));
\item[(S2)] $f$ belongs to the residual set $\mathcal O$ of Hausdorff continuity
points of the fibres (residual uniform recovery theorem of \cite{PrA});
\item[(S3)] $\Cset(f)=\cl\Cert^{\rm sh}(f)$; then every mate is recovered
along every sequence $f_n\to f$ (Theorem~\ref{thm:shifted}); more generally
$\Cset(f)=\mathcal K(f)$ when $a\in c_{00}$ (Proposition~\ref{prop:intrinsic});
\item[(S4)] $I$ is finite and $f$ satisfies \textup{(BT)}: finite base support,
finitely many strict non-peaks, no degenerate peaks and \textup{(MS)}, with
an \emph{arbitrary} contact set (Corollary~\ref{cor:BTrecovered}; this
contains Theorem~\ref{thm:tame} and the first row of
Section~\ref{sec:resonance});
\item[(S5)] $f$ is a limit of points of $\Rec$ along which the support
functions $\tilde r$ are asymptotically at least those of $f$
(Theorem~\ref{thm:transitivity});
\item[(S6)] $T$ is the signature-ladder operator of
Definition~\ref{def:SLD}, $I=\{1,\dots,N\}$, and $f\in\mathcal R_0$, i.e.\
$F$ is finite and the signature sets keep room (Theorem~\ref{thm:R0}); no
condition on the blocks of $f$ is needed, and $\mathcal R_0$ contains all
norm-attaining and all base-tame first rows.
\end{enumerate}
For an individual mate $g\in\Cset(f)$, $(f,g)\in\cl\NA$ whenever $g$ lies
in $\cl\Cert^{\rm sh}(f)$; this covers finite and shifted certificates
(Theorems~\ref{thm:transport}, \ref{thm:shifted}), weighted directions with
$O(\sigma)$ box tails (Theorem~\ref{thm:weighted}), locally split mates and
split-contractive operators when $I$ is finite (Theorem~\ref{thm:locsplit},
Corollary~\ref{cor:splitcontractive}), every mate that is a certificate of
radius $\gtrsim t$ up to an $O(t)$ remainder at all small scales
(Corollary~\ref{cor:ladder}), and carrier mates with non-summable gaps
(Theorem~\ref{thm:carriers}).  When $a\in c_{00}$, also every balanced
finite certificate with $\Gamma_w\le1$ is recovered
(Theorem~\ref{thm:transfer}).  When $I$ and $F$ are finite, every mate with
two-piece data (one-sided linear decompositions, possibly through infinitely
many contacts) with $\kappa_w\le1$ is recovered along engineered
norm-attaining approximants if $\Delta d_m\ge0$ in every block
(Corollary~\ref{cor:D1}), and if the blocks with $\Delta d_m<0$ satisfy the
scrambling condition \textup{(SC)} (Theorem~\ref{thm:engineered}); no
hypothesis on the number of active blocks, on the contact set or on rates of
$T$ is needed.  At \textup{(BT)} points the exact one-sided second-order
coefficients are known (Theorem~\ref{thm:onesided}).  For $\ell_2^{d+1}$ we
only prove the joint recovery of tuples of finite certificates under the
hypothesis of Theorem~\ref{thm:joint}.

On the negative side, the intrinsic mechanisms of
Sections~\ref{sec:certificates}--\ref{sec:second} do not suffice in general:
for some admissible $T$ there is a non-attaining $f$ at which the defect
$\Cset(f)\setminus\mathcal K(f)$ spans an infinite-dimensional space
(Theorem~\ref{thm:defect}), and its explicit elements are not recovered
along the canonical truncations (Proposition~\ref{prop:nonrecovery}).  This
is a defect of the intrinsic mechanisms only: that first row satisfies
\textup{(BT)}, and every one of its mates is recovered along engineered
norm-attaining approximants (Corollary~\ref{cor:exampleRecovered}).  At
every point where $a\in c_{00}$ and $K$ is finite (in particular at every
norm-attaining point) exact resonances do not exist
(Proposition~\ref{prop:exact}).  Nothing proved here points to a
counterexample; a counterexample would produce a two-dimensional Hilbert
space without Lindenstrauss' property B, answering Question~6 of Johnson and
Wolfe \cite{JW} negatively.

\subsection{What remains}
By Lemma~\ref{lem:pair} and Corollary~\ref{cor:lsc}(c), the density problem
for $p_N$ is equivalent to the following statement, in which only the defect
matters.

\begin{problem}\label{prob:main}
Let $N\ge1$, let $T$ be admissible, and let $f\in S_{p_N^*}$.  Is it true
that for every $g\in\Def(f)$ and every $\rho<1$ there are norm-attaining
$f_n\to f$ with $\rho g\in\Li_n\Cset(f_n)$?  Equivalently, is every pair
$(f,\rho g)$ with $g\in\Cset(f)\setminus\cl\Cert^{\rm sh}(f)$ a limit of
norm-attaining operators into $\ell_2^2$?
\end{problem}

A positive answer for infinitely many $N$ gives density for Mart\'in's
norm (Lemma~\ref{lem:martintail}); a negative answer for one $N$ and one
admissible $T$ gives a two-dimensional range without property B.  By
Proposition~\ref{prop:nonrecovery}, the approximants in
Problem~\ref{prob:main} cannot in general be the canonical truncations.

\emph{For the designed operator.}  For the operator $T$ of
Definition~\ref{def:SLD} and every $N$, Problem~\ref{prob:main} is
equivalent to Lemma~Z (Problem~\ref{prob:lemmaZ} and
Theorem~\ref{thm:reductionZ}): every mate of every first row must be
approximated by mates at first rows in $\mathcal R_0$, which need not attain
their norms.  What remains is base-side engineering at signature-resonant
first rows (contacts or near-contacts swallowing signature sets, as at the
first row of Section~\ref{sec:resonance}) and at first rows with infinite
base support (Remark~\ref{rem:lemmaZ}).

\emph{For an arbitrary admissible $T$.}  The following classes are open.

\begin{problem}[Negative $d$-mismatch]\label{prob:resonant}
Let $I$ be finite, $f\in S_{p^*}$ with $F$ finite, and let $g\in\Cset(f)$
carry two-piece data with $\kappa_w\le1$ such that the blocks with
$\Delta d_m<0$ do not satisfy \textup{(SC)}.  Is $(f,g)\in\cl\NA$?
\end{problem}

Theorem~\ref{thm:engineered} settles $\Delta d\ge0$ (for any number of
active blocks and without steering) and $\Delta d<0$ under \textup{(SC)};
Remark~\ref{rem:withoutSC}(a) gives a quantitative partial answer, and
Remark~\ref{rem:withoutSC}(b) indicates that \textup{(SC)} can fail for
admissible $T$ (sketch).  The explicit switching mates of
Theorem~\ref{thm:defect} are $d$-neutral and are recovered.

\begin{problem}[Generic supports]\label{prob:generic}
Let $F$ be finite and let some $Q_m$ be infinite, or \textup{(MS)} fail, or
degenerate peaks exist.  Is every $g\in\Cset(f)$ a limit of mates
$g_n\in\Cset(f_n)$, $f_n\to f$, carrying two-piece data to which
Theorem~\ref{thm:engineered} applies?
\end{problem}

At such $f$ the lower bound of Theorem~\ref{thm:onesided} is not available
(the block expansion is multiscale), and mates need not have two-piece data;
mixed mates combining deep weighted parts with switching parts are not
covered either.

\begin{problem}[Second-order defect]\label{prob:secondorder}
Is there $f$ with $a\in c_{00}$ and a balanced finite certificate
$g\in\Cset(f)\cap S(f)$ with $\Gamma_w(g)>1$?  By
Proposition~\ref{prop:lowerbound} such an $f$ must have degenerate peaks,
infinitely many strict non-peaks, failure of (MS), or infinitely many
contacts; in the last case a finite-dimensional analogue suggests that the
second-order coefficient can drop below $\Gamma_w$ (Remark~\ref{rem:kinks}).
At \textup{(BT)} points the question does not affect recoverability:
Corollary~\ref{cor:BTrecovered} recovers every mate, and the relevant
coefficients are the one-sided $\gamma^\pm\le\Gamma_w$
(Theorem~\ref{thm:onesided}).
\end{problem}

\begin{problem}[Approximate resonances]\label{prob:approx}
Do switching mates carried, on both sides of $t=0$, by near-threshold
strict non-peaks with summable gaps or by weak peaks belong to
$\cl\Cert(f)$, or at least to the closure of the mates recovered in
Section~\ref{sec:engineered}?  There is no linear obstruction for them
(Remark~\ref{rem:summablegaps}).  One way to prove non-membership in
$\cl\Cert(f)$ is to exhibit a sequence $f_n\to f$ with
$g\notin\Li_n\Cset(f_n)$ (Corollary~\ref{cor:lsc}).
\end{problem}

Such scale-dependent switching is the main remaining mechanism for an
arbitrary admissible $T$.  We note that operators $T$ have been proposed whose
near-threshold peak carriers are one-sided with bounded conversion capacity;
for them no recovery argument is known (the available arguments leave a
capacity requirement that does not improve as $\rho\to1$), and no
counterexample either.

\begin{problem}[Infinite base support]\label{prob:infiniteF}
Extend Theorems~\ref{thm:onesided}, \ref{thm:engineered} and~\ref{thm:R0}
to first rows with $a\notin c_{00}$.  All results of
Sections~\ref{sec:engineered} and~\ref{sec:designed} use that $F$ is finite
(no flips on $F$ below a fixed scale).
\end{problem}

\begin{problem}[Universality of the defect]\label{prob:universal}
Does \emph{every} admissible $T$ (in particular every $T$ produced by the
proof of Mart\'in's Lemma~B) admit a first row with nonempty defect?  The
linear obstruction requires a first row with few strict non-peaks, which an
adversarial $T$ may prevent; for the designed operator of
Section~\ref{sec:designed} the defect is empty at every first row in
$\mathcal R_0$ (Corollary~\ref{cor:nodefectR0}).
\end{problem}

Finally we note a gap between necessary and sufficient conditions in the
reduction itself.  If $g_1,\dots,g_d\in\Cset(f)$, then
$(g_1,\dots,g_d)/\sqrt d\in\Cset_d(f)$, and density into $\ell_2^{d+1}$
implies that such tuples are recovered along a common norm-attaining
sequence (by the argument of Lemma~\ref{lem:pair}, with rotations of
$\ell_2^{d+1}$).  But we do not
know whether density into $\ell_2^2$ alone forces the existence, for each
$f$, of one norm-attaining sequence along which the whole fibre $\Cset(f)$ is
recovered (condition (a) of Corollary~\ref{cor:finite}); only the
sufficiency of that condition is used.

\begin{thebibliography}{99}
\bibitem{DGS}
G.~Debs, G.~Godefroy and J.~Saint Raymond,
\emph{Topological properties of the set of norm-attaining linear
functionals}, Canad. J. Math. \textbf{47} (1995), 318--329.
\bibitem{JW}
W.~B.~Johnson and J.~Wolfe, \emph{Norm attaining operators},
Studia Math. \textbf{65} (1979), 7--19.
\bibitem{KLMW}
V.~Kadets, G.~L\'opez, M.~Mart\'in and D.~Werner,
\emph{Norm attaining operators of finite rank}, arXiv:1905.08272.
\bibitem{KLMW20}
V.~Kadets, G.~L\'opez, M.~Mart\'in and D.~Werner,
\emph{Equivalent norms with an extremely nonlineable set of norm attaining
functionals}, J. Inst. Math. Jussieu \textbf{19} (2020), 259--279.
\bibitem{Rockafellar}
R.~T.~Rockafellar, \emph{Convex Analysis}, Princeton University Press,
Princeton, 1970.
\bibitem{Martin}
M.~Mart\'in,
\emph{A Banach space whose set of norm-attaining functionals is
algebraically trivial}, J. Funct. Anal. \textbf{288} (2025);
arXiv:2406.07273.
\bibitem{PrA}
\emph{Residual recovery and unbounded component lifts for Mart\'in norms},
accompanying research note (Preprint A), unpublished manuscript (authors
and date not recorded on the manuscript).
\bibitem{PrB}
M.~Han, M.~Mart\'in and D.~L.~Rodr\'iguez-Vidanes,
\emph{Norm attaining operators on renormings of $c_0$}, accompanying
preprint (Preprint B), unpublished manuscript (no year or arXiv number
available); numbering as in the available version.
\bibitem{Check}
\emph{Checking the positive density route for Mart\'in norms},
accompanying research note, unpublished (cited in \cite{PrA}; its local
certificate theorem is re-proved here as Corollary~\ref{cor:check}).
\end{thebibliography}
\end{document}
